\documentclass[10pt]{article}
\usepackage[utf8]{inputenc}
\usepackage[T1]{fontenc}
\usepackage{graphicx}
\usepackage[export]{adjustbox}
\graphicspath{ {./images/} }
\usepackage{hyperref}
\hypersetup{colorlinks=true, linkcolor=blue, filecolor=magenta, urlcolor=cyan,}
\urlstyle{same}
\usepackage{amsmath}
\usepackage{amsfonts}
\usepackage{amssymb}
\usepackage[version=4]{mhchem}
\usepackage{stmaryrd}
\usepackage{bbold}
\usepackage{mathrsfs}
\usepackage{caption}
\usepackage{multirow}

%New command to display footnote whose markers will always be hidden
\let\svthefootnote\thefootnote
\newcommand\blfootnotetext[1]{%
  \let\thefootnote\relax\footnote{#1}%
  \addtocounter{footnote}{-1}%
  \let\thefootnote\svthefootnote%
}
%Overriding the \footnotetext command to hide the marker if its value is `0`
\let\svfootnotetext\footnotetext
\renewcommand\footnotetext[2][?]{%
  \if\relax#1\relax%
    \ifnum\value{footnote}=0\blfootnotetext{#2}\else\svfootnotetext{#2}\fi%
  \else%
    \if?#1\ifnum\value{footnote}=0\blfootnotetext{#2}\else\svfootnotetext{#2}\fi%
    \else\svfootnotetext[#1]{#2}\fi%
  \fi
}
\DeclareUnicodeCharacter{22C5}{\ifmmode\cdot\else{$\cdot$}\fi}
\DeclareUnicodeCharacter{0394}{\ifmmode\Delta\else{$\Delta$}\fi}
\DeclareUnicodeCharacter{22EE}{\ifmmode\vdots\else{$\vdots$}\fi}
\title{Chapter 22: Duration Analysis}

\author{}
\date{}

\begin{document}
\maketitle

\subsection*{22.1 Introduction}
Some response variables in economics come in the form of a duration, which is the time elapsed until a certain event occurs. A few examples include weeks unemployed, months spent on welfare, days until arrest after incarceration, and quarters until an Internet firm files for bankruptcy.

The recent literature on duration analysis is quite rich. In this chapter we focus on the developments that have been used most often in applied work. In addition to providing a rigorous introduction to modern duration analysis, this chapter should prepare you for more advanced treatments, such as Lancaster's (1990) monograph, van den Berg (2001), and Cameron and Trivedi (2005).

Duration analysis has its origins in what is typically called survival analysis, where the duration of interest is survival time of a subject. In survival analysis we are interested in how various treatments or demographic characteristics affect survival times. In the social sciences, we are interested in any situation where an individualor family, or firm, and so on-begins in an initial state and is either observed to exit the state or is censored. (We will discuss the exact nature of censoring in Sections 22.3 and 22.4.) The calendar dates on which units enter the initial state do not have to be the same. (When we introduce covariates in Section 22.2.2, we note how dummy variables for different calendar dates can be included in the covariates, if necessary, to allow for systematic differences in durations by starting date.)

Traditional duration analysis begins by specifying a population distribution for the duration, usually conditional on some explanatory variables (covariates) observed at the beginning of the duration. For example, for the population of people who became unemployed during a particular period, we might observe education levels, experience, marital status-all measured when the person becomes unemployed-wage on prior job, and a measure of unemployment benefits. Then we specify a distribution for the unemployment duration conditional on the covariates. Any reasonable distribution reflects the fact that an unemployment duration is nonnegative. Once a complete conditional distribution has been specified, the same maximum likelihood methods that we studied in Chapter 19 for censored regression models can be used. In this framework, we are typically interested in estimating the effects of the covariates on the expected duration.

Recent treatments of duration analysis tend to focus on the hazard function. The hazard function allows us to approximate the probability of exiting the initial state within a short interval, conditional on having survived up to the starting time of the interval. In econometric applications, hazard functions are usually conditional on some covariates. An important feature for policy analysis is allowing the hazard function to depend on covariates that change over time.

In Section 22.2 we define and discuss hazard functions, and we settle certain issues involved with introducing covariates into hazard functions. In Section 22.3 we show how censored regression models apply to standard duration models with single-cycle flow data, when all covariates are time constant. We also discuss the most common way of introducing unobserved heterogeneity into traditional duration analysis. Given parametric assumptions, we can test for duration dependence-which means that the probability of exiting the initial state depends on the length of time in the state-as well as for the presence of unobserved heterogeneity.

In Section 22.4 we study methods that allow flexible estimation of a hazard function, with both time-constant and time-varying covariates. We assume that we have grouped data; this term means that durations are observed to fall into fixed intervals (often weekly or monthly intervals) and that any time-varying covariates are assumed to be constant within an interval. We focus attention on the case with two states, with everyone in the population starting in the initial state, and single-cycle data, where each person either exits the initial state or is censored before exiting. We also show how heterogeneity can be included when the covariates are strictly exogenous. We touch on some additional issues in Section 22.5.

\subsection*{22.2 Hazard Functions}
The hazard function plays a central role in modern duration analysis. In this section, we discuss various features of the hazard function, both with and without covariates, and provide some examples.

\subsection*{22.2.1 Hazard Functions without Covariates}
Often in this chapter it is convenient to distinguish random variables from particular outcomes of random variables. Let $T \geq 0$ denote the duration, which has some distribution in the population; $t$ denotes a particular value of $T$. (As with any econometric analysis, it is important to be very clear about the relevant population, a topic we consider in Section 22.3.) In survival analysis, $T$ is the length of time a subject lives. Much of the current terminology in duration analysis comes from survival applications. For us, $T$ is the time at which a person (or family, firm, and so on) leaves the initial state. For example, if the initial state is unemployment, $T$ would be the time, measured in, say, weeks, until a person becomes employed.

The cumulative distribution function (cdf) of $T$ is defined as


\begin{equation*}
F(t)=\mathrm{P}(T \leq t), \quad t \geq 0 \tag{22.1}
\end{equation*}


The survivor function is defined as $S(t) \equiv 1-F(t)=\mathrm{P}(T>t)$, and this is the probability of "surviving" past time $t$. We assume in the rest of this section that $T$ is continuous-and, in fact, has a differentiable cdf-because this assumption simplifies statements of certain probabilities. Discreteness in observed durations can be viewed as a consequence of the sampling scheme, as we discuss in Section 22.4. Denote the density of $T$ by $f(t)=\frac{d F}{d t}(t)$.

For $h>0$,


\begin{equation*}
\mathrm{P}(t \leq T<t+h \mid T \geq t) \tag{22.2}
\end{equation*}


is the probabilty of leaving the initial state in the interval $[t, t+h)$ given survival up until time $t$. The hazard function for $T$ is defined as


\begin{equation*}
\lambda(t)=\lim _{h \downarrow 0} \frac{\mathrm{P}(t \leq T<t+h \mid T \geq t)}{h} \tag{22.3}
\end{equation*}


For each $t, \lambda(t)$ is the instantaneous rate of leaving per unit of time. From equation (22.3) it follows that, for "small" $h$,


\begin{equation*}
\mathrm{P}(t \leq T<t+h \mid T \geq t) \approx \lambda(t) h \tag{22.4}
\end{equation*}


Thus the hazard function can be used to approximate a conditional probability in much the same way that the height of the density of $T$ can be used to approximate an unconditional probability.

Example 22.1 (Unemployment Duration): If $T$ is length of time unemployed, measured in weeks, then $\lambda(20)$ is (approximately) the probability of becoming employed between weeks 20 and 21. The phrase "becoming employed" reflects the fact that the person was unemployed up through week 20 . That is, $\lambda(20)$ is roughly the probability of becoming employed between weeks 20 and 21 conditional on having been unemployed through week 20.

Example 22.2 (Recidivism Duration): Suppose $T$ is the number of months before a former prisoner is arrested for a crime. Then $\lambda(12)$ is roughly the probability of being arrested during the 13th month conditional on not having been arrested during the first year.

We can express the hazard function in terms of the density and cdf very simply. First, write\\
$\mathrm{P}(t \leq T<t+h \mid T \geq t)=\mathrm{P}(t \leq T<t+h) / \mathrm{P}(T \geq t)=\frac{F(t+h)-F(t)}{1-F(t)}$

When the cdf is differentiable, we can take the limit of the right-hand side, divided by $h$, as $h$ approaches zero from above:


\begin{equation*}
\lambda(t)=\lim _{h \downarrow 0} \frac{F(t+h)-F(t)}{h} \cdot \frac{1}{1-F(t)}=\frac{f(t)}{1-F(t)}=\frac{f(t)}{S(t)} \tag{22.5}
\end{equation*}


Because the derivative of $S(t)$ is $-f(t)$, we have


\begin{equation*}
\lambda(t)=-\frac{d \log S(t)}{\mathrm{d} t} \tag{22.6}
\end{equation*}


and, using $F(0)=0$, we can integrate to get


\begin{equation*}
F(t)=1-\exp \left[-\int_{0}^{t} \lambda(s) d s\right], \quad t \geq 0 \tag{22.7}
\end{equation*}


Straightforward differentiation of equation (22.7) gives the density of $T$ as


\begin{equation*}
f(t)=\lambda(t) \exp \left[-\int_{0}^{t} \lambda(s) d s\right] \tag{22.8}
\end{equation*}


Therefore, all probabilities can be computed using the hazard function. For example, for points $a_{1}<a_{2}$,

$$
\mathrm{P}\left(T \geq a_{2} \mid T \geq a_{1}\right)=\frac{1-F\left(a_{2}\right)}{1-F\left(a_{1}\right)}=\exp \left[-\int_{a_{1}}^{a_{2}} \lambda(s) d s\right]
$$

and


\begin{equation*}
\mathrm{P}\left(a_{1} \leq T<a_{2} \mid T \geq a_{1}\right)=1-\exp \left[-\int_{a_{1}}^{a_{2}} \lambda(s) d s\right] \tag{22.9}
\end{equation*}


This last expression is especially useful for constructing the log-likelihood functions needed in Section 22.4.

The shape of the hazard function is of primary interest in many empirical applications. In the simplest case, the hazard function is constant:


\begin{equation*}
\lambda(t)=\lambda, \quad \text { all } t \geq 0 \tag{22.10}
\end{equation*}


This function means that the process driving $T$ is memoryless: the probability of exit in the next interval does not depend on how much time has been spent in the initial state. From equation (22.7), a constant hazard implies


\begin{equation*}
F(t)=1-\exp (-\lambda t) \tag{22.11}
\end{equation*}


which is the cdf of the exponential distribution. Conversely, if $T$ has an exponential distribution, it has a constant hazard.

When the hazard function is not constant, we say that the process exhibits duration dependence. Assuming that $\lambda(\cdot)$ is differentiable, there is positive duration dependence at time $t$ if $\mathrm{d} \lambda(t) / \mathrm{d} t>0$; if $\mathrm{d} \lambda(t) / \mathrm{d} t>0$ for all $t>0$, then the process exhibits positive duration dependence. With positive duration dependence, the probability of exiting the initial state increases the longer one is in the initial state. If the derivative is negative, then there is negative duration dependence.

Example 22.3 (Weibull Distribution): If $T$ has a Weibull distribution, its cdf is given by $F(t)=1-\exp \left(-\gamma t^{\alpha}\right)$, where $\gamma$ and $\alpha$ are nonnegative parameters. The density is $f(t)=\gamma \alpha t^{\alpha-1} \exp \left(-\gamma t^{\alpha}\right)$. By equation (22.5), the hazard function is


\begin{equation*}
\lambda(t)=f(t) / S(t)=\gamma \alpha t^{\alpha-1} \tag{22.12}
\end{equation*}


When $\alpha=1$, the Weibull distribution reduces to the exponential with $\lambda=\gamma$. If $\alpha>1$, the hazard is monotonically increasing, so the hazard everywhere exhibits positive duration dependence; for $\alpha<1$, the hazard is monotonically decreasing. Provided we think the hazard is monotonically increasing or decreasing, the Weibull distribution is a relatively simple way to capture duration dependence.

We often want to specify the hazard directly, in which case we can use equation (22.7) to determine the duration distribution.

Example 22.4 (Log-Logistic Hazard Function): The log-logistic hazard function is specified as


\begin{equation*}
\lambda(t)=\frac{\gamma \alpha t^{\alpha-1}}{1+\gamma t^{\alpha}} \tag{22.13}
\end{equation*}


where $\gamma$ and $\alpha$ are positive parameters. When $\alpha=1$, the hazard is monotonically decreasing from $\gamma$ at $t=0$ to zero as $t \rightarrow \infty$; when $\alpha<1$, the hazard is also monotonically decreasing to zero as $t \rightarrow \infty$, but the hazard is unbounded as $t$ approaches zero. When $\alpha>1$, the hazard is increasing until $t=[(\alpha-1) / \gamma]^{1-\alpha}$, and then it decreases to zero.

Straightforward integration gives\\
$\int_{0}^{t} \lambda(s) d s=\log \left(1+\gamma t^{\alpha}\right)=-\log \left[\left(1+\gamma t^{\alpha}\right)^{-1}\right]$\\
so that, by equation (22.7),


\begin{equation*}
F(t)=1-\left(1+\gamma t^{\alpha}\right)^{-1}, \quad t \geq 0 \tag{22.14}
\end{equation*}


Differentiating with respect to $t$ gives\\
$f(t)=\gamma \alpha t^{\alpha-1}\left(1+\gamma t^{\alpha}\right)^{-2}$\\
Using this density, it can be shown that $Y \equiv \log (T)$ has density $g(y)= \alpha \exp [\alpha(y-\mu)] /\{1+\exp [\alpha(y-\mu)]\}^{2}$, where $\mu=-\alpha^{-1} \log (y)$ is the mean of $Y$. In other words, $\log (T)$ has a logistic distribution with mean $\mu$ and variance $\pi^{2} /\left(3 \alpha^{2}\right)$ (hence the name "log-logistic").

\subsection*{22.2.2 Hazard Functions Conditional on Time-Invariant Covariates}
Usually in economics we are interested in hazard functions conditional on a set of covariates or regressors. When these do not change over time-as is often the case given the way many duration data sets are collected - then we simply define the hazard (and all other features of $T$ ) conditional on the covariates. Thus the conditional hazard is

$$
\lambda(t ; \mathbf{x})=\lim _{h \downarrow 0} \frac{\mathrm{P}(t \leq T<t+h \mid T \geq t, \mathbf{x})}{h}
$$

where $\mathbf{x}$ is a vector of explanatory variables. All the formulas from the previous subsection continue to hold provided the cdf and density are defined conditional on $\mathbf{x}$. For example, if the conditional $\operatorname{cdf} F(\cdot \mid \mathbf{x})$ is differentiable, we have


\begin{equation*}
\lambda(t ; \mathbf{x})=\frac{f(t \mid \mathbf{x})}{1-F(t \mid \mathbf{x})} \tag{22.15}
\end{equation*}


where $f(\cdot \mid \mathbf{x})$ is the density of $T$ given $\mathbf{x}$. Often we are interested in the partial effects of the $x_{j}$ on $\lambda(t ; \mathbf{x})$, which are defined as partial derivatives for continuous $x_{j}$ and as differences for discrete $x_{j}$.

If the durations start at different calendar dates-which is usually the case-we can include indicators for different starting dates in the covariates. These allow us to control for seasonal differences in duration distributions.

An especially important class of models with time-invariant regressors consists of proportional hazard models. A proportional hazard can be written as


\begin{equation*}
\lambda(t ; \mathbf{x})=\kappa(\mathbf{x}) \lambda_{0}(t) \tag{22.16}
\end{equation*}


where $\kappa(\cdot)>0$ is a positive function of $\mathbf{x}$ and $\lambda_{0}(t)>0$ is called the baseline hazard. The baseline hazard is common to all units in the population; individual hazard functions differ proportionately based on a function $\kappa(\mathbf{x})$ of observed covariates.

Typically, $\kappa(\cdot)$ is parameterized as $\kappa(\mathbf{x})=\exp (\mathbf{x} \boldsymbol{\beta})$, where $\boldsymbol{\beta}$ is a vector of parameters. Then


\begin{equation*}
\log \lambda(t ; \mathbf{x})=\mathbf{x} \boldsymbol{\beta}+\log \lambda_{0}(t) \tag{22.17}
\end{equation*}


and $\beta_{j}$ measures the semielasticity of the hazard with respect to $x_{j}$. [If $x_{j}$ is the log of an underlying variable, say $x_{j}=\log \left(z_{j}\right), \beta_{j}$ is the elasticity of the hazard with respect to $z_{j}$.]

Occasionally we are interested only in how the covariates shift the hazard function, in which case estimation of $\lambda_{0}$ is not necessary. Cox (1972) obtained a partial maximum likelihood estimator for $\boldsymbol{\beta}$ that does not require estimating $\lambda_{0}(\cdot)$. We discuss Cox's approach briefly in Section 22.5. In economics, much of the time we are interested in the shape of the baseline hazard. We discuss estimation of proportional hazard models with a flexible baseline hazard in Section 22.4.

If in the Weibull hazard function (22.12) we replace $\gamma$ with $\exp (\mathbf{x} \boldsymbol{\beta})$, where the first element of $\mathbf{x}$ is unity, we obtain a proportional hazard model with $\lambda_{0}(t) \equiv \alpha t^{\alpha-1}$. However, if we replace $\gamma$ in equation (22.13) with $\exp (\mathbf{x} \boldsymbol{\beta})$-which is the most common way of introducing covariates into the log-logistic model-we do not obtain a hazard with the proportional hazard form.

Example 22.1 (continued): If $T$ is an unemployment duration, $\mathbf{x}$ might contain education, labor market experience, marital status, race, and number of children, all measured at the beginning of the unemployment spell. Policy variables in $\mathbf{x}$ might reflect the rules governing unemployment benefits, where these are known before each person's unemployment duration.

Example 22.2 (continued): To explain the length of time before arrest after release from prison, the covariates might include participation in a work program while in prison, years of education, marital status, race, time served, and past number of convictions.

\subsection*{22.2.3 Hazard Functions Conditional on Time-Varying Covariates}
Studying hazard functions is more complicated when we wish to model the effects of time-varying covariates on the hazard function. For one thing, it makes no sense to specify the distribution of the duration $T$ conditional on the covariates at only one time period. Nevertheless, we can still define the appropriate conditional probabilities that lead to a conditional hazard function.

Let $\mathbf{x}(t)$ denote the vector of regressors at time $t$; again, this is the random vector describing the population. For $t \geq 0$, let $\mathbf{X}(t), t \geq 0$, denote the covariate path up\\
through time $t: \mathbf{X}(t) \equiv\{\mathbf{x}(s): 0 \leq s \leq t\}$. Following Lancaster (1990, Chapter 2), we define the conditional hazard function at time $t$ by


\begin{equation*}
\lambda[t ; \mathbf{X}(t)]=\lim _{h \downarrow 0} \frac{\mathbf{P}[t \leq T<t+h \mid T \geq t, \mathbf{X}(t+h)]}{h} \tag{22.18}
\end{equation*}


assuming that this limit exists. A discussion of assumptions that ensure existence of equation (22.18) is well beyond the scope of this book; see Lancaster (1990, Chapter 2 ). One case where this limit exists very generally occurs when $T$ is continuous and, for each $t, \mathbf{x}(t+h)$ is constant for all $h \in[0, \eta(t)]$ for some function $\eta(t)>0$. Then we can replace $\mathbf{X}(t+h)$ with $\mathbf{X}(t)$ in equation (22.18) [because $\mathbf{X}(t+h)=\mathbf{X}(t)$ for $h$ sufficiently small]. For reasons we will see in Section 22.4, we must assume that timevarying covariates are constant over the interval of observation (such as a week or a month), anyway, in which case there is no problem in defining equation (22.18).

For certain purposes, it is important to know whether time-varying covariates are strictly exogenous. With the hazard defined as in equation (22.18), Lancaster (1990, Definition 2.1) provides a definition that rules out feedback from the duration to future values of the covariates. Specifically, if $\mathbf{X}(t, t+h)$ denotes the covariate path from time $t$ to $t+h$, then Lancaster's strict exogeneity condition is


\begin{equation*}
\mathrm{P}[\mathbf{X}(t, t+h) \mid T \geq t+h, \mathbf{X}(t)]=\mathrm{P}[\mathbf{X}(t, t+h) \mid \mathbf{X}(t)] \tag{22.19}
\end{equation*}


for all $t \geq 0, h>0$. Actually, when condition (22.19) holds, Lancaster says $\{\mathbf{x}(t)$ : $t>0\}$ is "exogenous." We prefer the name "strictly exogenous" because condition (22.19) is closely related to the notions of strict exogeneity that we have encountered throughout this book. Plus, it is important to see that condition (22.19) has nothing to do with contemporaneous endogeneity: by definition, the covariates are sequentially exogenous (see Section 7.4) because, by specifying $\lambda[t ; \mathbf{X}(t)]$, we are conditioning on current and past covariates.

Equation (22.19) applies to covariates whose entire path is well defined whether or not the agent is in the initial state. One such class of covariates, called external covariates by Kalbfleisch and Prentice (1980), has the feature that the covariate path is independent of whether any particular agent has or has not left the initial state. In modeling time until arrest, these covariates might include law enforcement per capita in the person's city of residence or the city unemployment rate.

Other covariates are not external to each agent but have paths that are still defined after the agent leaves the initial state. For example, marital status is well defined before and after someone is arrested, but it is possibly related to whether someone has been arrested. Whether marital status satisfies condition (22.19) is an empirical issue.

The definition of strict exogeneity in condition (22.19) cannot be applied to timevarying covariates whose path is not defined once the agent leaves the initial state. Kalbfleisch and Prentice (1980) call these internal covariates. Lancaster (1990, p. 28) gives the example of job tenure duration, where a time-varying covariate is wage paid on the job: if a person leaves the job, it makes no sense to define the future wage path in that job. As a second example, in modeling the time until a former prisoner is arrested, a time-varying covariate at time $t$ might be wage income in the previous month, $t-1$. If someone is arrested and reincarcerated, it makes little sense to define future labor income.

It is pretty clear that internal covariates cannot satisfy any reasonable strict exogeneity assumption. This fact will be important in Section 22.4 when we discuss estimation of duration models with unobserved heterogeneity and grouped duration data. We will actually use a slightly different notion of strict exogeneity that is directly relevant for conditional maximum likelihood estimation. Nevertheless, it is in the same spirit as condition (22.19).

With time-varying covariates there is not, strictly speaking, such a thing as a proportional hazard model. Nevertheless, it has become common in econometrics to call a hazard of the form


\begin{equation*}
\lambda[t ; \mathbf{x}(t)]=\kappa[\mathbf{x}(t)] \lambda_{0}(t) \tag{22.20}
\end{equation*}


a proportional hazard with time-varying covariates. The function multiplying the baseline hazard is usually $\kappa[\mathbf{x}(t)]=\exp [\mathbf{x}(t) \boldsymbol{\beta}]$; for notational reasons, we show this depending only on $\mathbf{x}(t)$ and not on past covariates [which can always be included in $\mathbf{x}(t)]$. We will discuss estimation of these models, without the strict exogeneity assumption, in Section 22.4.2. In Section 22.4.3, when we multiply equation (22.20) by unobserved heterogeneity, strict exogeneity becomes very important.

The log-logistic hazard is also easily modified to have time-varying covariates. One way to include time-varying covariates parametrically is\\
$\lambda[t ; \mathbf{x}(t)]=\exp [\mathbf{x}(t) \boldsymbol{\beta}] \alpha t^{\alpha-1} /\left\{1+\exp [\mathbf{x}(t) \boldsymbol{\beta}] t^{\alpha}\right\}$\\
We will see how to estimate $\alpha$ and $\boldsymbol{\beta}$ in Section 22.4.2.

\subsection*{22.3 Analysis of Single-Spell Data with Time-Invariant Covariates}
We assume that the population of interest is individuals entering the initial state during a given interval of time, say $[0, b]$, where $b>0$ is a known constant. (Naturally, "individual" can be replaced with any population unit of interest, such as "family"\\
or "firm.") As in all econometric contexts, it is very important to be explicit about the underlying population. By convention, we let zero denote the earliest calendar date that an individual can enter the initial state, and $b$ is the last possible date. For example, if we are interested in the population of U.S. workers who became unemployed at any time during 1998, and unemployment duration is measured in years (with .5 meaning half a year), then $b=1$. If duration is measured in weeks, then $b=52$; if duration is measured in days, then $b=365$; and so on.

In using the methods of this section, we typically ignore the fact that durations are often grouped into discrete intervals-for example, measured to the nearest week or month-and treat them as continuously distributed. If we want to explicitly recognize the discreteness of the measured durations, we should treat them as grouped data, as we do in Section 22.4.

We restrict attention to single-spell data. That is, we use, at most, one completed spell per individual. If, after leaving the initial state, an individual subsequently reenters the initial state in the interval $[0, b]$, we ignore this information. In addition, the covariates in the analysis are time invariant, which means we collect covariates on individuals at a given point in time-usually, at the beginning of the spell-and we do not re-collect data on the covariates during the course of the spell. Time-varying covariates are more naturally handled in the context of grouped duration data in Section 22.4.

We study two general types of sampling from the population that we have described. The most common, and the easiest to handle, is flow sampling. In Section 22.3.3 we briefly consider various kinds of stock sampling.

\subsection*{22.3.1 Flow Sampling}
With flow sampling, we sample individuals who enter the state at some point during the interval $[0, b]$, and we record the length of time each individual is in the initial state. We collect data on covariates known at the time the individual entered the initial state. For example, suppose we are interested in the population of U.S. workers who became unemployed at any time during 1998, and we randomly sample from U.S. male workers who became unemployed during 1998. At the beginning of the unemployment spell we might obtain information on tenure in last job, wage on last job, gender, marital status, and information on unemployment benefits.

There are two common ways to collect flow data on unemployment spells. First, we may randomly sample individuals from a large population, say, all working-age individuals in the United States for a given year, say, 1998. Some fraction of these people will be in the labor force and will become unemployed during 1998-that is,\\
enter the initial state of unemployment during the specified interval-and this group of people who become unemployed is our random sample of all workers who become unemployed during 1998. Another possibility is retrospective sampling. For example, suppose that, for a given state in the United States, we have access to unemployment records for 1998. We can then obtain a random sample of all workers who became unemployed during 1998.

Flow data are usually subject to right censoring. That is, after a certain amount of time, we stop following the individuals in the sample, which we must do in order to analyze the data. (Right censoring is the only kind that occurs with flow data, so we will often refer to right censoring as "censoring" in this and the next subsection.) For individuals who have completed their spells in the initial state, we observe the exact duration. But for those still in the initial state, we only know that the duration lasted as long as the tracking period. In the unemployment duration example, we might follow each individual for a fixed length of time, say, two years. If unemployment spells are measured in weeks, we would have right censoring at 104 weeks. Alternatively, we might stop tracking individuals at a fixed calendar date, say, the last week in 1999. Because individuals can become unemployed at any time during 1998, calendar-date censoring results in censoring times that differ across individuals.

\subsection*{22.3.2 Maximum Likelihood Estimation with Censored Flow Data}
For a random draw $i$ from the population, let $a_{i} \in[0, b]$ denote the time at which individual $i$ enters the initial state (the "starting time"), let $t_{i}^{*}$ denote the length of time in the initial state (the duration), and let $\mathbf{x}_{i}$ denote the vector of observed covariates. We assume that $t_{i}^{*}$ has a continuous conditional density $f\left(t \mid \mathbf{x}_{i} ; \boldsymbol{\theta}\right), t \geq 0$, where $\boldsymbol{\theta}$ is the vector of unknown parameters.

Without right censoring we would observe a random sample on $\left(a_{i}, t_{i}^{*}, \mathbf{x}_{i}\right)$, and estimation would be a standard exercise in conditional maximum likelihood. To account for right censoring, we assume that the observed duration, $t_{i}$, is obtained as


\begin{equation*}
t_{i}=\min \left(t_{i}^{*}, c_{i}\right) \tag{22.21}
\end{equation*}


where $c_{i}$ is the censoring time for individual $i$. In some cases, $c_{i}$ is constant across $i$. For example, suppose $t_{i}^{*}$ is unemployment duration for person $i$, measured in weeks. If the sample design specifies that we follow each person for at most two years, at which point all people remaining unemployed after two years are censored, then $c=$ 104. If we have a fixed calendar date at which we stop tracking individuals, the censoring time differs by individual because the workers typically would become unemployed on different calendar dates. If $b=52$ weeks and we censor everyone at two years from the start of the study, the censoring times could range from 52 to 104 weeks.

We assume that, conditional on the covariates, the true duration is independent of the starting point, $a_{i}$, and the censoring time, $c_{i}$ :


\begin{equation*}
\mathrm{D}\left(t_{i}^{*} \mid \mathbf{x}_{i}, a_{i}, c_{i}\right)=\mathrm{D}\left(t_{i}^{*} \mid \mathbf{x}_{i}\right) \tag{22.22}
\end{equation*}


where $\mathrm{D}(\cdot \mid \cdot)$ denotes conditional distribution. Assumption (22.22) clearly holds when $a_{i}$ and $c_{i}$ are constant for all $i$, but it holds under much weaker assumptions. Sometimes $c_{i}$ is constant for all $i$, in which case assumption (22.22) holds when the duration is independent of the starting time, conditional on $\mathbf{x}_{i}$. If there are seasonal effects on the duration-for example, unemployment durations that start in the summer have a different expected length than durations that start at other times of the year-then we may have to put dummy variables for different starting dates in $\mathbf{x}_{i}$ to ensure that assumption (22.22) holds. This approach would also ensure that assumption (22.22) holds when a fixed calendar date is used for censoring, implying that $c_{i}$ is not constant across $i$. Assumption (22.22) holds for certain nonstandard censoring schemes, too. For example, if an element of $\mathbf{x}_{i}$ is education, assumption (22.22) holds if, say, individuals with more education are censored more quickly.

Under assumption (22.22), the distribution of $t_{i}^{*}$ given ( $\mathbf{x}_{i}, a_{i}, c_{i}$ ) does not depend on $\left(a_{i}, c_{i}\right)$. Therefore, if the duration is not censored, the density of $t_{i}=t_{i}^{*}$ given $\left(\mathbf{x}_{i}, a_{i}, c_{i}\right)$ is simply $f\left(t \mid \mathbf{x}_{i} ; \boldsymbol{\theta}\right)$. The probability that $t_{i}$ is censored is\\
$\mathrm{P}\left(t_{i}^{*} \geq c_{i} \mid \mathbf{x}_{i}\right)=1-F\left(c_{i} \mid \mathbf{x}_{i} ; \boldsymbol{\theta}\right)$\\
where $F\left(t \mid \mathbf{x}_{i} ; \boldsymbol{\theta}\right)$ is the conditional cdf of $t_{i}^{*}$ given $\mathbf{x}_{i}$. Letting $d_{i}$ be a censoring indicator ( $d_{i}=1$ if uncensored, $d_{i}=0$ if censored), the conditional likelihood for observation $i$ can be written as


\begin{equation*}
f\left(t_{i} \mid \mathbf{x}_{i} ; \boldsymbol{\theta}\right)^{d_{i}}\left[1-F\left(t_{i} \mid \mathbf{x}_{i} ; \boldsymbol{\theta}\right)\right]^{\left(1-d_{i}\right)} \tag{22.23}
\end{equation*}


Importantly, neither the starting times, $a_{i}$, nor the length of the interval, $b$, plays a role in the analysis. (In fact, in the vast majority of treatments of flow data, $b$ and $a_{i}$ are not even introduced. However, it is important to know that the reason $a_{i}$ is not relevant for the analysis of flow data is the conditional independence assumption in equation (22.22).) By contrast, the censoring times $c_{i}$ do appear in the likelihood for censored observations because then $t_{i}=c_{i}$. Given data on ( $t_{i}, d_{i}, \mathbf{x}_{i}$ ) for a random sample of size $N$, the maximum likelihood estimator of $\boldsymbol{\theta}$ is obtained by maximizing


\begin{equation*}
\sum_{i=1}^{N}\left\{d_{i} \log \left[f\left(t_{i} \mid \mathbf{x}_{i} ; \boldsymbol{\theta}\right)\right]+\left(1-d_{i}\right) \log \left[1-F\left(t_{i} \mid \mathbf{x}_{i} ; \boldsymbol{\theta}\right)\right]\right\} \tag{22.24}
\end{equation*}


For the choices of $f(\cdot \mid \mathbf{x} ; \boldsymbol{\theta})$ used in practice, the conditional MLE regularity conditions-see Chapter 13-hold, and the MLE is $\sqrt{N}$-consistent and asymptotically normal. (If there is no censoring, $d_{i}=1$ for all $i$ and the second term in expression (22.24) is simply dropped.)

Because the hazard function can be expressed as in equation (22.15), once we specify $f$, the hazard function can be estimated once we have the MLE, $\hat{\boldsymbol{\theta}}$. For example, the Weibull distribution with covariates has conditional density


\begin{equation*}
f\left(t \mid \mathbf{x}_{i} ; \boldsymbol{\theta}\right)=\exp \left(\mathbf{x}_{i} \boldsymbol{\beta}\right) \alpha t^{\alpha-1} \exp \left[-\exp \left(\mathbf{x}_{i} \boldsymbol{\beta}\right) t^{\alpha}\right] \tag{22.25}
\end{equation*}


where $\mathbf{x}_{i}$ contains unity as its first element for all $i$. (We obtain this density from Example 22.3 with $\gamma$ replaced by $\exp \left(\mathbf{x}_{i} \boldsymbol{\beta}\right)$.) The hazard function in this case is simply $\lambda(t ; \mathbf{x})=\exp (\mathbf{x} \boldsymbol{\beta}) \alpha t^{\alpha-1}$.

Example 22.5 (Weibull Model for Recidivism Duration): Let durat be the length of time, in months, until an inmate is arrested after being released from prison. Although the duration is rounded to the nearest month, we treat durat as a continuous variable with a Weibull distribution. We are interested in how certain covariates affect the hazard function for recidivism, and also whether there is positive or negative duration dependence, once we have conditioned on the covariates. The variable workprg - a binary indicator for participation in a prison work program-is of particular interest.

The data in RECID.RAW, which comes from Chung, Schmidt, and Witte (1991), are flow data because it is a random sample of convicts released from prison during the period July 1, 1977, through June 30, 1978. The data are retrospective in that they were obtained by looking at records in April 1984, which served as the common censoring date. Because of the different starting times, the censoring times, $c_{i}$, vary from 70 to 81 months. The results of the Weibull estimation are in Table 22.1.

In interpreting the estimates, we use equation (22.17). For small $\hat{\beta}_{j}$, we can multiply the coefficient by 100 to obtain the semielasticity of the hazard with respect to $x_{j}$. (No covariates appear in logarithmic form, so there are no elasticities among the $\hat{\beta}_{j}$.) For example, if tserved increases by one month, the hazard shifts up by about 1.4 percent, and the effect is statistically significant. Another year of education reduces the hazard by about 2.3 percent, but the effect is insignificant at even the 10 percent level against a two-sided alternative.

The sign of the workprg coefficient is unexpected, at least if we expect the work program to have positive benefits after the inmates are released from prison. (The result is not statistically different from zero.) The reason could be that the program is ineffective or that there is self-selection into the program.

\begin{table}[h]
\begin{center}
\captionsetup{labelformat=empty}
\caption{Table 22.1\\
Weibull Estimation of Criminal Recidivism}
\begin{tabular}{|l|l|}
\hline
Explanatory Variable & Coefficient (Standard Error) \\
\hline
workprg & \begin{tabular}{l}
. 091 \\
(.091) \\
\end{tabular} \\
\hline
priors & \begin{tabular}{l}
. 089 \\
(.013) \\
\end{tabular} \\
\hline
tserved & \begin{tabular}{l}
. 014 \\
(.002) \\
\end{tabular} \\
\hline
felon & \begin{tabular}{l}
-. 299 \\
(.106) \\
\end{tabular} \\
\hline
alcohol & \begin{tabular}{l}
. 447 \\
(.106) \\
\end{tabular} \\
\hline
drugs & \begin{tabular}{l}
. 281 \\
(.098) \\
\end{tabular} \\
\hline
black & \begin{tabular}{l}
. 454 \\
(.088) \\
\end{tabular} \\
\hline
married & \begin{tabular}{l}
-. 152 \\
(.109) \\
\end{tabular} \\
\hline
educ & \begin{tabular}{l}
-. 023 \\
(.019) \\
\end{tabular} \\
\hline
age & \begin{tabular}{l}
-. 0037 \\
(.0005) \\
\end{tabular} \\
\hline
constant & \begin{tabular}{l}
-3.402 \\
(0.301) \\
\end{tabular} \\
\hline
$\hat{\alpha}$ & \begin{tabular}{l}
. 806 \\
(.031) \\
\end{tabular} \\
\hline
Observations & 1,445 \\
\hline
Log likelihood & -1,633.03 \\
\hline
\end{tabular}
\end{center}
\end{table}

For large $\hat{\beta}_{j}$, we should exponentiate and subtract unity to obtain the proportionate change. For example, at any point in time, the hazard is about $100[\exp (.447)-1] =56.4$ percent greater for someone with an alcohol problem than for someone without.

The estimate of $\alpha$ is .806, and the standard error of $\hat{\alpha}$ leads to a strong rejection of $\mathrm{H}_{0}: \alpha=1$ against $\mathrm{H}_{0}: \alpha<1$. Therefore, there is evidence of negative duration dependence, conditional on the covariates. This means that, for a particular ex-convict, the instantaneous rate of being arrested decreases with the length of time out of prison. Figure 22.1 provides a graph of the estimated Weibull hazard.

When the Weibull model is estimated without the covariates, $\hat{\alpha}=.770$ ( $\mathrm{se}=.031$ ), which shows slightly more negative duration dependence. This is a typical finding in

\begin{figure}[h]
\begin{center}
  \includegraphics[alt={},max width=\textwidth]{2e17339c-7635-4ef8-aded-2191d8d853ff-1028_913_1256_330_288}
\captionsetup{labelformat=empty}
\caption{Figure 22.1\\
Weibull hazard of recidivism}
\end{center}
\end{figure}

applications of Weibull duration models: estimated $\alpha$ without covariate tends to be less than the estimate with covariates. Lancaster (1990, Section 10.2) contains a theoretical discussion based on unobserved heterogeneity.

When we are primarily interested in the effects of covariates on the expected duration (rather than on the hazard), we can apply a censored Tobit analysis to the $\log$ of the duration. A Tobit analysis assumes that, for each random draw $i, \log \left(t_{i}^{*}\right)$ given $\mathbf{x}_{i}$ has a $\operatorname{Normal}\left(\mathbf{x}_{i} \boldsymbol{\delta}, \sigma^{2}\right)$ distribution, which implies that $t_{i}^{*}$ given $\mathbf{x}_{i}$ has a lognormal distribution. (The first element of $\mathbf{x}_{i}$ is unity.) The hazard function for a lognormal distribution, conditional on $\mathbf{x}$, is $\lambda(t ; \mathbf{x})=h[(\log t-\mathbf{x} \boldsymbol{\delta}) / \sigma] / \sigma t$, where $h(z) \equiv \phi(z) /[1-\Phi(z)], \phi(\cdot)$ is the standard normal probability density function (pdf), and $\Phi(\cdot)$ is the standard normal cdf. The lognormal hazard function is not monotonic and does not have the proportional hazard form. Nevertheless, the estimates of the $\delta_{j}$ are easy to interpret because the model is equivalent to


\begin{equation*}
\log \left(t_{i}^{*}\right)=\mathbf{x}_{i} \boldsymbol{\delta}+e_{i} \tag{22.26}
\end{equation*}


where $e_{i}$ is independent of $\mathbf{x}_{i}$ and normally distributed. Therefore, the $\delta_{j}$ are semielasticities-or elasticities if the covariates are in logarithmic form-of the covariates on the expected duration.

The Weibull model can also be represented in regression form. When $t_{i}^{*}$ given $\mathbf{x}_{i}$ has density (22.25), $\exp \left(\mathbf{x}_{i} \boldsymbol{\beta}\right)\left(t_{i}^{*}\right)^{\alpha}$ is independent of $\mathbf{x}_{i}$ and has a unit exponential distribution. Therefore, its natural $\log$ has a type I extreme value distribution; therefore, we can write $\alpha \log \left(t_{i}^{*}\right)=-\mathbf{x}_{i} \boldsymbol{\beta}+u_{i}$, where $u_{i}$ is independent of $\mathbf{x}_{i}$ and has density $g(u)=\exp (u) \exp \{\exp (-u)\}$. The mean of $u_{i}$ is not zero, but, because $u_{i}$ is independent of $\mathbf{x}_{i}$, we can write $\log \left(t_{i}^{*}\right)$ exactly as in equation (22.26), where the slope coefficents are given by $\delta_{j}=-\beta_{j} / \alpha$, and the intercept is more complicated. Now, $e_{i}$ does not have a normal distribution, but it is independent of $\mathbf{x}_{i}$ with zero mean. Censoring can be handled by maximum likelihood estimation. The estimated coefficients can be compared with the censored Tobit estimates described previously to see if the estimates are sensitive to the distributional assumption.

In Example 22.5, we can obtain the Weibull estimates of the $\delta_{j}$ as $\hat{\delta}_{j}=-\hat{\beta}_{j} / \hat{\alpha}$. (Some econometrics packages, such as Stata, allow direct estimation of the $\delta_{j}$ and provide standard errors.) For example, $\hat{\delta}_{\text {drugs }}=-.281 / .806 \approx-.349$. When the lognormal model is used, the coefficient on drugs is somewhat smaller in magnitude, about -.298. As another example, $\hat{\delta}_{\text {age }}=.0046$ in the Weibull estimation and $\hat{\delta}_{\text {age }}=.0039$ in the lognormal estimation. In both cases, the estimates have $t$ statistics greater than six. For obtaining estimates on the expected duration, the Weibull and lognormal models give similar results. The lognormal model fits the data notably better, with log likelihood $=-1,597.06$. This result is consistent with the findings of Chung, Schmidt, and Witte (1991).

Importantly, the shapes of the lognormal and Weibull hazards are quite different. The lognormal hazard, which is plotted in Figure 22.2 with the covariates set at their mean values, first increases until about six months and then decreases thereafter. This shape implies that for a short period-roughly six months-the instantaneous probability of being arrested for a new crime increases the longer an ex-convict has been released from prison. After that, the hazard falls to about .001 at 80 months. Therefore, conditional on an ex-convict being out of prison $6 \frac{1}{2}$ years, the probability of being arrested for a new crime in the 81st month is roughly .001 . The Weibull hazard implies that there is initially a greater than .01 instantaneous probability of being arrested. At 80 months, the probability of being arrested during the 81st month is about .005 , or five times larger than that obtained from the lognormal model.

Sometimes we begin by specifying a parametric model for the hazard conditional on $\mathbf{x}$ and then use the formulas from Section 22.2 to obtain the cdf and density. This

\begin{figure}[h]
\begin{center}
  \includegraphics[alt={},max width=\textwidth]{2e17339c-7635-4ef8-aded-2191d8d853ff-1030_915_1258_325_286}
\captionsetup{labelformat=empty}
\caption{Figure 22.2\\
Lognormal hazard of recidivism}
\end{center}
\end{figure}

approach is easiest when the hazard leads to a tractable duration distribution, but there is no reason the hazard function must be of the proportional hazard form.

Example 22.6 (Log-Logistic Hazard with Covariates): A log-logistic hazard function with covariates is


\begin{equation*}
\lambda(t ; \mathbf{x})=\exp (\mathbf{x} \boldsymbol{\beta}) \alpha t^{\alpha-1} /\left[1+\exp (\mathbf{x} \boldsymbol{\beta}) t^{\alpha}\right] \tag{22.27}
\end{equation*}


where $x_{1} \equiv 1$. From equation (22.14) with $\gamma=\exp (\mathbf{x} \boldsymbol{\beta})$, the $\operatorname{cdf}$ is


\begin{equation*}
F(t \mid \mathbf{x} ; \boldsymbol{\theta})=1-\left[1+\exp (\mathbf{x} \boldsymbol{\beta}) t^{\alpha}\right]^{-1}, \quad t \geq 0 \tag{22.28}
\end{equation*}


The distribution of $\log \left(t_{i}^{*}\right)$ given $\mathbf{x}_{i}$ is logistic with mean $-\alpha^{-1} \log \{\exp (\mathbf{x} \boldsymbol{\beta})\}= -\alpha^{-1} \mathbf{x} \boldsymbol{\beta}$ and variance $\pi^{2} /\left(3 \alpha^{2}\right)$. Therefore, $\log \left(t_{i}^{*}\right)$ can be written as in equation (22.26) where $e_{i}$ has a zero mean logistic distribution and is independent of $\mathbf{x}_{i}$ and $\boldsymbol{\delta}=-\alpha^{-1} \boldsymbol{\beta}$. This is another example where the effects of the covariates on the mean\\
duration can be obtained by an OLS regression when there is no censoring. With censoring, the distribution of $e_{i}$ must be accounted for using the log likelihood in expression (22.24).

\subsection*{22.3.3 Stock Sampling}
Flow data with right censoring are common, but other sampling schemes are also used. With stock sampling we randomly sample from individuals that are in the initial state at a given point in time. The population is again individuals who enter the initial state during a specified interval, $[0, b]$. However, rather than observe a random sample of people flowing into the initial state, we can only obtain a random sample of individuals that are in the initial state at time $b$. In addition to the possibility of right censoring, we may also face the problem of left censoring, which occurs when some or all of the starting times, $a_{i}$, are not observed. For now, we assume that (1) we observe the starting times $a_{i}$ for all individuals we sample at time $b$ and (2) we can follow sampled individuals for a certain length of time after we observe them at time $b$. We also allow for right censoring.

In the unemployment duration example, where the population comprises workers who became unemployed at some point during 1998, stock sampling would occur if we randomly sampled from workers who were unemployed during the last week of 1998. This kind of sampling causes a clear sample selection problem: we necessarily exclude from our sample any individual whose unemployment spell ended before the last week of 1998. Because these spells were necessarily shorter than a year, we cannot just assume that the missing observations are randomly missing.

The sample selection problem caused by stock sampling is essentially the same situation we faced in Section 19.5, where we covered the truncated regression model. Therefore, we will call this the left truncation problem. Kiefer (1988) calls it lengthbiased sampling.

Under the assumptions that we observe the $a_{i}$ and can observe some spells past the sampling date $b$, left truncation is fairly easy to deal with. With the exception of replacing flow sampling with stock sampling, we make the same assumptions as in Section 22.3.2.

To account for the truncated sampling, we must modify the density in equation (22.23) to reflect the fact that part of the population is systematically omitted from the sample. Let $\left(a_{i}, c_{i}, \mathbf{x}_{i}, t_{i}\right)$ denote a random draw from the population of all spells starting in $[0, b]$. We observe this vector if and only if the person is still in the initial state at time $b$, that is, if and only if $a_{i}+t_{i}^{*} \geq b$ or $t_{i}^{*} \geq b-a_{i}$, where $t_{i}^{*}$ is the true duration. But, under the conditional independence assumption (22.22),


\begin{equation*}
\mathrm{P}\left(t_{i}^{*} \geq b-a_{i} \mid a_{i}, c_{i}, \mathbf{x}_{i}\right)=1-F\left(b-a_{i} \mid \mathbf{x}_{i} ; \boldsymbol{\theta}\right) \tag{22.29}
\end{equation*}


where $F\left(\cdot \mid \mathbf{x}_{i} ; \boldsymbol{\theta}\right)$ is the cdf of $t_{i}^{*}$ given $\mathbf{x}_{i}$, as before. The correct conditional density function is obtained by dividing equation (22.23) by equation (22.29). In Problem 22.5 you are asked to adapt the arguments in Section 19.5 to also allow for right censoring. The log-likelihood function can be written as


\begin{equation*}
\sum_{i=1}^{N}\left\{d_{i} \log \left[f\left(t_{i} \mid \mathbf{x}_{i} ; \boldsymbol{\theta}\right)\right]+\left(1-d_{i}\right) \log \left[1-F\left(t_{i} \mid \mathbf{x}_{i} ; \boldsymbol{\theta}\right)\right]-\log \left[1-F\left(b-a_{i} \mid \mathbf{x}_{i} ; \boldsymbol{\theta}\right)\right]\right\} \tag{22.30}
\end{equation*}


where, again, $t_{i}=c_{i}$ when $d_{i}=0$. Unlike in the case of flow sampling, with stock sampling both the starting dates, $a_{i}$, and the length of the sampling interval, $b$, appear in the conditional likelihood function. Their presence makes it clear that specifying the interval $[0, b]$ is important for analyzing stock data. [Lancaster (1990, p. 183) essentially derives equation (22.30) under a slightly different sampling scheme; see also Lancaster (1979).]

Equation (22.30) has an interesting implication. If observation $i$ is right censored at calendar date $b$-that is, if we do not follow the spell after the initial data collectionthen the censoring time is $c_{i}=b-a_{i}$. Because $d_{i}=0$ for censored observations, the log likelihood for such an observation is $\log \left[1-F\left(c_{i} \mid \mathbf{x}_{i} ; \boldsymbol{\theta}\right)\right]-\log \left[1-F\left(b-a_{i} \mid \mathbf{x}_{i} ; \boldsymbol{\theta}\right)\right]=$ 0 . In other words, observations that are right censored at the data collection time provide no information for estimating $\boldsymbol{\theta}$, at least when we use equation (22.30). Consequently, the log likelihood in equation (22.30) does not identify $\boldsymbol{\theta}$ if all units are right censored at the interview date: equation ( 22.30 ) is identically zero. The intuition for why equation (22.30) fails in this case is fairly clear: our data consist only of $\left(a_{i}, \mathbf{x}_{i}\right)$, and equation (22.30) is a log likelihood that is conditional on $\left(a_{i}, \mathbf{x}_{i}\right)$. Effectively, there is no random response variable.

Even when we censor all observed durations at the interview date, we can still estimate $\boldsymbol{\theta}$, provided - at least in a parametric context-we specify a model for the conditional distribution of the starting times, $\mathrm{D}\left(a_{i} \mid \mathbf{x}_{i}\right)$. (This is essentially the problem analyzed by Nickell, 1979.) We are still assuming that we observe the $a_{i}$. So, for example, we randomly sample from the pool of people unemployed in the last week of 1998 and find out when their unemployment spells began (along with covariates). We do not follow any spells past the interview date. (As an aside, if we sample unemployed people during the last week of 1998, we are likely to obtain some observations where spells began before 1998. For the population we have specified, these people would simply be discarded. If we want to include people whose spells began\\
prior to 1998, we need to redefine the interval. For example, if durations are measured in weeks and if we want to consider durations beginning in the five-year period prior to the end of 1998, then $b=260$.)

For concreteness, we assume that $\mathrm{D}\left(a_{i} \mid \mathbf{x}_{i}\right)$ is continuous on $[0, b]$ with density $k\left(\cdot \mid \mathbf{x}_{i} ; \boldsymbol{\eta}\right)$. Let $s_{i}$ denote a sample selection indicator, which is unity if we observe random draw $i$, that is, if $t_{i}^{*} \geq b-a_{i}$. Estimation of $\boldsymbol{\theta}$ (and $\boldsymbol{\eta}$ ) can proceed by applying CMLE to the density of $a_{i}$ conditional on $\mathbf{x}_{i}$ and $s_{i}=1$. (Note that this is the only density we can hope to estimate, as our sample only consists of observations $\left(a_{i}, \mathbf{x}_{i}\right)$ when $s_{i}=1$.) This density is informative for $\boldsymbol{\theta}$ even if $\boldsymbol{\eta}$ is not functionally related to $\boldsymbol{\theta}$ (as would typically be assumed) because there are some durations that started and ended in $[0, b]$; we simply do not observe them. Knowing something about the starting time distribution gives us information about the duration distribution. (In the context of flow sampling, when $\boldsymbol{\eta}$ is not functionally related to $\boldsymbol{\theta}$, the density of $a_{i}$ given $\mathbf{x}_{i}$ is uninformative for estimating $\boldsymbol{\theta}$; in other words, $a_{i}$ is ancillary for $\boldsymbol{\theta}$.)

In Problem 22.6 you are asked to show that the density of $a_{i}$ conditional on observing $\left(a_{i}, \mathbf{x}_{i}\right)$ is


\begin{equation*}
p\left(a \mid \mathbf{x}_{i}, s_{i}=1\right)=k\left(a \mid \mathbf{x}_{i} ; \boldsymbol{\eta}\right)\left[1-F\left(b-a \mid \mathbf{x}_{i} ; \boldsymbol{\theta}\right)\right] / \mathrm{P}\left(s_{i}=1 \mid \mathbf{x}_{i} ; \boldsymbol{\theta}, \boldsymbol{\eta}\right) \tag{22.31}
\end{equation*}


$0<a<b$, where


\begin{equation*}
\mathrm{P}\left(s_{i}=1 \mid \mathbf{x}_{i} ; \boldsymbol{\theta}, \boldsymbol{\eta}\right)=\int_{0}^{b}\left[1-F\left(b-u \mid \mathbf{x}_{i} ; \boldsymbol{\theta}\right)\right] k\left(u \mid \mathbf{x}_{i} ; \boldsymbol{\eta}\right) \mathrm{d} u \tag{22.32}
\end{equation*}


(Lancaster (1990, Section 8.3.3) essentially obtains the right-hand side of equation (22.31) but uses the notion of backward recurrence time. The argument in Problem 22.6 is more straightforward because it is based on a standard truncation argument.) Once we have specified the duration cdf, $F$, and the starting time density, $k$, we can use conditional MLE to estimate $\boldsymbol{\theta}$ and $\boldsymbol{\eta}$ : the log likelihood for observation $i$ is just the $\log$ of equation (22.31), evaluated at $a_{i}$. If we assume that $a_{i}$ is independent of $\mathbf{x}_{i}$ and has a uniform distribution on $[0, b]$, the estimation simplifies somewhat; see Problem 22.6. Allowing for a discontinuous starting time density $k\left(\cdot \mid \mathbf{x}_{i} ; \boldsymbol{\eta}\right)$ does not materially affect equation (22.31). For example, if the interval $[0,1]$ represents one year, we might want to allow different entry rates over the different seasons. This approach would correspond to a uniform distribution over each subinterval that we choose.

We now turn to the problem of left censoring, which arises with stock sampling when we do not actually know when any spell began. In other words, the $a_{i}$ are not observed, and therefore neither are the true durations, $t_{i}^{*}$. However, we assume that\\
we can follow spells after the interview date. Without right censoring, this assumption means we can observe the time in the current spell since the interview date, say, $r_{i}$, which we can write as $r_{i}=t_{i}^{*}+a_{i}-b$. We still have a left truncation problem because we only observe $r_{i}$ when $t_{i}^{*}>b-a_{i}$, that is, when $r_{i}>0$. The general approach is the same as with the earlier problems: we obtain the density of the variable that we can at least partially observe, $r_{i}$ in this case, conditional on observing $r_{i}$. Problem 22.8 asks you to fill in the details, accounting also for possible right censoring.

We can easily combine stock sampling and flow sampling. For example, in the case that we observe the starting times, $a_{i}$, suppose that, at time $m<b$, we sample a stock of individuals already in the initial state. In addition to following spells of individuals already in the initial state, suppose we can randomly sample individuals flowing into the initial state between times $m$ and $b$. Then we follow all the individuals appearing in the sample, at least until right censoring. For starting dates after $m\left(a_{i} \geq m\right)$, there is no truncation, and so the log likelihood for these observations is just as in equation (22.24). For $a_{i}<m$, the $\log$ likelihood is identical to equation (22.30) except that $m$ replaces $b$. Other combinations are easy to infer from the preceding results.

\subsection*{22.3.4 Unobserved Heterogeneity}
One way to obtain more general duration models is to introduce unobserved heterogeneity into fairly simple duration models. In addition, we sometimes want to test for duration dependence conditional on observed covariates and unobserved heterogeneity. The key assumptions used in most models that incorporate unobserved heterogeneity are that (1) the heterogeneity is independent of the observed covariates, as well as starting times and censoring times; (2) the heterogeneity has a distribution known up to a finite number of parameters; and (3) the heterogeneity enters the hazard function multiplicatively. We will make these assumptions. In the context of single-spell flow data, it is difficult to relax any of these assumptions. (In the special case of a lognormal duration distribution, we can relax assumption 1 by using Tobit methods with endogenous explanatory variables; see Section 17.5.2.)

In some fields, particularly those concerned with modeling survival times (such as biostatistics), the unobserved heterogeneity is called frailty. Then the hazard conditional on the unobserved frailty is the instantaneous probability of dying conditional on surviving up through time $t$ for an individual with a given frailty. In this chapter, we usually use the term heterogeneity. Often, a model that explicitly introduces heterogeneity into a hazard function is called a mixture model.

Before we cover the general case, it is useful to cover an example due to Lancaster (1979). For a random draw $i$ from the population, a Weibull hazard function conditional on observed covariates $\mathbf{x}_{i}$ and unobserved heterogeneity $v_{i}$ is


\begin{equation*}
\lambda\left(t ; \mathbf{x}_{i}, v_{i}\right)=v_{i} \exp \left(\mathbf{x}_{i} \boldsymbol{\beta}\right) \alpha t^{\alpha-1} \tag{22.33}
\end{equation*}


where $x_{i 1} \equiv 1$ and $v_{i}>0$. Lancaster (1990) calls equation (22.33) a conditional hazard, because it conditions on the unobserved heterogeneity (or frailty) $v_{i}$. Technically, almost all hazards in econometrics are conditional because we almost always condition on observed covariates. Notice how $v_{i}$ enters equation (22.33) multiplicatively. To identify the parameters $\alpha$ and $\boldsymbol{\beta}$ we need a normalization on the distribution of $v_{i}$; we use the most common, $\mathrm{E}\left(v_{i}\right)=1$. This implies that, for a given vector $\mathbf{x}$, the average hazard is $\exp (\mathbf{x} \boldsymbol{\beta}) \alpha t^{\alpha-1}$. An interesting hypothesis is $\mathrm{H}_{0}: \alpha=1$, which means that, conditional on $\mathbf{x}_{i}$ and $v_{i}$, there is no duration dependence.

In the general case where the cdf of $t_{i}^{*}$ given ( $\mathbf{x}_{i}, v_{i}$ ) is $F\left(t \mid \mathbf{x}_{i}, v_{i} ; \boldsymbol{\theta}\right)$, we can obtain the distribution of $t_{i}^{*}$ given $\mathbf{x}_{i}$ by integrating out the unobserved effect. Because $v_{i}$ and $\mathbf{x}_{i}$ are independent, the cdf of $t_{i}^{*}$ given $\mathbf{x}_{i}$ is


\begin{equation*}
G\left(t \mid \mathbf{x}_{i} ; \boldsymbol{\theta}, \boldsymbol{\rho}\right)=\int_{0}^{\infty} F\left(t \mid \mathbf{x}_{i}, v ; \boldsymbol{\theta}\right) h(v ; \boldsymbol{\rho}) d v \tag{22.34}
\end{equation*}


where, for concreteness, the density of $v_{i}, h(\cdot ; \boldsymbol{\rho})$, is assumed to be continuous and depends on the unknown parameters $\boldsymbol{\rho}$. From equation (22.34) the density of $t_{i}^{*}$ given $\mathbf{x}_{i}, g\left(t \mid \mathbf{x}_{i} ; \boldsymbol{\theta}, \boldsymbol{\rho}\right)$, is easily obtained. We can now use the methods of Sections 22.3.2 and 22.3.3. For flow data, the log-likelihood function is as in equation (22.24), but with $G\left(t \mid \mathbf{x}_{i} ; \boldsymbol{\theta}, \boldsymbol{\rho}\right)$ replacing $F\left(t \mid \mathbf{x}_{i} ; \boldsymbol{\theta}\right)$ and $g\left(t \mid \mathbf{x}_{i} ; \boldsymbol{\theta}, \boldsymbol{\rho}\right)$ replacing $f\left(t \mid \mathbf{x}_{i} ; \boldsymbol{\theta}\right)$. We should assume that $\mathrm{D}\left(t_{i}^{*} \mid \mathbf{x}_{i}, v_{i}, a_{i}, c_{i}\right)=\mathrm{D}\left(t_{i}^{*} \mid \mathbf{x}_{i}, v_{i}\right)$ and $\mathrm{D}\left(v_{i} \mid \mathbf{x}_{i}, a_{i}, c_{i}\right)=\mathrm{D}\left(v_{i}\right)$; these assumptions ensure that the key condition (22.22) holds. The methods for stock sampling described in Section 22.3.3 also apply to the integrated cdf and density.

If we assume gamma-distributed heterogeneity-that is, $v_{i} \sim \operatorname{Gamma}(\delta, \delta)$, so that $\mathrm{E}\left(v_{i}\right)=1$ and $\operatorname{Var}\left(v_{i}\right)=1 / \delta$-we can find the distribution of $t_{i}^{*}$ given $\mathbf{x}_{i}$ for a broad class of hazard functions with multiplicative heterogeneity. Suppose that the hazard function is $\lambda\left(t ; \mathbf{x}_{i}, v_{i}\right)=v_{i} \kappa\left(t ; \mathbf{x}_{i}\right)$, where $\kappa(t ; \mathbf{x})>0$ (and need not have the proportional hazard form). For simplicity, we suppress the dependence of $\kappa(\cdot ; \cdot)$ on unknown parameters. From equation (22.7), the $\operatorname{cdf}$ of $t_{i}^{*}$ given ( $\mathbf{x}_{i}, v_{i}$ ) is


\begin{equation*}
F\left(t \mid \mathbf{x}_{i}, v_{i}\right)=1-\exp \left[-v_{i} \int_{0}^{t} \kappa\left(s ; \mathbf{x}_{i}\right) d s\right] \equiv 1-\exp \left[-v_{i} \xi\left(t ; \mathbf{x}_{i}\right)\right] \tag{22.35}
\end{equation*}


where $\xi\left(t ; \mathbf{x}_{i}\right) \equiv \int_{0}^{t} \kappa\left(s ; \mathbf{x}_{i}\right) d s$. We can obtain the cdf of $t_{i}^{*}$ given $\mathbf{x}_{i}$ by using equation (22.34). The density of $v_{i}$ is $h(v)=\delta^{\delta} v^{\delta-1} \exp (-\delta v) / \Gamma(\delta)$, where $\operatorname{Var}\left(v_{i}\right)=1 / \delta$ and\\
$\Gamma(\cdot)$ is the gamma function. Let $\xi_{i} \equiv \xi\left(t ; \mathbf{x}_{i}\right)$ for given $t$. Then

$$
\begin{aligned}
& \int_{0}^{\infty} \exp \left(-\xi_{i} v\right) \delta^{\delta} v^{\delta-1} \exp (-\delta v) / \Gamma(\delta) d v \\
& \quad=\left[\delta /\left(\delta+\xi_{i}\right)\right]^{\delta} \int_{0}^{\infty}\left(\delta+\xi_{i}\right)^{\delta} v^{\delta-1} \exp \left[-\left(\delta+\xi_{i}\right) v\right] / \Gamma(\delta) d v \\
& \quad=\left[\delta /\left(\delta+\xi_{i}\right)\right]^{\delta}=\left(1+\xi_{i} / \delta\right)^{-\delta}
\end{aligned}
$$

where the second-to-last equality follows because the integrand is the Gamma ( $\delta$, $\delta+\xi_{i}$ ) density and must integrate to unity. Now we use equation (22.34):


\begin{equation*}
G\left(t \mid \mathbf{x}_{i}\right)=1-\left[1+\xi\left(t ; \mathbf{x}_{i}\right) / \delta\right]^{-\delta} \tag{22.36}
\end{equation*}


Taking the derivative of equation (22.36) with respect to $t$, using the fact that $\kappa\left(t ; \mathbf{x}_{i}\right)$ is the derivative of $\xi\left(t ; \mathbf{x}_{i}\right)$, yields the density of $t_{i}^{*}$ given $\mathbf{x}_{i}$ as


\begin{equation*}
g\left(t \mid \mathbf{x}_{i}\right)=\kappa\left(t ; \mathbf{x}_{i}\right)\left[1+\xi\left(t ; \mathbf{x}_{i}\right) / \delta\right]^{-\delta-1} \tag{22.37}
\end{equation*}


The function $\kappa(t ; \mathbf{x})$ depends on parameters $\boldsymbol{\theta}$, and so $g(t \mid \mathbf{x})$ should be $g(t \mid \mathbf{x} ; \boldsymbol{\theta}, \delta)$. With censored data the vector $\boldsymbol{\theta}$ can be estimated along with $\delta$ by using the loglikelihood function in equation (22.24) (again, with $G$ replacing $F$ ).

The hazard associated with the density $g(t \mid \mathbf{x})$ is typically called the unconditional hazard because it does not condition on the unobserved heterogeneity (but, as always, does condition on the observed covariates). It is useful to think of the unconditional hazard as something that can be estimated quite generally because it involves the distribution of an observed outcome conditional on observed covariates. In fact, without censoring, estimation is easily done without making parametric restrictions and without even introducing the notion of unobserved heterogeneity.

When the hazard function has the Weibull form in equation (22.33), $\xi(t ; \mathbf{x})= \exp (\mathbf{x} \boldsymbol{\beta}) t^{\alpha}$, which leads to a very tractable analysis when plugged into equations (22.36) and (22.37). The resulting duration distribution is called the Burr distribution. Its hazard function-that is, the unconditional hazard function when the conditional hazard is Weibull and the heterogeneity has a gamma distribution-is $\exp (\mathbf{x} \boldsymbol{\beta}) \alpha t^{\alpha-1}\left[1+\exp (\mathbf{x} \boldsymbol{\beta}) t^{\alpha} / \delta\right]^{-1}$. It is useful to reparameterize the Burr distribution by letting $\eta=1 / \delta$ and then writing the hazard as\\
$\exp (\mathbf{x} \boldsymbol{\beta}) \alpha t^{\alpha-1} /\left[1+\eta \exp (\mathbf{x} \boldsymbol{\beta}) t^{\alpha}\right]$\\
Then $\eta=0$-that is, $\operatorname{Var}\left(v_{i}\right)=0$-leads to the Weibull hazard, as expected. Further, $\eta=1$ (so that $\operatorname{Var}\left(v_{i}\right)=E\left(v_{i}\right)$ ) gives the log-logistic hazard in equation (22.27).

Therefore, even if we ignore that we derived the Burr distribution using gamma heterogeneity, we see that it nests two important special cases.

If we use a log-logistic hazard for $\kappa(t ; \mathbf{x})$-which, recall, implies that $\kappa(t ; \mathbf{x})$ is not multiplicatively separable in $t$ and $\mathbf{x}$-and assume gamma heterogeneity, then we plug $\kappa(t ; \mathbf{x})=\exp (\mathbf{x} \boldsymbol{\beta}) \alpha t^{\alpha-1}\left[1+\exp (\mathbf{x} \boldsymbol{\beta}) t^{\alpha}\right]^{-1}$ and $\xi(t ; \mathbf{x})=\log \left[1+\exp (\mathbf{x} \boldsymbol{\beta}) t^{\alpha}\right]$ into equation (22.37) to obtain the density (as a function of the parameters $\boldsymbol{\beta}, \alpha$, and $\delta$ ); we plug these expressions into equation (22.36) to obtain the cdf. Again, maximum likelihood estimation with right censoring is fairly straightforward.

Before presenting an example, we should recall why we might want to explicitly introduce unobserved heterogeneity when the heterogeneity is assumed to be independent of the observed covariates. The strongest case is seen when we are interested in testing for duration dependence conditional on observed covariates and unobserved heterogeneity, where the unobserved heterogeneity enters the hazard multiplicatively. As carefully exposited by Lancaster (1990, Section 10.2), ignoring multiplicative heterogeneity in the Weibull model results in asymptotically underestimating $\alpha$. Therefore, we could very well conclude that there is negative duration dependence conditional on $\mathbf{x}$, whereas there is no duration dependence ( $\alpha=1$ ) conditional on $\mathbf{x}$ and $v$ or even positive duration dependence ( $\alpha>1$ ).

In a general sense, it is somewhat heroic to think we can distinguish between duration dependence and unobserved heterogeneity when we observe only a single cycle for each agent. The problem is simple to describe: because we can only estimate the distribution of $T$ given $\mathbf{x}$, we cannot uncover the distribution of $T$ given ( $\mathbf{x}, v$ ) unless we make extra assumptions, a point Lancaster (1990, Section 10.1) illustrates with an example. Therefore, we cannot tell whether the hazard describing $T$ given $(\mathbf{x}, v)$ exhibits duration dependence. But, when the hazard has the proportional hazard form $\lambda(t ; \mathbf{x}, v)=v \kappa(\mathbf{x}) \lambda_{0}(t)$, it is possible to identify the function $\kappa(\cdot)$ and the baseline hazard $\lambda_{0}(\cdot)$ quite generally (along with the distribution of $v$ ). See Lancaster (1990, Section 7.3) for a presentation of the results of Elbers and Ridder (1982). More recently, Horowitz (1999) demonstrated how to nonparametrically estimate the baseline hazard and the distribution of the unobserved heterogeneity under fairly weak assumptions.

When interest centers on how the observed covariates affect the mean duration, explicitly modeling unobserved heterogeneity is less compelling. Adding unobserved heterogeneity to equation (22.26) does not change the mean effects; it merely changes the error distribution. Without censoring, we would probably estimate $\boldsymbol{\beta}$ in equation (22.26) by OLS (rather than MLE) so that the estimators would be robust to distributional misspecification. With censoring, to perform maximum likelihood, we must know the distribution of $t_{i}^{*}$ given $\mathbf{x}_{i}$, and this depends on the distribution of $v_{i}$

\begin{table}[h]
\begin{center}
\captionsetup{labelformat=empty}
\caption{Table 22.2\\
Weibull Hazard with Gamma Heterogeneity, Criminal Recidivism}
\begin{tabular}{|l|l|}
\hline
Explanatory Variable & Coefficient (Standard Error) \\
\hline
workprg & \begin{tabular}{l}
. 007 \\
(.204) \\
\end{tabular} \\
\hline
priors & \begin{tabular}{l}
. 243 \\
(.042) \\
\end{tabular} \\
\hline
tserved & \begin{tabular}{l}
. 035 \\
(.007) \\
\end{tabular} \\
\hline
felon & \begin{tabular}{l}
-. 791 \\
(.267) \\
\end{tabular} \\
\hline
alcohol & \begin{tabular}{l}
1.174 \\
(.281) \\
\end{tabular} \\
\hline
drugs & \begin{tabular}{l}
. 285 \\
(.223) \\
\end{tabular} \\
\hline
black & \begin{tabular}{l}
. 772 \\
(.204) \\
\end{tabular} \\
\hline
married & \begin{tabular}{l}
-. 806 \\
(.258) \\
\end{tabular} \\
\hline
educ & \begin{tabular}{l}
-. 027 \\
(.045) \\
\end{tabular} \\
\hline
age & \begin{tabular}{l}
-. 005 \\
(.001) \\
\end{tabular} \\
\hline
constant & \begin{tabular}{l}
-5.394 \\
(.720) \\
\end{tabular} \\
\hline
$\hat{\alpha}$ & \begin{tabular}{l}
1.708 \\
(.162) \\
\end{tabular} \\
\hline
$\hat{\delta}$ & \begin{tabular}{l}
5.991 \\
(1.071) \\
\end{tabular} \\
\hline
Observations & 1,445 \\
\hline
Log likelihood & -1,584.92 \\
\hline
\end{tabular}
\end{center}
\end{table}

when we explicitly introduce unobserved heterogeneity. But again introducing unobserved heterogeneity is indistinguishable from simply allowing a more flexible duration distribution.

We now apply a model with a Weibull baseline hazard and gamma heterogeneity to the recidivism data. Individuals with a larger heterogeneity, $v_{i}$, have a higher probability of being arrested after release from prison in every interval, conditional on "surviving" up to that point.

The maximum likelihood estimates are given in Table 22.2. Notice how the estimate of $\alpha, 1.708$, is well above unity. Therefore, the conditional hazard exhibits

\begin{figure}[h]
\begin{center}
  \includegraphics[alt={},max width=\textwidth]{2e17339c-7635-4ef8-aded-2191d8d853ff-1039_917_1258_325_460}
\captionsetup{labelformat=empty}
\caption{Figure 22.3\\
Conditional Weibull hazard}
\end{center}
\end{figure}

positive duration dependence: for a given individual, the instantaneous probability of being arrested after prison release monotonically increases with the time out of prison. The situation is displayed in Figure 22.3. By contrast, in the Weibull model without heterogeneity, we estimated $\alpha$ to be .806 .

Allowing for heterogeneity effects an important change on the shape of the hazard. The qualitative effects of the covariates are similar to those in Table 22.1, although the magnitudes of some of the variables-for example, the number of priors, the felon dummy variable, and the marriage dummy variable-increase nontrivially. The estimate of $\delta$ is about 6.0 , and it is very statistically significant. Thus the Weibull model without heterogeneity is strongly rejected in favor of the Weibull model with heterogeneity.

The unconditional hazard-that is, the hazard for the duration distribution that integrates out the heterogeneity-is plotted in Figure 22.4. The covariates are set at mean values. Its shape is more like that in the lognormal model, although the peak

\begin{figure}[h]
\begin{center}
  \includegraphics[alt={},max width=\textwidth]{2e17339c-7635-4ef8-aded-2191d8d853ff-1040_914_1256_326_288}
\captionsetup{labelformat=empty}
\caption{Figure 22.4\\
Unconditional Burr hazard}
\end{center}
\end{figure}

of the hazard in Figure 22.4 is at roughly 12 months, rather than six months. The monotonicity of the conditional hazard coupled with the hump shape of the unconditional hazard is common in applications. The reasoning behind the hump shape for the unconditional hazard goes something like this. (For simplicity, assume that there are no observed covariates, or that we are conditioning on a set of values.) Initially, all types of men are in the risk set, and so the shape of the unconditional hazard-that is, the aggregate hazard across the entire population-mimics that of the conditional hazard. But in the early periods, men with high proclivities to commit crimes-that is, with higher $v_{i}$-will tend to be arrested. Those arrested drop out of the risk set in subsequent time periods. Therefore, the men left in the risk set in later time periods tend to have the lower predispositions to repeat crimes. As the length of time until arrest increases, men with smaller values of $v_{i}$ are left in the risk set, and this result translates into a declining unconditional hazard.

As with all applications that assume the proportional hazard form (in both observed covariates and unobserved heterogeneity), we are left with a conundrum: should we accept that the conditional hazard has the shape in Figure 22.3-for any heterogeneity and any value of $\mathbf{x}$-or should the conditional hazard be more complicated? As mentioned earlier, unless we take the multiplicative structure $\lambda\left(t ; \mathbf{x}_{i}, v_{i}\right)=v_{i} \kappa\left(\mathbf{x}_{i}\right) \lambda_{0}(t)$ as given, there is a fundamental lack of identification. If we replace $\kappa(\mathbf{x}) \lambda_{0}(t)$ with the general nonseparable function $\kappa(t ; \mathbf{x})$, identification of $\kappa(\cdot ; \mathbf{x})$ for given values of $\mathbf{x}$ becomes very difficult. Further, more complicated ways of introducing heterogeneity lead into very rough waters. For example, suppose we specify $\lambda\left(t ; \mathbf{x}_{i}, \mathbf{b}_{i}\right)=\exp \left(\mathbf{x}_{i} \mathbf{b}_{i}\right) \lambda_{0}(t)$, where $\mathbf{b}_{i}$ is a $K \times 1$ random vector. (Letting $x_{i 1} \equiv 1$ encompasses the usual case of multiplicative heterogeneity.) Even if we assume $\mathbf{b}_{i}$ is independent of $\mathbf{x}_{i}$, it is not clear that one can identify $\lambda_{0}(t)$ without assuming a full distribution for $\mathbf{b}_{i}$. Even if we took such an approach, we would never be able to distinguish between relatively simple models with unobserved heterogeneity and complicated hazards that cannot be written as a multiple of $\lambda_{0}(t)$.

\subsection*{22.4 Analysis of Grouped Duration Data}
Continuously distributed durations are, strictly speaking, rare in social science applications. Even if an underlying duration is properly viewed as being continuous, measurements are necessarily discrete. When the measurements are fairly precise, it is sensible to treat the durations as continuous random variables. But when the measurements are coarse-such as monthly, or perhaps even weekly-it can be important to account for the discreteness in the estimation.

Grouped duration data arise when each duration is only known to fall into a certain time interval, such as a week, a month, or even a year. For example, unemployment durations are often measured to the nearest week. In Example 22.2 the time until next arrest is measured to the nearest month. Even with grouped data we can generally estimate the parameters of the duration distribution.

The approach we take here to analyzing grouped data summarizes the information on staying in the initial state or exiting in each time interval in a sequence of binary outcomes. (Kiefer, 1988; Han and Hausman, 1990; Meyer, 1990; Lancaster, 1990; McCall, 1994; and Sueyoshi, 1995, all take this approach.) In effect, we have a panel data set where each cross section observation is a vector of binary responses, along with covariates. In addition to allowing us to treat grouped durations, the panel data approach has at least two additional advantages. First, in a proportional hazard specification, it leads to easy methods for estimating flexible hazard functions. Sec-\\
ond, because of the sequential nature of the data, time-varying covariates are easily introduced.

We assume flow sampling so that we do not have to address the sample selection problem that arises with stock sampling. We divide the time line into $M+1$ intervals, $\left[0, a_{1}\right),\left[a_{1}, a_{2}\right), \ldots,\left[a_{M-1}, a_{M}\right),\left[a_{M}, \infty\right)$, where the $a_{m}$ are known constants. For example, we might have $a_{1}=1, a_{2}=2, a_{3}=3$, and so on, but unequally spaced intervals are allowed. The last interval, [ $a_{M}, \infty$ ), is chosen so that any duration falling into it is censored at $a_{M}$ : no observed durations are greater than $a_{M}$. For a random draw from the population, let $c_{m}$ be a binary censoring indicator equal to unity if the duration is censored in interval $m$, and zero otherwise. Notice that $c_{m}=1$ implies $c_{m+1}=1$ : if the duration was censored in interval $m$, it is still censored in interval $m+1$. Because durations lasting into the last interval are censored, $c_{M+1} \equiv 1$. Similarly, $y_{m}$ is a binary indicator equal to unity if the duration ends in the $m$ th interval and zero otherwise. Thus, $y_{m+1}=1$ if $y_{m}=1$. If the duration is censored in interval $m\left(c_{m}=1\right)$, we set $y_{m} \equiv 1$ by convention.

As in Section 22.3, we allow individuals to enter the initial state at different calendar times. In order to keep the notation simple, we do not explicitly show the conditioning on these starting times, as the starting times play no role under flow sampling when we assume that, conditional on the covariates, the starting times are independent of the duration and any unobserved heterogeneity. If necessary, startingtime dummies can be included in the covariates.

For each person $i$, we observe $\left(y_{i 1}, c_{i 1}\right), \ldots,\left(y_{i M}, c_{i M}\right)$, which is a balanced panel data set. To avoid confusion with our notation for a duration ( $T$ for the random variable, $t$ for a particular outcome on $T$ ), we use $m$ to index the time intervals. The string of binary indicators for any individual is not unrestricted: we must observe a string of zeros followed by a string of ones. The important information is the interval in which $y_{\text {im }}$ becomes unity for the first time, and whether that represents a true exit from the initial state or censoring.

\subsection*{22.4.1 Time-Invariant Covariates}
With time-invariant covariates, each random draw from the population consists of information on $\left\{\left(y_{1}, c_{1}\right), \ldots,\left(y_{M}, c_{M}\right), \mathbf{x}\right\}$. We assume that a parametric hazard function is specified as $\lambda(t ; \mathbf{x}, \boldsymbol{\theta})$, where $\boldsymbol{\theta}$ is the vector of unknown parameters. Let $T$ denote the time until exit from the initial state. While we do not fully observe $T$, either we know which interval it falls into, or we know whether it was censored in a particular interval. This knowledge is enough to obtain the probability that $y_{m}$ takes on the value unity given $\left(y_{m-1}, \ldots, y_{1}\right),\left(c_{m}, \ldots, c_{1}\right)$, and $\mathbf{x}$. In fact, by definition this\\
probability depends only on $y_{m-1}, c_{m}$, and $\mathbf{x}$, and only two combinations yield probabilities that are not identically zero or one. These probabilities are


\begin{equation*}
\mathrm{P}\left(y_{m}=0 \mid y_{m-1}=0, \mathbf{x}, c_{m}=0\right) \tag{22.38}
\end{equation*}



\begin{equation*}
\mathrm{P}\left(y_{m}=1 \mid y_{m-1}=0, \mathbf{x}, c_{m}=0\right), \quad m=1, \ldots, M \tag{22.39}
\end{equation*}


(We define $y_{0} \equiv 0$ so that these equations hold for all $m \geq 1$.) To compute these probabilities in terms of the hazard for $T$, we assume that the duration is conditionally independent of censoring:\\
$T$ is independent of $c_{1}, \ldots, c_{M}$, given $\mathbf{x}$

This assumption allows the censoring to depend on $\mathbf{x}$ but rules out censoring that depends on unobservables, after conditioning on $\mathbf{x}$. Condition (22.40) holds for fixed censoring or completely randomized censoring. (It may not hold if censoring is due to nonrandom attrition.) Under assumption (22.40) we have, from equation (22.9),


\begin{align*}
\mathrm{P}\left(y_{m}=1 \mid y_{m-1}=0, \mathbf{x}, c_{m}=0\right) & =\mathrm{P}\left(a_{m-1} \leq T<a_{m} \mid T \geq a_{m-1}, \mathbf{x}\right) \\
& =1-\exp \left[-\int_{a_{m-1}}^{a_{m}} \lambda(s ; \mathbf{x}, \boldsymbol{\theta}) d s\right] \equiv 1-\alpha_{m}(\mathbf{x}, \boldsymbol{\theta}) \tag{22.41}
\end{align*}


for $m=1,2, \ldots, M$, where


\begin{equation*}
\alpha_{m}(\mathbf{x}, \boldsymbol{\theta}) \equiv \exp \left[-\int_{a_{m-1}}^{a_{m}} \lambda(s ; \mathbf{x}, \boldsymbol{\theta}) d s\right] \tag{22.42}
\end{equation*}


Therefore,


\begin{equation*}
\mathrm{P}\left(y_{m}=0 \mid y_{m-1}=0, \mathbf{x}, c_{m}=0\right)=\alpha_{m}(\mathbf{x}, \boldsymbol{\theta}) \tag{22.43}
\end{equation*}


We can use these probabilities to construct the likelihood function. If, for observation $i$, uncensored exit occurs in interval $m_{i}$, the likelihood is


\begin{equation*}
\left[\prod_{h=1}^{m_{i}-1} \alpha_{h}\left(\mathbf{x}_{i}, \boldsymbol{\theta}\right)\right]\left[1-\alpha_{m_{i}}\left(\mathbf{x}_{i}, \boldsymbol{\theta}\right)\right] \tag{22.44}
\end{equation*}


The first term represents the probability of remaining in the initial state for the first $m_{i}-1$ intervals, and the second term is the (conditional) probability that $T$ falls into interval $m_{i}$. (Because an uncensored duration must have $m_{i} \leq M$, expression (22.44) at most depends on $\alpha_{1}\left(\mathbf{x}_{i}, \boldsymbol{\theta}\right), \ldots, \alpha_{M}\left(\mathbf{x}_{i}, \boldsymbol{\theta}\right)$.) If the duration is censored in interval\\
$m_{i}$, we know only that exit did not occur in the first $m_{i}-1$ intervals, and the likelihood consists of only the first term in expression (22.44).

If $d_{i}$ is a censoring indicator equal to one if duration $i$ is uncensored, the log likelihood for observation $i$ can be written as


\begin{equation*}
\sum_{h=1}^{m_{i}-1} \log \left[\alpha_{h}\left(\mathbf{x}_{i}, \boldsymbol{\theta}\right)\right]+d_{i} \log \left[1-\alpha_{m_{i}}\left(\mathbf{x}_{i}, \boldsymbol{\theta}\right)\right] \tag{22.45}
\end{equation*}


The log likelihood for the entire sample is obtained by summing expression (22.45) across all $i=1, \ldots, N$. Under the assumptions made, this log likelihood represents the density of $\left(y_{1}, \ldots, y_{M}\right)$ given $\left(c_{1}, \ldots, c_{M}\right)$ and $\mathbf{x}$, and so the conditional maximum likelihood theory covered in Chapter 13 applies directly. The various ways of estimating asymptotic variances and computing test statistics are available.

To implement conditional MLE, we must specify a hazard function. One hazard function that has become popular because of its flexibility is a piecewise-constant proportional hazard: for $m=1, \ldots, M$,


\begin{equation*}
\lambda(t ; \mathbf{x}, \boldsymbol{\theta})=\kappa(\mathbf{x}, \boldsymbol{\beta}) \lambda_{m}, \quad a_{m-1} \leq t<a_{m} \tag{22.46}
\end{equation*}


where $\kappa(\mathbf{x}, \boldsymbol{\beta})>0$ (and typically $\kappa(\mathbf{x}, \boldsymbol{\beta})=\exp (\mathbf{x} \boldsymbol{\beta})$ ). This specification allows the hazard to be different (albeit constant) over each time interval. The parameters to be estimated are $\boldsymbol{\beta}$ and $\lambda$, where the latter is the vector of $\lambda_{m}, m=1, \ldots, M$. (Because durations in $\left[a_{M}, \infty\right)$ are censored at $a_{M}$, we cannot estimate the hazard over the interval $\left[a_{M}, \infty\right)$.) As an example, if we have unemployment duration measured in weeks, the hazard can be different in each week. If the durations are sparse, we might assume a different hazard rate for every two or three weeks (this assumption places restrictions on the $\left.\lambda_{m}\right)$. With the piecewise-constant hazard and $\kappa(\mathbf{x}, \boldsymbol{\beta})=\exp (\mathbf{x} \boldsymbol{\beta})$, for $m=1, \ldots, M$, we have


\begin{equation*}
\alpha_{m}(\mathbf{x}, \boldsymbol{\theta}) \equiv \exp \left[-\exp (\mathbf{x} \boldsymbol{\beta}) \lambda_{m}\left(a_{m}-a_{m-1}\right)\right] \tag{22.47}
\end{equation*}


Remember, the $a_{m}$ are known constants (often $a_{m}=m$ ) and not parameters to be estimated. Usually the $\lambda_{m}$ are unrestricted, in which case $\mathbf{x}$ does not contain an intercept.

The piecewise-constant hazard implies that the duration distribution is discontinuous at the endpoints, whereas in our discussion in Section 22.2 we assumed that the duration had a continuous distribution. A piecewise-continuous distribution causes no real problems, and the log likelihood is exactly as specified previously. Alternatively, as in Han and Hausman (1990) and Meyer (1990), we can assume that $T$\\
has a proportional hazard as in equation (22.16) with continuous baseline hazard, $\lambda_{0}(\cdot)$. Then, we can estimate $\boldsymbol{\beta}$ along with the parameters\\
$\int_{a_{m-1}}^{a_{m}} \lambda_{0}(s) d s, \quad m=1,2, \ldots, M$\\
In practice, the approaches are the same, and it is easiest to just assume a piecewiseconstant proportional hazard, as in equation (22.46).

Once the $\lambda_{m}$ have been estimated along with $\boldsymbol{\beta}$, an estimated hazard function is easily plotted: graph $\hat{\lambda}_{m}$ at the midpoint of the interval $\left[a_{m-1}, a_{m}\right)$, and connect the points.

Without covariates, maximum likelihood estimation of the $\lambda_{m}$ leads to a wellknown estimator of the survivor function. Rather than derive the MLE of the survivor function, it is easier to motivate the estimator from the representation of the survivor function as a product of conditional probabilities. For $m=1, \ldots, M$, the survivor function at time $a_{m}$ can be written as


\begin{equation*}
S\left(a_{m}\right)=\mathrm{P}\left(T>a_{m}\right)=\prod_{r=1}^{m} \mathrm{P}\left(T>a_{r} \mid T>a_{r-1}\right) \tag{22.48}
\end{equation*}


(Because $a_{0}=0$ and $\mathrm{P}(T>0)=1$, the $r=1$ term on the right-hand side of equation (22.48) is simply $\mathrm{P}\left(T>a_{1}\right)$.) Now, for each $r=1,2, \ldots, M$, let $N_{r}$ denote the number of people in the risk set for interval $r: N_{r}$ is the number of people who have neither left the initial state nor been censored at time $a_{r-1}$, which is the beginning of interval $r$. Therefore, $N_{1}$ is the number of individuals in the initial random sample; $N_{2}$ is the number of individuals who did not exit the initial state in the first interval, less the number of individuals censored in the first interval; and so on. Let $E_{r}$ be the number of people observed to leave in the $r$ th interval-that is, in the interval $\left[a_{r-1}, a_{r}\right)$. A consistent estimator of $\mathrm{P}\left(T>a_{r} \mid T>a_{r-1}\right)$ is $\left(N_{r}-E_{r}\right) / N_{r}, r=1,2, \ldots, M$. (We must use the fact that the censoring is ignorable in the sense of assumption (22.40), so that there is no sample selection bias in using only the uncensored observations.) It follows from equation (22.48) that a consistent estimator of the survivor function at time $a_{m}$ is


\begin{equation*}
\hat{S}\left(a_{m}\right)=\prod_{r=1}^{m}\left[\left(N_{r}-E_{r}\right) / N_{r}\right], \quad m=1,2, \ldots, M \tag{22.49}
\end{equation*}


This is the Kaplan-Meier estimator of the survivor function (at the points $a_{1}, a_{2}$, ..., $a_{M}$ ). Lancaster (1990, Section 8.2) contains a proof that maximum likelihood\\
estimation of the $\lambda_{m}$ (without covariates) leads to the Kaplan-Meier estimator of the survivor function. If there are no censored durations before time $a_{m}, \hat{S}\left(a_{m}\right)$ is simply the fraction of people who have not left the initial state at time $a_{m}$, which is obviously consistent for $\mathrm{P}\left(T>a_{m}\right)=S\left(a_{m}\right)$.

In the general model, we do not need to assume a proportional hazard specification within each interval. For example, we could assume a log-logistic hazard within each interval, with different parameters for each $m$. Because the hazard in such cases does not depend on the covariates multiplicatively, we must plug in values of $\mathbf{x}$ in order to plot the hazard. Sueyoshi (1995) studies such models in detail.

If the intervals $\left[a_{m-1}, a_{m}\right)$ are coarser than the data-for example, unemployment is measured in weeks, but we choose $\left[a_{m-1}, a_{m}\right)$ to be four weeks for all $m$-then we can specify nonconstant hazards within each interval. The piecewise-constant hazard corresponds to an exponential distribution within each interval. But we could specify, say, a Weibull distribution within each interval. See Sueyoshi (1995) for details.

\subsection*{22.4.2 Time-Varying Covariates}
Deriving the log likelihood is more complicated with time-varying covariates, especially when we do not assume that the covariates are strictly exogenous. Nevertheless, we will show that, if the covariates are constant within each time interval $\left[a_{m-1}, a_{m}\right)$, the form of the log likelihood is the same as expression (22.45), provided $\mathbf{x}_{i}$ is replaced with $\mathbf{x}_{i m}$ in interval $m$.

For the population, let $\mathbf{x}_{1}, \mathbf{x}_{2}, \ldots, \mathbf{x}_{M}$ denote the outcomes of the covariates in each of the $M$ time intervals, where we assume that the covariates are constant within an interval. This assumption is clearly an oversimplification, but we cannot get very far without it (and it reflects how data sets with time-varying covariates are usually constructed). When the covariates are internal and are not necessarily defined after exit from the initial state, the definition of the covariates in the time intervals is irrelevant; but it is useful to list covariates for all $M$ time periods.

We assume that the hazard at time $t$ conditional on the covariates up through time $t$ depends only on the covariates at time $t$. If past values of the covariates matter, they can simply be included in the covariates at time $t$. The conditional independence assumption on the censoring indicators is now stated as


\begin{equation*}
\mathrm{D}\left(T \mid T \geq a_{m-1}, \mathbf{x}_{m}, c_{m}\right)=\mathrm{D}\left(T \mid T \geq a_{m-1}, \mathbf{x}_{m}\right), \quad m=1, \ldots, M \tag{22.50}
\end{equation*}


This assumption allows the censoring decision to depend on the covariates during the time interval (as well as past covariates, provided they are either included in $\mathbf{x}_{m}$ or do not affect the distribution of $T$ given $\mathbf{x}_{m}$ ). Under this assumption, the probability of exit (without censoring) is


\begin{align*}
\mathrm{P}\left(y_{m}=1 \mid y_{m-1}=0, \mathbf{x}_{m}, c_{m}=0\right) & =\mathrm{P}\left(a_{m-1} \leq T<a_{m} \mid T \geq a_{m-1}, \mathbf{x}_{m}\right) \\
& =1-\exp \left[-\int_{a_{m-1}}^{a_{m}} \lambda\left(s ; \mathbf{x}_{m}, \boldsymbol{\theta}\right) d s\right] \equiv 1-\alpha_{m}\left(\mathbf{x}_{m}, \boldsymbol{\theta}\right) \tag{22.51}
\end{align*}


We can use equation (22.51), along with $\mathrm{P}\left(y_{m}=0 \mid y_{m-1}=0, \mathbf{x}_{m}, c_{m}=0\right)=\alpha_{m}\left(\mathbf{x}_{m}, \boldsymbol{\theta}\right)$, to build up a partial $\log$ likelihood for person $i$. As we discussed in Section 13.8, this is only a partial likelihood because we are not necessarily modeling the joint distribution of $\left(y_{1}, \ldots, y_{M}\right)$ given $\left\{\left(\mathbf{x}_{1}, c_{1}\right), \ldots,\left(c_{M}, \mathbf{x}_{M}\right)\right\}$.

For someone censored in interval $m$, the information on the duration is contained in $y_{i 1}=0, \ldots, y_{i, m-1}=0$. For someone who truly exits in interval $m$, there is additional information in $y_{i m}=1$. Therefore, the partial log likelihood is given by expression (22.45), but, to reflect the time-varying covariates, $\alpha_{h}\left(\mathbf{x}_{i}, \boldsymbol{\theta}\right)$ is replaced by $\alpha_{h}\left(\mathbf{x}_{i h}, \boldsymbol{\theta}\right)$ and $\alpha_{m_{i}}\left(\mathbf{x}_{i}, \boldsymbol{\theta}\right)$ is replaced by $\alpha_{m_{i}}\left(\mathbf{x}_{i, m_{i}}, \boldsymbol{\theta}\right)$.

Each term in the partial $\log$ likelihood represents the distribution of $y_{m}$ given $\left(y_{m-1}, \ldots, y_{1}\right),\left(\mathbf{x}_{m}, \ldots, \mathbf{x}_{1}\right)$, and $\left(c_{m}, \ldots, c_{1}\right)$. (Most of the probabilities in this conditional distribution are either zero or one; only the probabilities that depend on $\boldsymbol{\theta}$ are shown in expression (22.45).) Therefore, the density is dynamically complete, in the terminology of Section 13.8.3. As shown there, the usual maximum likelihood variance matrix estimators and statistics are asymptotically valid, even though we need not have the full conditional distribution of $\mathbf{y}$ given ( $\mathbf{x}, \mathbf{c}$ ). This result would change if, for some reason, we chose not to include past covariates when in fact they affect the current probability of exit even after conditioning on the current covariates. Then the robust forms of the statistics covered in Section 13.8 should be used. In most duration applications we want dynamic completeness.

If the covariates are strictly exogenous and if the censoring is strictly exogenous, then the partial likelihood is the full conditional likelihood. The precise strict exogeneity assumption is


\begin{equation*}
\mathrm{D}\left(T \mid T \geq a_{m-1}, \mathbf{x}, \mathbf{c}\right)=\mathrm{D}\left(T \mid T \geq a_{m-1}, \mathbf{x}_{m}\right), \quad m=1, \ldots, M \tag{22.52}
\end{equation*}


where $\mathbf{x}$ is the vector of covariates across all time periods and $\mathbf{c}$ is the vector of censoring indicators. There are two parts to this assumption. Ignoring the censoring, assumption (22.52) means that neither future nor past covariates appear in the hazard, once current covariates are controlled for. The second implication of assumption (22.52) is that the censoring is also strictly exogenous.

With time-varying covariates, the hazard specification


\begin{equation*}
\lambda\left(t ; \mathbf{x}_{m}, \boldsymbol{\theta}\right)=\kappa\left(\mathbf{x}_{m}, \boldsymbol{\beta}\right) \lambda_{m}, \quad a_{m-1} \leq t<a_{m} \tag{22.53}
\end{equation*}


$m=1, \ldots, M$, is still attractive. It implies that the covariates have a multiplicative effect in each time interval, and it allows the baseline hazard-the part common to all members of the population-to be flexible.

Meyer (1990) essentially uses the specification (22.53) to estimate the effect of unemployment insurance on unemployment spells. McCall (1994) shows how to allow for time-varying coefficients when $\kappa\left(\mathbf{x}_{m}, \boldsymbol{\beta}\right)=\exp \left(\mathbf{x}_{m} \boldsymbol{\beta}\right)$. In other words, $\boldsymbol{\beta}$ is replaced with $\boldsymbol{\beta}_{m}, m=1, \ldots, M$.

\subsection*{22.4.3 Unobserved Heterogeneity}
We can also add unobserved heterogeneity to hazards specified for grouped data, even if we have time-varying covariates. With time-varying covariates and unobserved heterogeneity, it is difficult to relax the strict exogeneity assumption. Also, with single-spell data, we cannot allow general correlation between the unobserved heterogeneity and the covariates. Therefore, we assume that the covariates are strictly exogenous conditional on unobserved heterogeneity and that the unobserved heterogeneity is independent of the covariates.

The precise assumptions are given by equation (22.52) but where unobserved heterogeneity, $v$, appears in both conditioning sets. In addition, we assume that $v$ is independent of ( $\mathbf{x}, \mathbf{c}$ ) (which is a further sense in which the censoring is exogenous).

In the leading case of the piecewise-constant baseline hazard, equation (22.53) becomes


\begin{equation*}
\lambda\left(t ; v, \mathbf{x}_{m}, \boldsymbol{\theta}\right)=v \kappa\left(\mathbf{x}_{m}, \boldsymbol{\beta}\right) \lambda_{m}, \quad a_{m-1} \leq t<a_{m} \tag{22.54}
\end{equation*}


where $v>0$ is a continuously distributed heterogeneity term. Using the same reasoning as in Sections 22.4.1 and 22.4.2, the density of $\left(y_{i 1}, \ldots, y_{i M}\right)$ given $\left(v_{i}, \mathbf{x}_{i}, \mathbf{c}_{i}\right)$ is


\begin{equation*}
\left[\prod_{h=1}^{m_{i}-1} \alpha_{h}\left(v_{i}, \mathbf{x}_{i h}, \boldsymbol{\theta}\right)\right]\left[1-\alpha_{m_{i}}\left(v_{i}, \mathbf{x}_{i, m_{i}}, \boldsymbol{\theta}\right)\right]^{d_{i}} \tag{22.55}
\end{equation*}


where $d_{i}=1$ if observation $i$ is uncensored. Because expression (22.55) depends on the unobserved heterogeneity, $v_{i}$, we cannot use it directly to consistently estimate $\boldsymbol{\theta}$. However, because $v_{i}$ is independent of $\left(\mathbf{x}_{i}, \mathbf{c}_{i}\right)$, with density $g(v ; \boldsymbol{\delta})$, we can integrate expression (22.55) against $g(\cdot ; \boldsymbol{\delta})$ to obtain the density of $\left(y_{i 1}, \ldots, y_{i M}\right)$ given $\left(\mathbf{x}_{i}, \mathbf{c}_{i}\right)$. This density depends on the observed data- $\left(m_{i}, d_{i}, \mathbf{x}_{i}\right)$-and the parameters $\boldsymbol{\theta}$ and $\boldsymbol{\delta}$. From this density, we construct the conditional log likelihood for observation $i$, and we can obtain the conditional MLE, just as in other nonlinear models with unobserved heterogeneity-see Chapters 15-19. Meyer (1990) assumes that the distribution of $v_{i}$ is gamma, with unit mean, and obtains the log-likelihood function in closed\\
form. McCall (1994) analyzes a heterogeneity distribution that contains the gamma as a special case.

It is possible to consistently estimate $\boldsymbol{\beta}$ and $\boldsymbol{\lambda}$ without specifying a parametric form for the heterogeneity distribution; this approach results in a semiparametric maximum likelihood estimator. Heckman and Singer (1984) first showed how to perform this method with a Weibull baseline hazard, and Meyer (1990) proved consistency when the hazard has the form (22.54). The estimated heterogeneity distribution is discrete and, in practice, has relatively few mass points. The consistency argument works by allowing the number of mass points to increase with the sample size. Computation is a difficult issue, and the asymptotic distribution of the semiparametric maximum likelihood estimator has not been worked out.

\subsection*{22.5 Further Issues}
The methods we have covered in this chapter have been applied in many contexts. Nevertheless, there are several important topics that we have neglected.

\subsection*{22.5.1 Cox's Partial Likelihood Method for the Proportional Hazard Model}
Cox (1972) suggested a partial likelihood method for estimating the parameters $\boldsymbol{\beta}$ in a proportional hazard model without specifying the baseline hazard. The strength of Cox's approach is that the effects of the covariates can be estimated very generally, provided the hazard is of the form (22.16). However, Cox's method is intended to be applied to flow data as opposed to grouped data. If we apply Cox's methods to grouped data, we must confront the practically important issue of individuals with identical observed durations. In addition, with time-varying covariates, Cox's method evidently requires the covariates to be strictly exogenous. Estimation of the hazard function itself is more complicated than the methods for grouped data that we covered in Section 22.4. See Amemiya (1985, Chapter 11) and Lancaster (1990, Chapter 9) for treatments of Cox's partial likelihood estimator.

\subsection*{22.5.2 Multiple-Spell Data}
All the methods we have covered assume a single spell for each sample unit. In other words, each individual begins in the initial state and then either is observed leaving the state or is censored. But at least some individuals might have multiple spells, especially if we follow them for long periods. For example, we may observe a person who is initially unemployed, becomes employed, and then after a time becomes\\
unemployed again. If we assume constancy across time about the process driving unemployment duration, we can use multiple spells to aid in identification, particularly in models with heterogeneity that can be correlated with time-varying covariates. Chamberlain (1985) and Honoré (1993b) contain identification results when multiple spells are observed. Chamberlain allowed for correlation between the heterogeneity and the time-varying covariates.

Multiple-spell data are also useful for estimating models with unobserved heterogeneity when the regressors are not strictly exogenous. Ham and Lalonde (1996) give an example in which participation in a job training program can be related to past unemployment duration, even though eligibility is randomly assigned. See also Wooldridge (2000) for a general framework that allows feedback to future explanatory variables in models with unobserved heterogeneity.

\subsection*{22.5.3 Competing Risks Models}
Another important topic is allowing for more than two possible states. Competing risks models allow for the possibility that an individual may exit into different alternatives. For example, a person working full-time may choose to retire completely or work part-time. Han and Hausman (1990) and Sueyoshi (1992) contain discussions of the assumptions needed to estimate competing risks models, with and without unobserved heterogeneity. See van den Berg (2001) for a detailed treatment.

\section*{Problems}
22.1. Use the data in RECID.RAW for this problem.\\
a. Using the covariates in Table 22.1, estimate equation (22.26) by censored Tobit. Verify that the log-likelihood value is $-1,597.06$.\\
b. Plug in the mean values for priors, tserved, educ, and age, and the values workprg $=0$, felon $=1$, alcohol $=1$, drugs $=1$, black $=0$, and married $=0$, and plot the estimated hazard for the lognormal distribution. Describe what you find.\\
c. Using only the uncensored observations, perform an OLS regression of $\log ($ durat $)$ on the covariates in Table 22.1. Compare the estimates on tserved and alcohol with those from part a. What do you conclude?\\
d. Now compute an OLS regression using all data-that is, treat the censored observations as if they are uncensored. Compare the estimates on tserved and alcohol from those in parts a and c.\\
22.2 Use the data in RECID.RAW to answer these questions:\\
a. To the Weibull model, add the variables super $(=1$ if release from prison was supervised) and rules (number of rules violations while in prison). Do the coefficient estimates on these new variables have the expected signs? Are they statistically significant?\\
b. Add super and rules to the lognormal model, and answer the same questions as in part a.\\
c. Compare the estimated effects of the rules variable on the expected duration for the Weibull and lognormal models. Are they practically different?\\
22.3. Consider the case of flow sampling, as in Section 22.3.2, but suppose that all durations are censored: $d_{i}=1, i=1, \ldots, N$.\\
a. Write down the log-likelihood function when all durations are censored.\\
b. Find the special case of the Weibull distribution in part a.\\
c. Consider the Weibull case where $\mathbf{x}_{i}$ only contains a constant, so that $F(t ; \alpha, \beta)= 1-\exp \left[-\exp (\beta) t^{\alpha}\right]$. Show that the Weibull $\log$ likelihood cannot be maximized for real numbers $\hat{\beta}$ and $\hat{\alpha}$.\\
d. From part c, what do you conclude about estimating duration models from flow data when all durations are right censored?\\
e. If the duration distribution is continuous, $c_{i}>b>0$ for some constant $b$, and $\mathrm{P}\left(t_{i}^{*}<t\right)>0$ for all $t>0$, is it likely, in a large random sample, to find that all durations have been censored?\\
22.4. Suppose that, in the context of flow sampling, we observe for each $i$ covariates $\mathbf{x}_{i}$, the censoring time $c_{i}$, and the binary indicator $d_{i}(=1$ if the observation is uncensored). We never observe $t_{i}^{*}$.\\
a. Show that the conditional likelihood function has the binary response form. What is the binary "response"?\\
b. Use the Weibull model to demonstrate the following when we only observe whether durations are censored: if the censoring times $c_{i}$ are constant, the parameters $\boldsymbol{\beta}$ and $\alpha$ are not identified. (Hint: Consider the same case as in Problem 22.3c, and show that the log likelihood depends only on the constant $\exp (\beta) c^{\alpha}$, where $c$ is the common censoring time.)\\
c. Use the lognormal model to argue that, provided the $c_{i}$ vary across $i$ in the population, the parameters are generally identified. (Hint: In the binary response model, what is the coefficient on $\log \left(c_{i}\right)$ ?)\\
22.5. In this problem you are to derive the log likelihood in equation (22.30). Assume that $c_{i}>b-a_{i}$ for all $i$, so that we always observe part of each spell after the sampling date, $b$. In what follows, we supress the parameter vector, $\boldsymbol{\theta}$.\\
a. For $b-a_{i}<t<c_{i}$, show that $\mathrm{P}\left(t_{i}^{*} \leq t \mid \mathbf{x}_{i}, a_{i}, c_{i}, s_{i}=1\right)=\left[F\left(t \mid \mathbf{x}_{i}\right)-F\left(b-a_{i} \mid \mathbf{x}_{i}\right)\right] / \left[1-F\left(b-a_{i} \mid \mathbf{x}_{i}\right)\right]$.\\
b. Use part a to obtain the density of $t_{i}^{*}$ conditional on ( $\mathbf{x}_{i}, a_{i}, c_{i}, s_{i}=1$ ) for $b-a_{i}<t<c_{i}$.\\
c. Show that $\mathrm{P}\left(t_{i}=c_{i} \mid \mathbf{x}_{i}, a_{i}, c_{i}, s_{i}=1\right)=\left[1-F\left(c_{i} \mid \mathbf{x}_{i}\right)\right] /\left[1-F\left(b-a_{i} \mid \mathbf{x}_{i}\right)\right]$.\\
d. Explain why parts b and c lead to equation (22.30).\\
22.6. Consider the problem of stock sampling where we do not follow spells after the sampling date, $b$, as described in Section 22.3.3. Let $F\left(\cdot \mid \mathbf{x}_{i}\right)$ denote the $\operatorname{cdf}$ of $t_{i}^{*}$ given $\mathbf{x}_{i}$, and let $k\left(\cdot \mid \mathbf{x}_{i}\right)$ denote the continuous density of $a_{i}$ given $\mathbf{x}_{i}$. We drop dependence on the parameters for most of the derivations. Assume that $t_{i}^{*}$ and $a_{i}$ are independent conditional on $\mathbf{x}_{i}$.\\
a. Let $s_{i}$ denote a selection indicator, so that $s_{i}=1\left(t_{i}^{*} \geq b-a_{i}\right)$. For any $0<a<b$, show that\\
$\mathrm{P}\left(a_{i} \leq a, s_{i}=1 \mid \mathbf{x}_{i}\right)=\int_{0}^{a} k\left(u \mid \mathbf{x}_{i}\right)\left[1-F\left(b-u \mid \mathbf{x}_{i}\right)\right] \mathrm{d} u$\\
b. Derive equation (22.32). (Hint: $\mathrm{P}\left(s_{i}=1 \mid \mathbf{x}_{i}\right)=\mathrm{E}\left(s_{i} \mid \mathbf{x}_{i}\right)=\mathrm{E}\left[\mathrm{E}\left(s_{i} \mid a_{i}, \mathbf{x}_{i}\right) \mid \mathbf{x}_{i}\right]$, and $\mathrm{E}\left(s_{i} \mid a_{i}, \mathbf{x}_{i}\right)=\mathrm{P}\left(t_{i}^{*} \geq b-a_{i} \mid \mathbf{x}_{i}\right)$.)\\
c. For $0<a<b$, what is the cdf of $a_{i}$ given $\mathbf{x}_{i}$ and $s_{i}=1$ ? Now derive equation (22.31).\\
d. Take $b=1$, and assume that the starting time distribution is uniform on $[0,1]$ (independent of $\mathbf{x}_{i}$ ). Find the density (22.31) in this case.\\
e. For the setup in part d , assume that the duration cdf has the Weibull form, $1-\exp \left[-\exp \left(\mathbf{x}_{i} \boldsymbol{\beta}\right) t^{\alpha}\right]$. What is the $\log$ likelihood for observation $i$ ?\\
22.7. Consider the original stock sampling problem that we covered in Section 22.3.3. There, we derived the log likelihood (22.30) by conditioning on the starting times, $a_{i}$. This approach is convenient because we do not have to specify a distribution for the starting times. But suppose we have an acceptable model for $k\left(\cdot \mid \mathbf{x}_{i} ; \boldsymbol{\eta}\right)$, the (continuous) density of $a_{i}$ given $\mathbf{x}_{i}$. Further, we maintain assumption (22.22) and assume $\mathrm{D}\left(a_{i} \mid c_{i}, \mathbf{x}_{i}\right)=\mathrm{D}\left(a_{i} \mid \mathbf{x}_{i}\right)$.\\
a. Show that the log-likelihood function conditional on $\mathbf{x}_{i}$, which accounts for truncation, is


\begin{gather*}
\sum_{i=1}^{N}\left\{d_{i} \log \left[f\left(t_{i} \mid \mathbf{x}_{i} ; \boldsymbol{\theta}\right)\right]+\left(1-d_{i}\right) \log \left[1-F\left(t_{i} \mid \mathbf{x}_{i} ; \boldsymbol{\theta}\right)\right]\right. \\
\left.+\log \left[k\left(a_{i} \mid \mathbf{x}_{i} ; \boldsymbol{\eta}\right)\right]-\log \left[\mathrm{P}\left(s_{i}=1 \mid \mathbf{x}_{i} ; \boldsymbol{\theta}, \boldsymbol{\eta}\right)\right]\right\} \tag{22.56}
\end{gather*}


where $\mathrm{P}\left(s_{i}=1 \mid \mathbf{x}_{i} ; \boldsymbol{\theta}, \boldsymbol{\eta}\right)$ is given in equation (22.32).\\
b. Discuss the trade-offs in using equation (22.30) or the log likelihood in (22.56).\\
22.8. In the context of stock sampling, where we are interested in the population of durations starting in $[0, b]$, suppose that we interview at date $b$, as usual, but we do not observe any starting times. (This assumption raises the issue of how we know individual $i$ 's starting time is in the specified interval, $[0, b]$. We assume that the interval is defined to make this condition true for all $i$.) Let $r_{i}^{*}=a_{i}+t_{i}^{*}-b$, which can be interpreted as the calendar date at which the spell ends minus the interview date. Even without right censoring, we observe $r_{i}^{*}$ only if $r_{i}^{*}>0$, in which case $r_{i}^{*}$ is simply the time in the spell since the interview date, $b$. Assume that $t_{i}^{*}$ and $a_{i}$ are independent conditional on $\mathbf{x}_{i}$.\\
a. Show that for $r>0$, the density of $r_{i}^{*}$ given $\mathbf{x}_{i}$ is\\
$h\left(r \mid \mathbf{x}_{i} ; \boldsymbol{\theta}, \boldsymbol{\eta}\right) \equiv \int_{0}^{b} k\left(u \mid \mathbf{x}_{i} ; \boldsymbol{\eta}\right) f\left(r+b-u \mid \mathbf{x}_{i} ; \boldsymbol{\theta}\right) d u$\\
where, as before, $k\left(a \mid \mathbf{x}_{i} ; \boldsymbol{\eta}\right)$ is the density of $a_{i}$ given $\mathbf{x}_{i}$ and $f\left(t \mid \mathbf{x}_{i} ; \boldsymbol{\theta}\right)$ is the duration density.\\
b. Let $q>0$ be a fixed censoring time after the interview date, and define $r_{i}= \min \left(r_{i}^{*}, q\right)$. Find $\mathrm{P}\left(r_{i}=q \mid \mathbf{x}_{i}\right)$ in terms of the cdf of $r_{i}^{*}$, say, $H\left(r \mid \mathbf{x}_{i} ; \boldsymbol{\theta}, \boldsymbol{\eta}\right)$.\\
c. Use parts $a$ and $b$, along with equation (22.32), to show that the log likelihood conditional on observing $\left(r_{i}, \mathbf{x}_{i}\right)$ is


\begin{align*}
& d_{i} \log \left[h\left(r_{i} \mid \mathbf{x}_{i} ; \boldsymbol{\theta}, \boldsymbol{\eta}\right)\right]+\left(1-d_{i}\right) \log \left[1-H\left(r_{i} \mid \mathbf{x}_{i} ; \boldsymbol{\theta}, \boldsymbol{\eta}\right)\right] \\
& -\log \left\{\int_{0}^{b}\left[1-F\left(b-u \mid \mathbf{x}_{i} ; \boldsymbol{\theta}\right)\right] k\left(u \mid \mathbf{x}_{i} ; \boldsymbol{\eta}\right) \mathrm{d} u\right\} \tag{22.57}
\end{align*}


where $d_{i}=1$ if observation $i$ has not been right censored.\\
d. Simplify the log likelihood from part c when $b=1$ and $k\left(a \mid \mathbf{x}_{i} ; \boldsymbol{\eta}\right)$ is the uniform density on $[0,1]$.\\
22.9. Consider the Weibull model with multiplicative heterogeneity, as in equation (22.33), but where $v_{i}$ takes on only two values. Think of there being two types of\\
people, A and B . Let $0<\eta<1$ be the value for type A people with $\rho=P\left(v_{i}=\eta\right)$, $0<\rho<1$.\\
a. Show that to ensure $E\left(v_{i}\right)=1$, the value of $v_{i}$ for type B people must be $(1-\rho \eta) /(1-\rho)$.\\
b. Find the cdf of $t_{i}^{*}$ conditional on $\mathbf{x}_{i}$ only. Call this $G\left(t \mid \mathbf{x}_{i} ; \alpha, \boldsymbol{\beta}, \eta, \rho\right)$.\\
c. Find the density of $t_{i}^{*}$ conditional on $\mathbf{x}_{i}, g\left(t \mid \mathbf{x}_{i} ; \alpha, \boldsymbol{\beta}, \eta, \rho\right)$. How would you estimate the parameters if you have right-censored data?\\
22.10. Let $0<a_{1}<a_{2}<\cdots<a_{M-1}<a_{M}$ be a positive, increasing set of constants, and let $T$ be a nonnegative random variable with $\mathrm{P}(T>0)=1$.\\
a. Show that, for any $m=1, \ldots, M, \mathrm{P}\left(T>a_{m}\right)=\mathrm{P}\left(T>a_{m} \mid T>a_{m-1}\right) \mathrm{P}\left(T>a_{m-1}\right)$.\\
b. Use part a to derive equation (22.48).\\
22.11. Use the data in RECID to answer the following questions.\\
a. Using the same explanatory variables as in Table 22.2, estimate a model with gamma-distributed heterogeneity and a log-logistic hazard for $\kappa(t ; \mathbf{x})$. Does the model fit better or worse than the Weibull/gamma mixture model reported in Table 22.2?\\
b. Graph the conditional hazard function (with $v=1$ ) and at the mean values of the covariates. How does its shape compare with the Weibull/gamma model?\\
c. Graph the unconditional hazard function and comment on its shape. How do the conditional and unconditional hazards differ?

\section*{References}
Abadie, A. (2003), "Semiparametric Instrumental Variable Estimation of Treatment Response Models," Journal of Econometrics 113, 231-263.\\
Abadie, A., and G. W. Imbens (2006), "Large Sample Properties of Matching Estimators for Average Treatment Effects," Econometrica 74, 235-267.\\
Abbring, J. H., and J. J. Heckman (2007), "Econometric Evaluation of Social Programs, Part III: Distributional Treatment Effects, Dynamic Treatment Effects and Dynamic Discrete Choice, and General Equilibrium Policy Evaluation," in Handbook of Econometrics, Volume 6B, ed. J. J. Heckman and E. Leamer. Amsterdam: North Holland, 5145-5303.

Abrevaya, J. (1997), "The Equivalence of Two Estimators for the Fixed Effects Logit Model," Economics Letters 55, 41-43.\\
Abrevaya, J., and C. Dahl (2008), "The Effects of Birth Inputs on Birthweight: Evidence from Quantile Estimation on Panel Data," Journal of Business and Economic Statistics 26, 379-397.\\
Adams, J. D., E. P. Chiang, and J. L. Jensen (2003), "The Influence of Federal Laboratory R\&D on Industrial Research," Review of Economics and Statistics 85, 1003-1020.\\
Ahn, S. C., Y. H. Lee, and P. Schmidt (2002), "GMM Estimation of Linear Panel Data Models with Time-Varying Individual Effects," Journal of Econometrics 101, 219-255.\\
Ahn, H., and J. L. Powell (1993), "Semiparametric Estimation of Censored Selection Models with a Nonparametric Selection Mechanism," Journal of Econometrics 58, 3-29.\\
Ahn, S. C., and P. Schmidt (1995), "Efficient Estimation of Models for Dynamic Panel Data," Journal of Econometrics 68, 5-27.\\
Ai, C. (1997), "A Semiparametric Maximum Likelihood Estimator," Econometrica 65, 933-963.\\
Aitchison, J., and S. D. Silvey (1958), "Maximum-Likelihood Estimation of Parameters Subject to Constraints," Annals of Mathematical Statistics 29, 813-828.\\
Altonji, J. G., T. E. Elder, and C. R. Tabor (2005), "An Evaluation of Instrumental Variable Strategies for Estimating the Effects of Catholic Schooling," Journal of Human Resources 40, 791-821.\\
Altonji, J., and R. L. Matzkin (2005), "Cross Section and Panel Data Estimators for Nonseparable Models with Endogenous Regressors," Econometrica 73, 1053-1102.\\
Altonji, J. G., and L. M. Segal (1996), "Small-Sample Bias in GMM Estimation of Covariance Structures," Journal of Business and Economic Statistics 14, 353-366.\\
Alvarez, J., and M. Arellano (2003), "Time Series and Cross-Section Asymptotics of Dynamic Panel Data Estimators," Econometrica 71, 121-159.\\
Amemiya, T. (1973), "Regression Analysis When the Dependent Variable Is Truncated Normal," Econometrica 41, 997-1016.\\
Amemiya, T. (1974), "The Nonlinear Two-Stage Least-Squares Estimator," Journal of Econometrics 2, 105-110.

Amemiya, T. (1985), Advanced Econometrics. Cambridge, MA: Harvard University Press.\\
Andersen, E. B. (1970), "Asymptotic Properties of Conditional Maximum Likelihood Estimators," Journal of the Royal Statistical Society, Series B, 32, 283-301.\\
Anderson, T. W., and C. Hsiao (1982), "Formulation and Estimation of Dynamic Models Using Panel Data," Journal of Econometrics 18, 67-82.\\
Andrews, D. W. K. (1989), "Power in Econometric Applications," Econometrica 57, 1059-1090.\\
Angrist, J. D. (1990), "Lifetime Earnings and the Vietnam Era Draft Lottery: Evidence from Social Security Administrative Records," American Economic Review 80, 313-336.\\
Angrist, J. D. (1991), "Instrumental Variables Estimation of Average Treatment Effects in Econometrics and Epidemiology," National Bureau of Economic Research Technical Working Paper Number 115.\\
Angrist, J. D. (1998), "Estimating the Labor Market Impact of Voluntary Military Service Using Social Security Data on Military Applicants," Econometrica 66, 249-288.

Angrist, J., V. Chernozhukov, and I. Fernandez-Val (2006), "Quantile Regression under Misspecification, with an Application to the U.S. Wage Structure," Econometrica 74, 539-563.\\
Angrist, J. D., and W. N. Evans (1998), "Children and Their Parents' Labor Supply: Evidence from Exogenous Variation in Family Size," American Economic Review 88, 450-477.\\
Angrist, J. D., G. W. Imbens, and D. B. Rubin (1996), "Identification and Causal Effects Using Instrumental Variables," Journal of the American Statistical Association 91, 444-455.\\
Angrist, J. D., and A. B. Krueger (1991), "Does Compulsory School Attendance Affect Schooling and Earnings?" Quarterly Journal of Economics 106, 979-1014.\\
Angrist, J. D., and V. Lavy (2002), "The Effect of High School Matriculation Awards: Evidence from Randomized Trials," National Bureau of Economic Research Working Paper 9389.\\
Angrist, J. D., and W. K. Newey (1991), "Overidentification Tests in Earnings Functions with Fixed Effects," Journal of Business and Economic Statistics 9, 317-323.\\
Angrist, J. D., and J.-S. Pischke (2009), Mostly Harmless Econometrics. Princeton, NJ: Princeton University Press.\\
Arellano, M. (1987), "Computing Robust Standard Errors for Within-Groups Estimators," Oxford Bulletin of Economics and Statistics 49, 431-434.\\
Arellano, M. (2003), Panel Data Econometrics. Oxford: Oxford University Press.\\
Arellano, M., and S. R. Bond (1991), "Some Specification Tests for Panel Data: Monte Carlo Evidence and an Application to Employment Equations," Review of Economic Studies 58, 277-298.\\
Arellano, M., and O. Bover (1995), "Another Look at the Instrumental Variables Estimation of ErrorComponent Models," Journal of Econometrics 68, 29-51.\\
Arellano, M., and B. E. Honoré (2001), "Panel Data: Some Recent Developments," in Handbook of Econometrics, Volume 5, ed. E. Leamer and J. J. Heckman. Amsterdam: North Holland, 3229-3296.\\
Ayers, I., and S. D. Levitt (1998), "Measuring Positive Externalities from Unobservable Victim Precaution: An Empirical Analysis of Lojack," Quarterly Journal of Economics 108, 43-77.\\
Baltagi, B. H. (1981), "Simultaneous Equations with Error Components," Journal of Econometrics 17, 189-200.\\
Baltagi, B. H. (2001), Econometric Analysis of Panel Data, second edition. West Sussex, UK: Wiley.\\
Baltagi, B. H., and Q. Li (1991), "A Joint Test for Serial Correlation and Random Individual Effects," Statistics and Probability Letters 11, 277-280.\\
Baltagi, B. H., and Q. Li (1995), "Testing AR(1) Against MA(1) Disturbances in an Error Component Model," Journal of Econometrics 68, 133-151.\\
Barnow, B., G. Cain, and A. Goldberger (1980), "Issues in the Analysis of Selectivity Bias," Evaluation Studies 5, 42-59.\\
Bartik, T. J. (1987), "The Estimation of Demand Parameters in Hedonic Price Models," Journal of Political Economy 95, 81-88.\\
Bartle, R. G. (1966), The Elements of Integration. New York: Wiley.\\
Bassett, G., and R. Koenker (1978), "Asymptotic Theory of Least Absolute Error Regression," Journal of the American Statistical Association 73, 618-622.\\
Bassi, L. J. (1984), "Estimating the Effect of Job Training Programs with Non-Random Selection," Review of Economics and Statistics 66, 36-43.\\
Bates, C. E., and H. White (1993), "Determination of Estimators with Minimum Asymptotic Covariances Matrices," Econometric Theory 9, 633-648.\\
Baum, C. F., M. E. Schaffer, and S. Stillman (2003), "Instrumental Variables and GMM: Estimation and Testing," The Stata Journal 3, 1-31.

Bekker, P. A. (1994), "Alternative Approximations to the Distributions of Instrumental Variable Estimators," Econometrica 62, 657-681.\\
Bera, A. K., and C. R. McKenzie (1986), "Alternative Forms and Properties of the Score Test," Journal of Applied Statistics 13, 13-25.\\
Berndt, E. R., B. H. Hall, R. E. Hall, and J. A. Hausman (1974), "Estimation and Inference in Nonlinear Structural Models," Annals of Economic and Social Measurement 3, 653-666.\\
Bhargava, A., L. Franzini, and W. Narendranathan (1982), "Serial Correlation and the Fixed Effects Model," Review of Economic Studies 49, 533-549.\\
Bhattacharya, D. (2005), "Asymptotic Inference from Multi-stage Samples," Journal of Econometrics 126, 145-171.\\
Biddle, J. E., and D. S. Hamermesh (1990), "Sleep and the Allocation of Time," Journal of Political Economy 98, 922-943.\\
Billingsley, P. (1979), Probability and Measure. New York: Wiley.\\
Blackburn, M. (2007), "Estimating Wage Differentials without Logarithms," Labour Economics 14, 73-98.\\
Blackburn, M., and D. Neumark (1992), "Unobserved Ability, Efficiency Wages, and Interindustry Wage Differentials," Quarterly Journal of Economics 107, 1421-1436.\\
Blank, R. M. (1988), "Simultaneously Modeling the Supply of Weeks and Hours among Female Household Heads," Journal of Labor Economics 6, 177-204.\\
Blundell, R., and S. Bond (1998), "Initial Conditions and Moment Restrictions in Dynamic Panel Data Models," Journal of Econometrics 87, 115-144.\\
Blundell, R., R. Griffith, and F. Windmeijer (1998), "Individual Effects and Dynamics in Count Data Models," mimeo, Institute of Fiscal Studies, London.\\
Blundell, R., and J. L. Powell (2003), "Endogeneity in Nonparametric and Semiparametric Regression Models," in Advances in Economics and Econometrics, ed. M. Dewatripont, L. P. Hansen, and S. J. Turnovsky. Cambridge: Cambridge University Press, 312-357.\\
Blundell, R., and J. L. Powell (2004), "Endogeneity in Semiparametric Binary Response Models," Review of Economic Studies 71, 655-679.\\
Bound, J., D. A. Jaeger, and R. M. Baker (1995), "Problems with Instrumental Variables Estimation When the Correlation between the Instruments and Endogenous Explanatory Variables Is Weak," Journal of the American Statistical Association 90, 443-450.\\
Breusch, T. S., G. E. Mizon, and P. Schmidt (1989), "Efficient Estimation Using Panel Data," Econometrica 57, 695-700.\\
Breusch, T. S., and A. R. Pagan (1979), "A Simple Test for Heteroskedasticity and Random Coefficient Variation," Econometrica 50, 987-1007.\\
Breusch, T. S., and A. R. Pagan (1980), "The LM Test and Its Applications to Model Specification in Econometrics," Review of Economic Studies 47, 239-254.\\
Breusch, T., H. Qian, P. Schmidt, and D. Wyhowski (1999), "Redundancy of Moment Conditions," Journal of Econometrics 91, 89-111.\\
Brown, B. W., and M. B. Walker (1995), "Stochastic Specification in Random Production Models of CostMinimizing Firms," Journal of Econometrics 66, 175-205.\\
Buchinsky, M. (1994), "Changes in the U.S. Wage Structure, 1963-1987: Application of Quantile Regression," Econometrica 62, 405-458.\\
Buchinsky, M. (1998), "Recent Advances in Quantile Regression Models: A Practical Guide for Empirical Research," Journal of Human Resources 33, 88-126.\\
Buchinsky, M., and J. Hahn (1998), "An Alternative Estimator for the Censored Quantile Regression Model," Econometrica 66, 653-671.

Burnett, N. (1997), "Gender Economics Courses in Liberal Arts Colleges," Journal of Economic Education 28, 369-377.\\
Butler, J. S., and R. A. Moffitt (1982), "A Computationally Efficient Quadrature Procedure for the OneFactor Multinomial Probit Model," Econometrica 50, 761-764.\\
Cameron, A. C., and P. K. Trivedi (1986), "Econometric Models Based on Count Data: Comparisons and Applications of Some Estimators and Tests," Journal of Applied Econometrics 1, 29-53.\\
Cameron, A. C., and P. K. Trivedi (2005), Microeconometrics: Methods and Applications. New York: Cambridge University Press.\\
Card, D. (1995), "Using Geographic Variation in College Proximity to Estimate the Return to Schooling," in Aspects of Labour Market Behavior: Essays in Honour of John Vanderkamp, ed. L. N. Christophides, E. K. Grant, and R. Swidinsky. Toronto: University of Toronto Press, 201-222.

Card, D. (2001), "Estimating the Return to Schooling: Progress on Some Persistent Econometric Problems," Econometrica 69, 1127-1160.\\
Card, D., and A. B. Krueger (1994), "Minimum Wages and Employment: A Case Study of the Fast-Food Industry in New Jersey and Pennsylvania," American Economic Review 84, 772-793.\\
Cardell, N. S., and M. M. Hopkins (1977), "Education, Income, and Ability: A Comment," Journal of Political Economy 85, 211-215.\\
Carpenter, C. (2004), "How Do Zero Tolerance Drunk Driving Laws Work?" Journal of Health Economics 23, 61-83.\\
Casella, G., and R. L. Berger (2002), Statistical Inference, second edition. Pacific Grove, CA: Duxbury.\\
Chamberlain, G. (1980), "Analysis of Covariance with Qualitative Data," Review of Economic Studies 47, 225-238.\\
Chamberlain, G. (1982), "Multivariate Regression Models for Panel Data," Journal of Econometrics 18, 5-46.\\
Chamberlain, G. (1984), "Panel Data," in Handbook of Econometrics, Volume 2, ed. Z. Griliches and M. D. Intriligator. Amsterdam: North Holland, 1247-1318.

Chamberlain, G. (1985), "Heterogeneity, Omitted Variable Bias, and Duration Dependence," in Longitudinal Analysis of Labor Market Data, ed. J. J. Heckman and B. Singer. Cambridge: Cambridge University Press, 3-38.\\
Chamberlain, G. (1987), "Asymptotic Efficiency in Estimation with Conditional Moment Restrictions," Journal of Econometrics 34, 305-334.\\
Chamberlain, G. (1992a), "Efficiency Bounds for Semiparametric Regression," Econometrica 60, 567-596.\\
Chamberlain, G. (1992b), "Comment: Sequential Moment Restrictions in Panel Data," Journal of Business and Economic Statistics 10, 20-26.\\
Chamberlain, G. (1994), "Quantile Regression, Censoring and the Structure of Wages," in Proceedings of the Sixth World Congress of the Econometric Society, ed. C. Sims and J. J. Laffont. New York: Cambridge University Press, 171-209.\\
Chay, K. Y., and D. R. Hyslop (2001), "Identification and Estimation of Dynamic Binary Response Models: Empirical Evidence Using Alternative Approaches," University of California, Berkeley, Department of Economics Working Paper.\\
Chesher, A., and R. Spady (1991), "Asymptotic Expansions of the Information Matrix Test Statistic," Econometrica 59, 787-815.\\
Chung, C.-F., and A. Goldberger (1984), "Proportional Projections in Limited Dependent Variable Models," Econometrica 52, 531-534.\\
Chung, C.-F., P. Schmidt, and A. D. Witte (1991), "Survival Analysis: A Survey," Journal of Quantitative Criminology 7, 59-98.\\
Cornwell, C., and D. Trumball (1994), "Estimating the Economic Model of Crime with Panel Data," Review of Economics and Statistics 76, 360-366.

Cosslett, S. R. (1981), "Efficient Estimation of Discrete-Choice Models," in Structural Analysis of Discrete Data with Econometric Applications, ed. C. F. Manski and D. McFadden. Cambridge, MA: MIT Press, 51-111.\\
Cosslett, S. R. (1993), "Estimation from Endogenously Stratified Samples," in Handbook of Statistics, Volume 11, ed. G. S. Maddala, C. R. Rao, and H. D. Vinod. Amsterdam: North Holland, 1-43.\\
Cox, D. R. (1961), "Tests of Separate Families of Hypotheses," in Proceedings of the Fourth Berkeley Symposium on Mathematical Statistics and Probability. Berkeley: University of California Press, 105-123.\\
Cox, D. R. (1962), "Further Results on Tests of Separate Families of Hypotheses," Journal of the Royal Statistical Society, Series B, 24, 406-424.\\
Cox, D. R. (1972), "Regression Models and Life Tables," Journal of the Royal Statistical Society, Series B, 34, 187-220.\\
Cragg, J. (1971), "Some Statistical Models for Limited Dependent Variables with Applications to the Demand for Durable Goods," Econometrica 39, 829-844.\\
Cragg, J. (1983), "More Efficient Estimation in the Presence of Heteroskedasticity of Unknown Form," Econometrica 51, 751-763.\\
Cragg, J. G., and S. G. Donald (1996), "Inferring the Rank of a Matrix," Journal of Econometrics 76, 223-250.\\
Currie, J., and D. Thomas (1995), "Does Head Start Make a Difference?" American Economic Review 85, 341-364.\\
Cutler, D. M., and E. L. Glaeser (1997), "Are Ghettos Good or Bad?" Quarterly Journal of Economics 112, 827-872.\\
Davidson, J. (1994), Stochastic Limit Theory. Oxford: Oxford University Press.\\
Davidson, R., and J. G. MacKinnon (1984), "Convenient Specification Tests for Logit and Probit Models," Journal of Econometrics 24, 241-262.\\
Davidson, R., and J. G. MacKinnon (1985), "Heteroskedasticity-Robust Tests in Regression Directions," Annale de l'INSÉÉ 59/60, 183-218.\\
Davidson, R., and J. G. MacKinnon (1992), "A New Form of the Information Matrix Test," Econometrica 60, 145-147.\\
Davidson, R., and J. G. MacKinnon (1993), Estimation and Inference in Econometrics. New York: Oxford University Press.\\
Deaton, A. (1995), "Data and Econometric Tools for Development Analysis," in Handbook of Development Economics, Volume 3A, ed. J. Berhman and T. N. Srinivasan. Amsterdam: North Holland, 1785-1882.\\
Dehejia, R. H., and S. Wahba (1999), "Causal Effects in Non-Experimental Studies: Evaluating the Evaluation of Training Programs," Journal of the American Statistical Association 94, 1053-1062.\\
Donald, S. G., and K. Lang (2007), "Inference with Difference-in-Differences and Other Panel Data," Review of Economics and Statistics 89, 221-233.\\
Donald, S. G., and H. J. Paarsch (1996), "Identification, Estimation, and Testing in Parametric Empirical Models of Auctions within the Independent Private Values Paradigm," Econometric Theory 12, 517-567.\\
Downes, T. M., and S. M. Greenstein (1996), "Understanding the Supply Decisions of Nonprofits: Modeling the Location of Private Schools," Rand Journal of Economics 27, 365-390.\\
Duan, N. (1983), "Smearing Estimate: A Nonparametric Retransformation Method," Journal of the American Statistical Association 78, 605-610.\\
Duan, N., W. G. Manning, C. N. Morris, and J. P. Newhouse (1984), "Choosing between the SampleSelection Model and the Multi-Part Model," Journal of Business and Economic Statistics 2, 283-289.\\
Durbin, J. (1954), "Errors in Variables," Review of the International Statistical Institute 22, 23-32.

Dustmann, C., and M. E. Rochina-Barrachina (2007), "Selection Correction in Panel Data Models: An Application to the Estimation of Females' Wage Equations," Econometrics Journal 10, 263-293.\\
Eicker, F. (1967), "Limit Theorems for Regressions with Unequal and Dependent Errors," Proceedings of the Fifth Berkeley Symposium on Mathematical Statistics and Probability 1, 59-82. Berkeley: University of California Press.\\
Elbers, C., and G. Ridder (1982), "True and Spurious Duration Dependence: The Identifiability of the Proportional Hazard Model," Review of Economic Studies 49, 403-410.\\
El Sayyad, G. M. (1973), "Bayesian and Classical Analysis of Poisson Regression," Journal of the Royal Statistical Society, Series B, 35, 445-451.\\
Engle, R. F. (1982), "Autoregressive Conditional Heteroskedasticity with Estimates of the Variance of U.K. Inflation," Econometrica 50, 987-1008.

Engle, R. F. (1984), "Wald, Likelihood Ratio, and Lagrange Multiplier Statistics in Econometrics," in Handbook of Econometrics, Volume 2, ed. Z. Griliches and M. D. Intriligator. Amsterdam: North Holland, 776-828.\\
Epple, D. (1987), "Hedonic Prices and Implicit Markets: Estimated Demand and Supply Functions for Differentiated Products," Journal of Political Economy 95, 59-80.\\
Estrella, A. (1998), "A New Measure of Fit for Equations with Dichotomous Dependent Variables," Journal of Business and Economic Statistics 16, 198-205.\\
Evans, W. N., W. E. Oates, and R. M. Schwab (1992), "Measuring Peer Group Effects: A Study of Teenage Behavior," Journal of Political Economy 100, 966-991.\\
Evans, W. N., and R. M. Schwab (1995), "Finishing High School and Starting College: Do Catholic Schools Make a Difference?" Quarterly Journal of Economics 110, 941-974.\\
Fernández-Val, I. (2009), "Fixed Effects Estimation of Structural Parameters and Marginal Effects in Panel Probit Models," Journal of Econometrics 150, 71-85.\\
Fin, T., and P. Schmidt (1984), "A Test of the Tobit Specification Against an Alternative Suggested by Cragg," Review of Economics and Statistics 66, 174-177.\\
Fisher, F. M. (1965), "Identifiability Criteria in Nonlinear Systems: A Further Note," Econometrica 33, 197-205.\\
Flores-Lagunes, A. (2007), "Finite Sample Evidence of IV Estimators under Weak Instruments," Journal of Applied Econometrics 22, 677-694.\\
Foster, A. D., and M. R. Rosenzweig (1995), "Learning by Doing and Learning from Others: Human Capital and Technical Change in Agriculture," Journal of Political Economy 103, 1176-1209.\\
Franses, P. H., and R. Paap (2001), Quantitative Models in Marketing Research. Cambridge: Cambridge University Press.\\
Friedberg, L. (1998), "Did Unilateral Divorce Raise Divorce Rates? Evidence from Panel Data," American Economic Review 88, 608-627.\\
Frölich, M. (2004), "Finite-Sample Properties of Propensity-Score Matching and Weighting Estimators," Review of Economics and Statistics 86, 77-90.\\
Gallant, A. R. (1987), Nonlinear Statistical Models. New York: Wiley.\\
Gallant, A. R., and H. White (1988), A Unified Theory of Estimation and Inference for Nonlinear Dynamic Models. New York: Blackwell.\\
Garen, J. (1984), "The Returns to Schooling: A Selectivity Bias Approach with a Continuous Choice Variable," Econometrica 52, 1199-1218.\\
Geronimus, A. T., and S. Korenman (1992), "The Socioeconomic Consequences of Teen Childbearing Reconsidered," Quarterly Journal of Economics 107, 1187-1214.\\
Geweke, J., and M. P. Keane (2001), "Computationally Intensive Methods for Integration in Economics," in Handbook of Econometrics, Volume 5, ed. E. Leamer and J. J. Heckman. Amsterdam: North Holland, 3463-3568.

Gill, R. D., and J. M. Robins (2001), "Causal Inference for Complex Longitudinal Data: The Continuous Case," Annals of Statistics 29, 1785-1811.\\
Goldberger, A. S. (1968), Topics in Regression Analysis. New York: Macmillan.\\
Goldberger, A. S. (1972), "Structural Equation Methods in the Social Sciences," Econometrica 40, 979-1001.\\
Goldberger, A. S. (1981), "Linear Regression after Selection," Journal of Econometrics 15, 357-366.\\
Goldberger, A. S. (1991), A Course in Econometrics. Cambridge, MA: Harvard University Press.\\
Gordy, M. B. (1999), "Hedging Winner's Curse with Multiple Bids: Evidence from the Portuguese Treasury Bill Auction," Review of Economics and Statistics 81, 448-465.\\
Gourieroux, G., and A. Monfort (1994), "Testing Non-nested Hypotheses," in Handbook of Econometrics, Volume 4, ed. R. F. Engle and D. L. McFadden. Amsterdam: North Holland, 2583-2637.\\
Gourieroux, C., A. Monfort, and C. Trognon (1984a), "Pseudo-Maximum Likelihood Methods: Theory," Econometrica 52, 681-700.\\
Gourieroux, C., A. Monfort, and C. Trognon (1984b), "Pseudo-Maximum Likelihood Methods: Applications to Poisson Models," Econometrica 52, 701-720.\\
Greene, W. (1997), Econometric Analysis, third edition. New York: Macmillan.\\
Greene, W. (1998), "Gender Economics Courses in Liberal Arts Colleges: Further Results," Journal of Economic Education 29, 291-300.\\
Greene, W. (2003), Econometric Analysis, fifth edition. Upper Saddle River, NJ: Prentice Hall.\\
Greene, W. (2004), "The Behaviour of the Maximum Likelihood Estimator of Limited Dependent Variable Models in the Presence of Fixed Effects," Econometrics Journal 7, 98-119.\\
Gregory, A. W., and M. R. Veall (1985), "On Formulating Wald Tests for Nonlinear Restrictions," Econometrica 53, 1465-1468.\\
Griliches, Z., B. H. Hall, and J. A. Hausman (1978), "Missing Data and Self-Selection in Large Panels," Annale de l'INSÉÉ 30/31, 137-176.\\
Griliches, Z., and J. A. Hausman (1986), "Errors in Variables in Panel Data," Journal of Econometrics 31, 93-118.\\
Griliches, Z., and W. M. Mason (1972), "Education, Income and Ability," Journal of Political Economy, Part II, 80, S74-S103.\\
Gronau, R. (1974), "Wage Comparisons—A Selectivity Bias," Journal of Political Economy 82, 1119-1143.\\
Gruber, J. (1994), "The Incidence of Mandated Maternity Benefits," American Economic Review 84, 622-641.\\
Gruber, J., and J. M. Poterba (1994), "Tax Incentives and the Decision to Purchase Health Insurance: Evidence from the Self-Employed," Quarterly Journal of Economics 109, 701-733.\\
Haavelmo, T. (1943), "The Statistical Implications of a System of Simultaneous Equations," Econometrica 11, 1-12.\\
Hagy, A. P. (1998), "The Demand for Child Care Quality: An Hedonic Price Approach," Journal of Human Resources 33, 683-710.\\
Hahn, J. (1998), "On the Role of the Propensity Score in Efficient Semiparametric Estimation of Average Treatment Effects," Econometrica 66, 315-331.\\
Hahn, J. (1999), "How Informative Is the Initial Condition in the Dynamic Panel Data Model with Fixed Effects?" Journal of Econometrics 93, 309-326.\\
Hahn, J., and W. K. Newey (2004), "Jackknife and Analytical Bias Reduction for Nonlinear Panel Models," Econometrica 72, 1295-1319.\\
Hahn, J., P. Tood, and W. van der Klaauw (2001), "Identification and Estimation of Treatment Effects with a Regression-Discontinuity Design," Econometrica 69, 201-209.

Hajivassiliou, V. A., and P. A. Ruud (1994), "Classical Estimation Methods for LDV Models Using Simulation," in Handbook of Econometrics, Volume 4, ed. R. F. Engle and D. McFadden. Amsterdam: North Holland, 2383-2441.\\
Hall, A. (1987), "The Information Matrix Test for the Linear Model," Review of Economic Studies 54, 257-263.\\
Ham, J. C., and R. J. Lalonde (1996), "The Effect of Sample Selection and Initial Conditions in Duration Models: Evidence from Experimental Data on Training," Econometrica 64, 175-205.\\
Hamilton, J. D. (1994), Time Series Analysis. Princeton, NJ: Princeton University Press.\\
Han, A. K., and J. A. Hausman (1990), "Flexible Parametric Estimation of Duration and Competing Risk Models," Journal of Applied Econometrics 5, 1-28.\\
Han, C., and P. C. B. Phillips (2006), "GMM with Many Moment Conditions," Econometrica 74, 147-192.\\
Hansen, C. B. (2007), "Asymptotic Properties of a Robust Variance Matrix Estimator for Panel Data When T Is Large," Journal of Econometrics 141, 597-620.\\
Hansen, L. P. (1982), "Large Sample Properties of Generalized Method of Moments Estimators," Econometrica 50, 1029-1054.\\
Härdle, W., and O. Linton (1994), "Applied Nonparametric Methods," in Handbook of Econometrics, Volume 4, ed. R. F. Engle and D. McFadden. Amsterdam: North Holland, 2295-2339.\\
Harris, R. D. F., and E. Tzavalis (1999), "Inference for Unit Roots in Dynamic Panels Where the Time Dimension Is Fixed," Journal of Econometrics 91, 201-226.\\
Hausman, J. A. (1978), "Specification Tests in Econometrics," Econometrica 46, 1251-1271.\\
Hausman, J. A. (1983), "Specification and Estimation of Simultaneous Equations Models," in Handbook of Econometrics, Volume 1, ed. Z. Griliches and M. D. Intriligator. Amsterdam: North Holland, 391-448.\\
Hausman, J. A., B. H. Hall, and Z. Griliches (1984), "Econometric Models for Count Data with an Application to the Patents-R\&D Relationship," Econometrica 52, 909-938.\\
Hausman, J. A., and D. L. McFadden (1984), "A Specification Test for the Multinomial Logit Model," Econometrica 52, 1219-1240.\\
Hausman, J. A., W. K. Newey, and W. E. Taylor (1987), "Efficient Estimation and Identification of Simultaneous Equation Models with Covariance Restrictions," Econometrica 55, 849-874.\\
Hausman, J. A., and W. E. Taylor (1981), "Panel Data and Unobservable Individual Effects," Econometrica 49, 1377-1398.\\
Hausman, J. A., and D. A. Wise (1977), "Social Experimentation, Truncated Distributions, and Efficient Estimation," Econometrica 45, 319-339.\\
Hausman, J. A., and D. A., Wise (1978), "A Conditional Probit Model for Qualitative Choice: Discrete Decisions Recognizing Interdependence and Heterogeneous Preferences," Econometrica 46, 403-426.\\
Hausman, J. A., and D. A. Wise (1981), "Stratification on an Endogenous Variable and Estimation: The Gary Income Maintenance Experiment," in Structural Analysis of Discrete Data with Econometric Applications, ed. C. F. Manski and D. McFadden. Cambridge, MA: MIT Press, 365-391.\\
Heckman, J. J. (1976), "The Common Structure of Statistical Models of Truncation, Sample Selection, and Limited Dependent Variables and a Simple Estimator for Such Models," Annals of Economic and Social Measurement 5, 475-492.\\
Heckman, J. J. (1978), "Dummy Endogenous Variables in a Simultaneous Equations System," Econometrica 46, 931-960.\\
Heckman, J. J. (1979), "Sample Selection Bias as a Specification Error," Econometrica 47, 153-161.\\
Heckman, J. J. (1981), "The Incidental Parameters Problem and the Problem of Initial Conditions in Estimating a Discrete Time-Discrete Data Stochastic Process," in Structural Analysis of Discrete Data with Econometric Applications, ed. C. F. Manski and D. McFadden. Cambridge, MA: MIT Press, 179-195.

Heckman, J. J. (1992), "Randomization and Social Program Evaluation," in Evaluating Welfare and Training Programs, ed. C. F. Manski and I. Garfinkel. Cambridge, MA: Harvard University Press, 201-230.\\
Heckman, J. J. (1997), "Instrumental Variables: A Study of Implicit Behavioral Assumptions Used in Making Program Evaluations," Journal of Human Resources 32, 441-462.\\
Heckman, J. J., and V. J. Hotz (1989), "Choosing among Alternative Nonexperimental Methods for Estimating the Impact of Social Programs: The Case of Manpower Training," Journal of the American Statistical Association 84, 862-875.\\
Heckman, J. J., H. Ichimura, and P. Todd (1998), "Matching as an Econometric Evaluation Estimator," Review of Economic Studies 65, 261-294.\\
Heckman, J. J., R. Lalonde, and J. Smith (2000), "The Economics and Econometrics of Active Labor Markets Programs," in Handbook of Labor Economics, Volume 3, ed. A. Ashenfelter and D. Card. New York: Elsevier Science.\\
Heckman, J. J., L. Lochner, and C. Taber (1998), "General-Equilibrium Treatment Effects," American Economic Review 88, 381-386.\\
Heckman, J. J., and R. Robb (1985), "Alternative Methods for Evaluating the Impact of Interventions," in Longitudinal Analysis of Labor Market Data, ed. J. J. Heckman and B. Singer. New York: Cambridge University Press, 156-245.\\
Heckman, J. J., and B. Singer (1984), "A Method for Minimizing the Impact of Distributional Assumptions in Econometric Models for Duration Data," Econometrica 52, 271-320.\\
Heckman, J. J., S. Urzua, and E. Vytlacil (2006), "Understanding Instrumental Variables in Models with Essential Heterogeneity," Review of Economics and Statistics 88, 389-432.\\
Heckman, J. J., and E. Vytlacil (1998), "Instrumental Variables Methods for the Correlated Random Coefficient Model," Journal of Human Resources 33, 974-987.\\
Heckman, J. J., and E. Vytlacil (2006), "Structural Equations, Treatment Effects, and Econometric Policy Evaluation," Econometrica 73, 669-738.\\
Heckman, J. J., and E. Vytlacil (2007a), "Econometric Evaluation of Social Programs, Part I: Causal Models, Structural Models and Econometric Policy Evaluation," in Handbook of Econometrics, Volume 6B, ed. J. Heckman and E. Leamer. New York: Elsevier Science, Chapter 70, 4779-4874.\\
Heckman, J. J., and E. Vytlacil (2007b), "Econometric Evaluation of Social Programs, Part II: Using the Marginal Treatment Effect to Organize Alternative Econometric Estimators to Evaluate Social Programs, and to Forecast Their Effects in New Environments," in Handbook of Econometrics, Volume 6B, ed. J. Heckman and E. Leamer. New York: Elsevier Science, Chapter 71, 4875-5143.\\
Hendry, D. F. (1984), "Monte Carlo Experimentation in Econometrics," in Handbook of Econometrics, Volume 2, ed. Z. Griliches and M. D. Intriligator. Amsterdam: North Holland, 937-976.\\
Hirano, K., G. W. Imbens, and G. Ridder (2003), "Efficient Estimation of Average Treatment Effects Using the Estimated Propensity Score," Econometrica 71, 1161-1189.\\
Holtz-Eakin, D., W. Newey, and H. S. Rosen (1988), "Estimating Vector Autoregressions with Panel Data," Econometrica 56, 1371-1395.\\
Holzer, H., R. Block, M. Cheatham, and J. Knott (1993), "Are Training Subsidies Effective? The Michigan Experience," Industrial and Labor Relations Review 46, 625-636.\\
Honoré, B. E. (1992), "Trimmed LAD and Least Squares Estimation of Truncated and Censored Regression Models with Fixed Effects," Econometrica 60, 533-565.\\
Honoré, B. E. (1993a), "Orthogonality Conditions for Tobit Models with Fixed Effects and Lagged Dependent Variables," Journal of Econometrics 59, 35-61.\\
Honoré, B. E. (1993b), "Identification Results for Duration Models with Multiple Spells," Review of Economic Studies 60, 241-246.\\
Honoré, B. E., and L. Hu (2004), "Estimation of Cross Sectional and Panel Data Censored Regression Models with Endogeneity," Journal of Econometrics 122, 293-316.

Honoré, B. E., and E. Kyriazidou (2000a), "Panel Data Discrete Choice Models with Lagged Dependent Variables," Econometrica 68, 839-874.\\
Honoré, B. E., and E. Kyriazidou (2000b), "Estimation of Tobit-Type Models with Individual Specific Effects," Econometric Reviews 19, 341-366.\\
Honoré, B. E., E. Kyriazidou, and C. Udry (1997), "Estimation of Type 3 Tobit Models Using Symmetric Trimming and Pairwise Comparisons," Journal of Econometrics 76, 107-128.\\
Horowitz, J. L. (1992), "A Smoothed Maximum Score Estimator for the Binary Response Model," Econometrica 60, 505-531.\\
Horowitz, J. L. (1999), "Semiparametric Estimation of a Proportional Hazard Model with Unobserved Heterogeneity," Econometrica 67, 1001-1028.\\
Horowitz, J. L. (2001), "The Bootstrap," in Handbook of Econometrics, Volume 5, ed. J. J. Heckman and E. Leamer. Amsterdam: North Holland, 3159-3228.

Horowitz, J. L., and C. F. Manski (1998), "Censoring of Outcomes and Regressors Due to Survey Nonresponse: Identification and Estimation Using Weights and Imputations," Journal of Econometrics 84, 37-58.\\
Horvitz, D., and D. Thompson (1952), "A Generalization of Sampling without Replacement from a Finite Population," Journal of the American Statistical Association 47, 663-685.\\
Hoxby, C. M. (1996), "How Teachers' Unions Affect Education Production," Quarterly Journal of Economics 111, 671-718.\\
Hoxby, C. M. (2000), "Does Competition among Public Schools Benefit Students and Taxpayers?" American Economic Review 90, 1209-1238.\\
Hsiao, C. (1986), Analysis of Panel Data. Cambridge: Cambridge University Press.\\
Hsiao, C. (2003), Analysis of Panel Data, second edition. Cambridge: Cambridge University Press.\\
Huber, P. J. (1967), "The Behavior of Maximum Likelihood Estimates under Nonstandard Conditions," in Proceedings of the Fifth Berkeley Symposium in Mathematical Statistics, Volume 1. Berkeley: University of California Press, 221-233.\\
Huber, P. J. (1981), Robust Statistics. New York: Wiley.\\
Ichimura, H. (1993), "Semiparametric Least Squares (SLS) and Weighted SLS Estimation of Single-Index Models," Journal of Econometrics 58, 71-120.\\
Im, K. S., S. C. Ahn, P. Schmidt, and J. M. Wooldridge (1999), "Efficient Estimation of Panel Data Models with Strictly Exogenous Explanatory Variables," Journal of Econometrics 93, 177-201.\\
Imbens, G. W. (1992), "An Efficient Method of Moments Estimator for Discrete Choice Models with Choice-Based Sampling," Econometrica 60, 1187-1214.\\
Imbens, G. W. (2000), "The Role of the Propensity Score in Estimating Dose-Response Functions," Biometrika 87, 706-710.\\
Imbens, G. W. (2004), "Nonparametric Estimation of Average Treatment Effects under Exogeneity: A Review," Review of Economics and Statistics 86, 4-29.\\
Imbens, G. W., and J. D. Angrist (1994), "Identification and Estimation of Local Average Treatment Effects," Econometrica 62, 467-476.\\
Imbens, G. W., and K. Kalyanaraman (2009), "Optimal Bandwidth Choice for the Regression Discontinuity Estimator," National Bureau of Economic Research Working Paper Number 14726.\\
Imbens, G. W., and T. Lancaster (1996), "Efficient Estimation and Stratified Sampling," Journal of Econometrics 74, 289-318.\\
Imbens, G. W., and T. Lemieux (2008), "Regression Discontinuity Designs: A Guide to Practice," Journal of Econometrics 142, 615-635.\\
Imbens, G. W., and D. B. Rubin (forthcoming), Causal Inference in Statistics. Cambridge: Cambridge University Press.

Imbens, G. W., and J. M. Wooldridge (2009), "Recent Developments in the Econometrics of Program Evaluation," Journal of Economic Literature 47(1), 5-86.\\
Kahn, S., and K. Lang (1988), "Efficient Estimation of Structural Hedonic Systems," International Economic Review 29, 157-166.\\
Kakwani, N. (1967), "The Unbiasedness of Zellner's Seemingly Unrelated Regressions Equation Estimators," Journal of the American Statistical Association 62, 141-142.\\
Kalbfleisch, J. D., and R. L. Prentice (1980), The Statistical Analysis of Failure Time Data. New York: Wiley.\\
Kane, T. J., and C. E. Rouse (1995), "Labor-Market Returns to Two- and Four-Year Colleges," American Economic Review 85, 600-614.\\
Kao, C. (1999), "Spurious Regression and Residual-Based Tests for Cointegration in Panel Data," Journal of Econometrics 90, 1-44.\\
Keane, M. P. (1993), "Simulation Estimation for Panel Data Models with Limited Dependent Variables," in Handbook of Statistics, Volume 11, ed. G. S. Maddala, C. R. Rao, and H. D. Vinod. Amsterdam: North Holland, 545-571.\\
Keane, M. P., and R. A. Moffitt (1998), "A Structural Model of Multiple Welfare Participation and Labor Supply," International Economic Review 39, 553-589.\\
Keane, M. P., and D. E. Runkle (1992), "On the Estimation of Panel Data Models with Serial Correlation When Instruments Are Not Strictly Exogenous," Journal of Business and Economic Statistics 10, 1-9.\\
Keane, M. P., and K. I. Wolpin (1997), "The Career Decisions of Young Men," Journal of Political Economy 105, 473-522.\\
Kiefer, N. M. (1980), "Estimation of Fixed Effect Models for Time Series of Cross-Sections with Arbitrary Intertemporal Covariance," Journal of Econometrics 14, 195-202.\\
Kiefer, N. M. (1988), "Economic Duration Data and Hazard Functions," Journal of Economic Literature 26, 646-679.\\
Kiefer, N. M. (1989), "The ET Interview: Arthur S. Goldberger," Econometric Theory 5, 133-160.\\
Kiel, K. A., and K. T. McClain (1995), "House Prices during Siting Decision Stages: The Case of an Incinerator from Rumor through Operation," Journal of Environmental Economics and Management 28, 241-255.\\
Kieschnick, R., and B. D. McCullough (2003), "Regression Analysis of Variates Observed on ( 0,1 ): Percentages, Proportions and Fractions," Statistical Modelling 3, 193-213.\\
Kinal, T. W. (1980), "The Existence of Moments of k-Class Estimators," Econometrica 48, 241-249.\\
Klein, R. W., and R. H. Spady (1993), "An Efficient Semiparametric Estimator for Discrete Choice Models," Econometrica 61, 387-421.\\
Koenker, R. (1981), "A Note on Studentizing a Test for Heteroskedasticity," Journal of Econometrics 17, 107-112.\\
Koenker, R. (2004), "Quantile Regression for Longitudinal Data," Journal of Multivariate Analysis 91, 74-89.

Koenker, R. (2005), Quantile Regression. Cambridge: Cambridge University Press.\\
Koenker, R., and G. Bassett (1978), "Regression Quantiles," Econometrica 46, 33-50.\\
Krueger, A. B. (1993), "How Computers Have Changed the Wage Structure: Evidence from Microdata, 1984-1989," Quarterly Journal of Economics 108, 33-60.\\
Kuksov, D., and J. M. Villas-Boas (2008), "Endogeneity and Individual Consumer Choice," Journal of Marketing Research 45, 702-714.\\
Kyriazidou, E. (1997), "Estimation of a Panel Data Sample Selection Model," Econometrica 65, 1335-1364.

Lahiri, K., and P. Schmidt (1978), "On the Estimation of Triangular Structural Systems," Econometrica 46, 1217-1221.\\
Lalonde, R. J. (1986), "Evaluating the Econometric Evaluations of Training Programs with Experimental Data," American Economic Review 76, 604-620.\\
Lancaster, T. (1979), "Econometric Methods for the Duration of Unemployment," Econometrica 47, 939956.

Lancaster, T. (1990), The Econometric Analysis of Transition Data. Cambridge: Cambridge University Press.\\
LeCam, L. (1953), "On Some Asymptotic Properties of Maximum Likelihood Estimates and Related Bayes Estimates," University of California Publications in Statistics 1, 277-328.\\
Lechner, M. (2001), "Identification and Estimation of Causal Effects of Multiple Treatments under the Conditional Independence Assumption," in Econometric Evaluation of Labour Market Policies, ed. M. Lechner, and F. Pfeiffer. Heidelberg: Physica, 43-58.\\
Lechner, M. (2004), "Sequential Matching Estimation of Dynamic Causal Effects," Discussion Paper 2004-06, Department of Economics, University of St. Gallen. St. Gallen, Switzerland.\\
Lechner, M., and R. Miquel (2005), "Identification of Effects of Dynamic Treatments by Sequential Conditional Independence Assumptions," Discussion Paper 2005-17, Department of Economics, University of St. Gallen. St. Gallen, Switzerland.\\
Lemieux, T. (1998), "Estimating the Effects of Unions on Wage Inequality in a Panel Data Model with Comparative Advantage and Nonrandom Selection," Journal of Labor Economics 16, 261-291.\\
Leung, S. F., and S. Yu (1996), "On the Choice between Sample Selection and Two-Part Models," Journal of Econometrics 72, 197-229.\\
Levinsohn, J., and A. Petrin (2003), "Estimating Production Functions Using Inputs to Control for Unobservables," Review of Economic Studies 70, 317-341.\\
Levitt, S. D. (1996), "The Effect of Prison Population Size on Crime Rates: Evidence from Prison Overcrowding Legislation," Quarterly Journal of Economics 111, 319-351.\\
Levitt, S. D. (1997), "Using Electoral Cycles in Police Hiring to Estimate the Effect of Police on Crime," American Economic Review 87, 270-290.\\
Lewbel, A. (2000), "Semiparametric Qualitative Response Model Estimation with Unknown Heteroscedasticity or Instrumental Variables," Journal of Econometrics 97, 145-177.\\
Li, Q., and J. S. Racine (2007), Nonparametric Econometrics. Princeton, NJ: Princeton University Press.\\
Li, Q., J. S. Racine, and J. M. Wooldridge (2008), "Estimating Average Treatment Effects with Continuous and Discrete Covariates: The Case of Swan-Ganz Catheterization," American Economic Review 98, 357-362.\\
Li, Q., J. S. Racine, and J. M. Wooldridge (2009), "Efficient Estimation of Average Treatment Effects with Mixed Categorical and Continuous Data," Journal of Business and Economic Statistics 27, 206-223.\\
Liang, Y.-K., and S. L. Zeger (1986), "Longitudinal Data Analysis Using Generalized Linear Models," Biometrika 73, 13-22.\\
Little, R. J. A., and D. B. Rubin (2002), Statistical Analysis with Missing Data, second edition. Hoboken, NJ: Wiley.\\
Loeb, S., and J. Bound (1996), "The Effect of Measured School Inputs on Academic Achievement: Evidence from the 1920s, 1930s and 1940s Birth Cohorts," Review of Economics and Statistics 78, 653-664.\\
Long, J. S., and J. Freese (2001), Regression Models for Categorical Dependent Variables Using Stata. College Station, TX: Stata Press.\\
MacKinnon, J. G., and H. White (1985), "Some Heteroskedasticity Consistent Covariance Matrix Estimators with Improved Finite Sample Properties," Journal of Econometrics 29, 305-325.\\
MaCurdy, T. E. (1982), "The Use of Time Series Processes to Model the Error Structure of Earnings in a Longitudinal Data Analysis," Journal of Econometrics 18, 83-114.

Maddala, G. S. (1983), Limited Dependent and Qualitative Variables in Econometrics. Cambridge: Cambridge University Press.\\
Maloney, M. T., and R. E. McCormick (1993), "An Examination of the Role That Intercollegiate Athletic Participation Plays in Academic Achievement: Athlete's Feats in the Classroom," Journal of Human Resources 28, 555-570.\\
Manski, C. F. (1975), "Maximum Score Estimation of the Stochastic Utility Model of Choice," Journal of Econometrics 3, 205-228.\\
Manski, C. F. (1987), "Semiparametric Analysis of Random Effects Linear Models from Binary Panel Data," Econometrica 55, 357-362.\\
Manski, C. F. (1988), Analog Estimation Methods in Econometrics. New York: Chapman and Hall.\\
Manski, C. F. (1996), "Learning about Treatment Effects from Experiments with Random Assignment of Treatments," Journal of Human Resources 31, 709-733.\\
Manski, C. F., and S. Lerman (1977), "The Estimation of Choice Probabilities from Choice-Based Samples," Econometrica 45, 1977-1988.\\
Manski, C. F., and D. McFadden (1981), "Alternative Estimators and Sample Designs for Discrete Choice Analysis," in Structural Analysis of Discrete Data with Econometric Applications, ed. C. F. Manski and D. McFadden. Cambridge, MA: MIT Press, 2-50.

McCall, B. P. (1994), "Testing the Proportional Hazards Assumption in the Presence of Unmeasured Heterogeneity," Journal of Applied Econometrics 9, 321-334.\\
McCullagh, P., and J. A. Nelder (1989), Generalized Linear Models, second edition. New York: Chapman and Hall.\\
McDonald, J. B. (1996), "An Application and Comparison of Some Flexible Parametric and SemiParametric Qualitative Response Models," Economics Letters 53, 145-152.\\
McDonald, J. F., and R. A. Moffitt (1980), "The Uses of Tobit Analysis," Review of Economics and Statistics 62, 318-321.\\
McFadden, D., and K. Train (2000), "Mixed MNL Models for Discrete Response," Journal of Applied Econometrics 15, 447-470.\\
McFadden, D. L. (1974), "Conditional Logit Analysis of Qualitative Choice Analysis," in Frontiers in Econometrics, ed. P. Zarembka. New York: Academic Press, 105-142.\\
McFadden, D. L. (1978), "Modeling the Choice of Residential Location," in Spatial Interaction Theory and Residential Location, ed. A. Karlqvist. Amsterdam: North Holland, 75-96.\\
McFadden, D. L. (1981), "Econometric Models of Probabilistic Choice," in Structural Analysis of Discrete Data with Econometric Applications, ed. C. F. Manski and D. McFadden. Cambridge, MA: MIT Press, 198-272.\\
McFadden, D. L. (1984), "Econometric Analysis of Qualitative Response Models," in Handbook of Econometrics, Volume 2, ed. Z. Griliches and M. D. Intriligator. Amsterdam: North Holland, 1395-1457.\\
McFadden, D. L. (1987), "Regression Based Specification Tests for the Multinomial Logit Model," Journal of Econometrics 34, 63-82.\\
Meyer, B. D. (1990), "Unemployment Insurance and Unemployment Spells," Econometrica 58, 757-782.\\
Meyer, B. D., W. K. Viscusi, and D. L. Durbin (1995), "Workers' Compensation and Injury Duration: Evidence from a Natural Experiment," American Economic Review 85, 322-340.\\
Model, K. E. (1993), "The Effect of Marijuana Decriminalization on Hospital Emergency Drug Episodes: 1975-1978," Journal of the American Statistical Association 88, 737-747.\\
Moffitt, R. A. (1996), "Identification of Causal Effects Using Instrumental Variables: Comment," Journal of the American Statistical Association 91, 462-465.\\
Moffitt, R., J. Fitzgerald, and P. Gottschalk (1999), "Sample Attrition in Panel Data: The Role of Selection on Observables," Annale d'Economie et de Statistique 55/56, 129-152.

Montgomery, E., K. Shaw, and M. E. Benedict (1992), "Pensions and Wages: An Hedonic Price Theory Approach," International Economic Review 33, 111-128.\\
Moon, C.-G. (1988), "Simultaneous Specification Test in a Binary Logit Model: Skewness and Heteroskedasticity," Communications in Statistics 17, 3361-3387.\\
Moon, H. R., and P. C. B. Phillips (2000), "Estimation of Autoregressive Roots Near Unity Using Panel Data," Econometric Theory 16, 927-997.\\
Moulton, B. (1990), "An Illustration of a Pitfall in Estimating the Effects of Aggregate Variables on Micro Units," Review of Economics and Statistics 72, 334-338.\\
Mroz, T. A. (1987), "The Sensitivity of an Empirical Model of Married Women's Hours of Work to Economic and Statistical Assumptions," Econometrica 55, 765-799.\\
Mullahy, J. (1997), "Instrumental-Variable Estimation of Count Data Models: Applications to Models of Cigarette Smoking Behavior," Review of Economics and Statistics 79, 586-593.\\
Mullahy, J. (1998), "Much Ado about Two: Reconsidering Retransformation and the Two-Part Model in Health Econometrics," Journal of Health Economics 17, 247-281.\\
Mundlak, Y. (1978), "On the Pooling of Time Series and Cross Section Data," Econometrica 46, 69-85.\\
Murtazashvili, I., and J. M. Wooldridge (2008), "Fixed Effects Instrumental Variables Estimation in Correlated Random Coefficient Panel Data Models," Journal of Econometrics 142, 539-552.\\
Newey, W. K. (1984), "A Method of Moments Interpretation of Sequential Estimators," Economics Letters 14, 201-206.\\
Newey, W. K. (1985), "Maximum Likelihood Specification Testing and Conditional Moment Tests," Econometrica 53, 1047-1070.\\
Newey, W. K. (1990), "Efficient Instrumental Variables Estimation of Nonlinear Models," Econometrica 58, 809-837.\\
Newey, W. K. (1993), "Efficient Estimation of Models with Conditional Moment Restrictions," in Handbook of Statistics, Volume 11, ed. G. S. Maddala, C. R. Rao, and H. D. Vinod. Amsterdam: North Holland, 419-454.\\
Newey, W. K., and D. McFadden (1994), "Large Sample Estimation and Hypothesis Testing," in Handbook of Econometrics, Volume 4, ed. R. F. Engle and D. McFadden. Amsterdam: North Holland, 2111-2245.\\
Newey, W. K., J. L. Powell, and F. Vella (1999), "Nonparametric Estimation of Triangular Simultaneous Equations Models," Econometrica 67, 565-603.\\
Newey, W. K., and K. D. West (1987), "A Simple, Positive Semi-Definite Heteroskedasticity and Autocorrelation Consistent Covariance Matrix," Econometrica 55, 703-708.\\
Nickell, S. (1979), "Estimating the Probability of Leaving Unemployment," Econometrica 47, 1249-1266.\\
Nijman, T., and M. Verbeek (1992), "Nonresponse in Panel Data: The Impact on Estimates of a Life Cycle Consumption Function," Journal of Applied Econometrics 7, 243-257.\\
Oaxaca, R. L., and M. R. Ransom (1994), "On Discrimination and the Decomposition of Wage Differentials," Journal of Econometrics 61, 5-21.\\
Olley, S., and A. Pakes (1996), "The Dynamics of Productivity in the Telecommunications Equipment Industry," Econometrica 64, 1263-1298.\\
Orme, C. (1990), "The Small Sample Performance of the Information Matrix Test," Journal of Econometrics 46, 309-331.\\
Orme, C. D., and T. Yamagata (2006), "The Asymptotic Distribution of the F-Test Statistic for Individual Effects," Econometrics Journal 9, 404-422.\\
Paarsch, H. J., and H. Hong (2006), An Introduction to the Structural Econometrics of Auction Data. Cambridge, MA: MIT Press.

Pagan, A. R. (1984), "Econometric Issues in the Analysis of Regressions with Generated Regressors," International Economic Review 25, 221-247.\\
Pagan, A. R., and F. Vella (1989), "Diagnostic Tests for Models Based on Individual Data: A Survey," Journal of Applied Econometrics 4, S29-S59.\\
Page, M. (1995), "Racial and Ethnic Discrimination in Urban Housing Markets: Evidence from a Recent Audit Study," Journal of Urban Economics 38, 183-206.\\
Papke, L. E. (1991), "Interstate Business Tax Differentials and New Firm Location," Journal of Public Economics 45, 47-68.\\
Papke, L. E. (1994), "Tax Policy and Urban Development: Evidence from the Indiana Enterprise Zone Program," Journal of Public Economics 54, 37-49.\\
Papke, L. E. (1998), "How Are Participants Directing Their Participant-Directed Individual Account Pension Plans?" American Economic Review 88, 212-216.\\
Papke, L. E. (2005), "The Effects of Spending on Test Pass Rates: Evidence from Michigan," Journal of Public Economics 89, 821-839.\\
Papke, L. E., and J. M. Wooldridge (1996), "Econometric Methods for Fractional Response Variables with an Application to 401(k) Plan Participation Rates," Journal of Applied Econometrics 11, 619-632.\\
Papke, L. E., and J. M. Wooldridge (2008), "Panel Data Methods for Fractional Response Variables with an Application to Test Pass Rates," Journal of Econometrics 145, 121-133.\\
Pearl, J. (2000), Causality: Models, Reasoning, and Inference. Cambridge: Cambridge University Press.\\
Perracchi, F. (2001), Econometrics. Chichester: Wiley.\\
Pesaran, M. H., and R. J. Smith (1995), "Estimating Long-Run Relationships from Dynamic Heterogeneous Panels," Journal of Econometrics 68, 79-113.\\
Petrin, A., and K. Train (2010), "A Control Function Approach to Endogeneity in Consumer Choice Models," Journal of Marketing Research 47, 3-13.\\
Phillips, P. C. B., and H. R. Moon (1999), "Linear Regression Limit Theory for Nonstationary Panel Data," Econometrica 67, 1057-1111.\\
Phillips, P. C. B., and H. R. Moon (2000), "Nonstationary Panel Data Analysis: An Overview of Some Recent Developments," Econometric Reviews 19, 263-268.\\
Phillips, P. C. B., and J. Y. Park (1988), "On the Formulation of Wald Tests for Nonlinear Restrictions," Econometrica 56, 1065-1083.\\
Polachek, S., and M.-K. Kim (1994), "Panel Estimates of the Gender Earnings Gap: Individual-Specific Intercepts and Individual-Specific Slope Models," Journal of Econometrics 61, 23-42.\\
Porter, J. (2002), "Efficiency of Covariance Matrix Estimators for Maximum Likelihood Estimation," Journal of Business and Economic Statistics 20, 431-440.\\
Powell, J. L. (1984), "Least Absolute Deviations Estimation for the Censored Regression Model," Journal of Econometrics 25, 303-325.\\
Powell, J. L. (1986), "Symmetrically Trimmed Least Squares Estimation for Tobit Models," Econometrica 54, 1435-1460.\\
Powell, J. L. (1991), "Estimation of Monotonic Regression Models under Quantile Restrictions," in Nonparametric and Semiparametric Methods in Econometrics and Statistics: Proceedings of the Fifth International Symposium in Economic Theory and Econometrics, ed. W. A. Barnett, J. L. Powell, and G. E. Cambridge: Cambridge University Press, 357-384.\\
Powell, J. L. (1994), "Estimation of Semiparametric Models," in Handbook of Econometrics, Volume 4, ed. R. F. Engle and D. McFadden. Amsterdam: North Holland, 2443-2521.

Powell, J. L., J. H. Stock, and T. M. Stoker (1989), "Semiparametric Estimation of Weighted Average Derivatives," Econometrica 57, 1403-1430.

Qian, H., and P. Schmidt (1999), "Improved Instrumental Variables and Generalized Method of Moments Estimators," Journal of Econometrics 91, 145-169.\\
Quah, D. (1994), "Exploiting Cross-Section Variations for Unit Root Inference in Dynamic Data,"Economics Letters 44, 9-19.\\
Quandt, R. E. (1983), "Computational Problems and Methods," in Handbook of Econometrics, Volume 1, ed. Z. Griliches and M. D. Intriligator. Amsterdam: North Holland, 699-764.\\
Ramsey, J. B. (1969), "Tests for Specification Errors in Classical Linear Least Squares Regression Analysis," Journal of the Royal Statistical Society, Series B, 31, 350-371.\\
Rao, C. R. (1948), "Large Sample Tests of Hypotheses Involving Several Parameters with Applications to Problems of Estimation," Proceedings of the Cambridge Philosophical Society 44, 50-57.\\
Raudenbush, S. W., and A. S. Bryk (2002), Hierarchical Linear Models: Applications and Data Analysis, second edition. Thousand Oaks, CA: Sage.\\
Rivers, D., and Q. H. Vuong (1988), "Limited Information Estimators and Exogeneity Tests for Simultaneous Probit Models," Journal of Econometrics 39, 347-366.\\
Robins, J. M., and A. Rotnitzky (1995), "Semiparametric Efficiency in Multivariate Regression Models," Journal of the American Statistical Association 90, 122-129.\\
Robins, J. A., A. Rotnitzky, and L. Zhao (1995), "Analysis of Semiparametric Regression Models for Repeated Outcomes in the Presence of Missing Data," Journal of the American Statistical Association 90, 106-121.\\
Robinson, P. M. (1988), "Root-N-Consistent Semiparametric Regression," Econometrica 56, 931-954.\\
Romer, D. (1993), "Openness and Inflation: Theory and Evidence," Quarterly Journal of Economics 108, 869-903.\\
Rose, N. L. (1990), "Profitability and Product Quality: Economic Determinants of Airline Safety Performance," Journal of Political Economy 98, 944-961.\\
Rosenbaum, P. R., and D. B. Rubin (1983), "The Central Role of the Propensity Score in Observational Studies for Causal Effects," Biometrika 70, 41-55.\\
Rouse, C. E. (1995), "Democratization or Diversion? The Effect of Community Colleges on Educational Attainment," Journal of Business and Economic Statistics 13, 217-224.\\
Rubin, D. B. (1974), "Estimating Causal Effects of Treatments in Randomized and Nonrandomized Studies," Journal of Educational Psychology 66, 688-701.\\
Rubin, D. B. (1976), "Inference and Missing Data," Biometrika 63, 581-592.\\
Rudin, W. (1976), Principles of Mathematical Analysis, 3rd edition. New York: McGraw-Hill.\\
Ruud, P. (1983), "Sufficient Conditions for Consistency of Maximum Likelihood Estimation Despite Misspecification of Distribution," Econometrica 51, 225-228.\\
Ruud, P. (1984), "Tests of Specification in Econometrics," Econometric Reviews 3, 211-242.\\
Ruud, P. (1986), "Consistent Estimation of Limited Dependent Variable Models Despite Misspecification of Distribution," Journal of Econometrics 32, 157-187.\\
Sander, W. (1992), "The Effect of Women's Schooling on Fertility," Economics Letters 40, 229-233.\\
Sargan, J. D. (1958), "The Estimation of Economic Relationships Using Instrumental Variables," Econometrica 26, 393-415.\\
Schmidt, P. (1976), Econometrics. New York: Marcel-Dekker.\\
Semykina, A., and J. M. Wooldridge (in press), "Estimating Panel Data Models in the Presence of Endogeneity and Selection," Journal of Econometrics.\\
Serfling, R. J. (2006), "Multivariate Symmetry and Asymmetry," in Encyclopedia of Statistical Sciences, ed. S. Kotz, N. Balakrishnan, C. B. Read, and B. Vidakovic. New York: Wiley, 5338-5345.\\
Shapiro, M. D. (1984), "The Permanent Income Hypothesis and the Real Interest Rate: Some Evidence from Panel Data," Economics Letters 14, 93-100.

Shea, J. (1995), "Union Contracts and the Life-Cycle/Permanent Income Hypothesis," American Economic Review 85, 186-200.\\
Shea, J. (1997), "Instrument Relevance in Multivariate Linear Models: A Simple Measure," Review of Economics and Statistics 79, 348-352.\\
Smith, R., and R. Blundell (1986), "An Exogeneity Test for a Simultaneous Equation Tobit Model with an Application to Labor Supply," Econometrica 54, 679-685.\\
Solon, G. (1985), "Comment on 'Benefits and Limitations of Panel Data' by C. Hsiao," Econometric Reviews 4, 183-186.\\
Staiger, D., and J. H. Stock (1997), "Instrumental Variables Regression with Weak Instruments," Econometrica 65, 557-586.\\
Stock, J. H., J. H. Wright, and M. Yogo (2002), "A Survey of Weak Instruments and Weak Identification in Generalized Method of Moments," Journal of Business and Economic Statistics 20, 518-529.\\
Stoker, T. M. (1986), "Consistent Estimation of Scaled Coefficients," Econometrica 54, 1461-1481.\\
Stoker, T. M. (1992), Lectures on Semiparametric Econometrics. Louvain-la-Neuve, Belgium: CORE Lecture Series.\\
Strauss, J., and D. Thomas (1995), "Human Resources: Empirical Modeling of Household and Family Decisions," in Handbook of Development Economics, Volume 3A, ed. J. Berhman and T. N. Srinivasan. Amsterdam: North Holland, 1883-2023.\\
Sueyoshi, G. T. (1992), "Semiparametric Proportional Hazards Estimation of Competing Risks Models with Time-Varying Covariates," Journal of Econometrics 51, 25-58.\\
Sueyoshi, G. T. (1995), "A Class of Binary Response Models for Grouped Duration Data," Journal of Applied Econometrics 10, 411-431.\\
Swait, J. (2003), "Flexible Covariance Structures for Categorical Dependent Variables through Finite Mixtures of Generalized Extreme Value Models," Journal of Business and Economic Statistics 21, 80-87.\\
Tauchen, G. (1985), "Diagnostic Testing and Evaluation of Maximum Likelihood Models," Journal of Econometrics 30, 415-443.\\
Tauchen, G. (1986), "Statistical Properties of Generalized Method-of-Moments Estimators of Structural Parameters Obtained from Financial Market Data," Journal of Business and Economic Statistics 4, 397-416.\\
Terza, J. V. (1998), "Estimating Count Models with Endogenous Switching: Sample Selection and Endogenous Treatment Effects," Journal of Econometrics 84, 129-154.\\
Theil, H. (1983), "Linear Algebra and Matrix Methods in Econometrics," in Handbook of Econometrics, Volume 1, ed. Z. Griliches and M. D. Intriligator. Amsterdam: North Holland, 5-65.\\
Thomas, D., J. Strauss, and M.-H. Henriques (1990), "Child Survival, Height for Age and Household Characteristics in Brazil," Journal of Development Economics 33, 197-234.\\
Tobin, J. (1958), "Estimation of Relationships for Limited Dependent Variables," Econometrica 26, 24-36. van den Berg, G. (2001), "Duration Models: Specification, Identification, and Multiple Durations," in Handbook of Econometrics, Volume 5, ed. J. J. Heckman and E. Leamer. Amsterdam: North Holland, 3381-3460.\\
van der Klaauw, W. (1996), "Female Labour Supply and Marital Status Decisions: A Life-Cyle Model," Review of Economic Studies 63, 199-235.\\
van der Klaauw, W. (2002), "Estimating the Effect of Financial Aid Offers on College Enrollment: A Regression-Discontinuity Approach," International Economic Review 43, 1249-1287.\\
van der Laan, M. J., and J. M. Robins (2003), Unified Methods for Censored Longitudinal and Causality. New York: Springer-Verlag.\\
Vella, F. (1992), "Simple Tests for Sample Selection Bias in Censored and Discrete Choice Models," Journal of Applied Econometrics 7, 413-421.

Vella, F. (1993), "A Simple Estimator for Simultaneous Models with Censored Endogenous Regressors," International Economic Review 34, 441-457.\\
Vella, F. (1998), "Estimating Models with Sample Selection Bias: A Survey," Journal of Human Resources 33, 127-169.\\
Vella, F., and M. Verbeek (1998), "Whose Wages Do Unions Raise? A Dynamic Model of Unionism and Wage Rate Determination for Young Men," Journal of Applied Econometrics 13, 163-183.\\
Vella, F., and M. Verbeek (1999), "Estimating and Interpreting Models with Endogenous Treatment Effects," Journal of Business and Economic Statistics 17, 473-478.\\
Villas-Boas, J. M., and R. S. Winer (1999), "Endogeneity in Brand Choice Models," Management Science 10, 1324-1338.\\
Vuong, Q. (1989), "Likelihood Ratio Tests for Model Selection and Nonnested Hypotheses," Econometrica 57, 307-333.\\
Wald, A. (1940), "The Fitting of Straight Lines if Both Variables Are Subject to Error," Annals of Mathematical Statistics 11, 284-300.\\
White, H. (1980a), "Nonlinear Regression on Cross Section Data," Econometrica 48, 721-746.\\
White, H. (1980b), "A Heteroskedasticity-Consistent Covariance Matrix Estimator and a Direct Test for Heteroskedasticity," Econometrica 48, 817-838.\\
White, H. (1981), "Consequences and Detection of Misspecified Nonlinear Regression Models," Journal of the American Statistical Association 76, 419-433.\\
White, H. (1982a), "Maximum Likelihood Estimation of Misspecified Models," Econometrica 50, 1-26.\\
White, H. (1982b), "Instrumental Variables Regression with Independent Observations," Econometrica 50, 483-499.\\
White, H. (1984), Asymptotic Theory for Econometricians. Orlando, FL: Academic Press.\\
White, H. (1994), Estimation, Inference and Specification Analysis. Cambridge: Cambridge University Press.\\
White, H. (2001), Asymptotic Theory for Econometricians, revised edition. San Diego, CA: Academic Press.\\
Windmeijer, F. (2000), "Moment Conditions for Fixed Effects Count Data Models with Endogenous Regressors," Economics Letters 68, 21-24.\\
Wolak, F. A. (1991), "The Local Nature of Hypothesis Tests Involving Inequality Constraints in Nonlinear Models," Econometrica 59, 981-995.\\
Wooldridge, J. M. (1990), "A Unified Approach to Robust, Regression-Based Specification Tests," Econometric Theory 6, 17-43.\\
Wooldridge, J. M. (1991a), "On the Application of Robust, Regression-Based Diagnostics to Models of Conditional Means and Conditional Variances," Journal of Econometrics 47, 5-46.\\
Wooldridge, J. M. (1991b), "Specification Testing and Quasi-Maximum Likelihood Estimation," Journal of Econometrics 48, 29-55.\\
Wooldridge, J. M. (1992), "Some Alternatives to the Box-Cox Regression Model," International Economic Review 33, 935-955.\\
Wooldridge, J. M. (1994a), "Estimation and Inference for Dependent Processes," in Handbook of Econometrics, Volume 4, ed. R. F. Engle and D. L. McFadden. Amsterdam: North-Holland, 2639-2738.\\
Wooldridge, J. M. (1994b), "Efficient Estimation under Heteroskedasticity," Econometric Theory 10, 223.\\
Wooldridge, J. M. (1995a), "Selection Corrections for Panel Data Models under Conditional Mean Independence Assumptions," Journal of Econometrics 68, 115-132.\\
Wooldridge, J. M. (1995b), "Score Diagnostics for Linear Models Estimated by Two Stage Least Squares," in Advances in Econometrics and Quantitative Economics, ed. G. S. Maddala, P. C. B. Phillips, and T. N. Srinivasan. Oxford: Blackwell, 66-87.

Wooldridge, J. M. (1996), "Estimating Systems of Equations with Different Instruments for Different Equations," Journal of Econometrics 74, 387-405.\\
Wooldridge, J. M. (1997a), "Multiplicative Panel Data Models without the Strict Exogeneity Assumption," Econometric Theory 13, 667-678.\\
Wooldridge, J. M. (1997b), "On Two Stage Least Squares Estimation of the Average Treatment Effect in a Random Coefficient Model," Economics Letters 56, 129-133.\\
Wooldridge, J. M. (1997c), "Quasi-Likelihood Methods for Count Data," in Handbook of Applied Econometrics, Volume 2, ed. M. H. Pesaran and P. Schmidt. Oxford: Blackwell, 352-406.\\
Wooldridge, J. M. (1999a), "Distribution-Free Estimation of Some Nonlinear Panel Data Models," Journal of Econometrics 90, 77-97.\\
Wooldridge, J. M. (1999b), "Asymptotic Properties of Weighted M-Estimators for Variable Probability Samples," Econometrica 67, 1385-1406.\\
Wooldridge, J. M. (1999c), "Estimating Average Partial Effects under Conditional Moment Independence Assumptions," mimeo, Michigan State University Department of Economics.\\
Wooldridge, J. M. (2000), "A Framework for Estimating Dynamic, Unobserved Effects Panel Data Models with Possible Feedback to Future Explanatory Variables," Economics Letters 68, 245-250.\\
Wooldridge, J. M. (2001), "Asymptotic Properties of Weighted M-Estimators for Standard Stratified Samples." Econometric Theory 17, 451-470.\\
Wooldridge, J. M. (2003a), "Cluster-Sample Methods in Applied Econometrics," American Economic Review 93 (May), 133-138.\\
Wooldridge, J. M. (2003b), "Further Results on Instrumental Variables Estimation of Average Treatment Effects in the Correlated Random Coefficient Model," Economics Letters 79 (May), 185-191.\\
Wooldridge, J. M. (2004), "Estimating Average Partial Effects Under Conditional Moment Independence Assumptions," CeMMAP Working Paper Number CWP03/04.\\
Wooldridge, J. M. (2005a), "Fixed Effects and Related Estimators for Correlated Random-Coefficient and Treatment Effect Panel Data Models," Review of Economics and Statistics 87, 385-390.\\
Wooldridge, J. M. (2005b), "Simple Solutions to the Initial Conditions Problem for Dynamic, Nonlinear Panel Data Models with Unobserved Heterogeneity," Journal of Applied Econometrics 20, 39-54.\\
Wooldridge, J. M. (2005c), "Unobserved Heterogeneity and Estimation of Average Partial Effects," in Identification and Inference for Econometric Models: Essays in Honor of Thomas Rothenberg, ed. D. W. K. Andrews and J. H. Stock. Cambridge: Cambridge University Press, 2005, 27-55.\\
Wooldridge, J. M. (2005d), "Violating Ignorability of Treatment by Controlling for Too Many Factors," Econometric Theory 21, 1026-1028.\\
Wooldridge, J. M. (2005e), "Instrumental Variables Estimation with Panel Data,"Econometric Theory 21, 865-869.\\
Wooldridge, J. M. (2007), "Inverse Probability Weighted M-Estimation for General Missing Data Problems," Journal of Econometrics 141, 1281-1301.\\
Wooldridge, J. M. (2008), "Instrumental Variables Estimation of the Average Treatment Effect in Correlated Random Coefficient Models," in Advances in Econometrics, Volume 21, ed. D. Millimet, J. Smith, and E. Vytlacil. Amsterdam: Elsevier, 93-117.\\
Wooldridge, J. M. (2009a), Introductory Econometrics: A Modern Approach, fourth edition. Mason, OH: South-Western.\\
Wooldridge, J. M. (2009b), "On Estimating Firm-Level Production Functions Using Proxy Variables to Control for Unobservables," Economics Letters 104, 112-114.\\
Wooldridge, J. M. (2009c), "Should Instrumental Variables Be Used as Matching Variables?" Mimeo, Michigan State University Department of Economics.

Wu, D. (1973), "Alternative Tests of Independence between Stochastic Regressors and Disturbances," Econometrica 41, 733-750.\\
Yatchew, A., and Z. Griliches (1985), "Specification Error in Probit Models," Review of Economics and Statistics 67, 134-39.\\
Zeldes, S. P. (1989), "Consumption and Liquidity Constraints: An Empirical Investigation," Journal of Political Economy 97, 305-346.\\
Zellner, A. (1962), "An Efficient Method of Estimating Seemingly Unrelated Regressions and Tests of Aggregation Bias," Journal of the American Statistical Association 57, 500-509.\\
Ziliak, J. P. (1997), "Efficient Estimation with Panel Data When Instruments Are Predetermined: An Empirical Comparison of Moment-Condition Estimators," Journal of Business and Economic Statistics 15, 419-431.\\
Ziliak, J. P., and T. J. Kniesner (1998), "The Importance of Sample Attrition in Life Cycle Labor Supply Estimation," Journal of Human Resources 33, 507-530.\\
Ziliak, J. P., B. Wilson, and J. Stone (1999), "Spatial Dynamics and Heterogeneity in the Cyclicality of Real Wages," Review of Economics and Statistics 81, 227-236.

\section*{Index}
Absolutely continuous, 522\\
Admissible linear transformation, 246\\
Amount decision, 690-691\\
Analogy principle, ordinary least squares (OLS) analysis, 57-58\\
Asymmetric absolute loss function, 450\\
Asymptotic analysis, 7, 37-47\\
asymptotic inference, 454-459\\
asymptotically normal distribution, 40-41, 43-44, 175-176, 405-409, 411-413, 728-732\\
boundedness in probability, 37-40\\
convergence in distribution, 40-41\\
convergence in probability, $38-40$\\
convergence of deterministic sequences, 37\\
estimating asymptotic variance, 413-420\\
adjustments for two-step M-estimator, 418-420\\
without nuisance parameters, 413-418\\
estimators in linear unobserved effects panel data models, asymptotic properties, 284-285\\
generalized instrumental variables (GIV) estimator, 222-224\\
generalized least squares (GLS) analysis, 175-182\\
generalized method of moments (GMM) estimator, 525-530\\
inference with fixed effects, 304-307\\
limit theorems for random samples, 41-42\\
linear panel data model, 197-198\\
ordinary least squares (OLS) analysis, 54-65\\
asymptotic inference, 59-60\\
consistency, 56-58\\
heteroskedasticity-robust inference, 60-62\\
Lagrange multiplier (score) tests, 62-65, 421-428\\
system ordinary least squares (SOLS), 167-172\\
properties of estimators, 42-45, 172-182\\
properties of test statistics, 45-47\\
two-stage least squares (2SLS) analysis, 101-104\\
heteroskedasticity-robust inference, 106-107\\
Asymptotic efficiency, 44, 103-104, 131, 229-231\\
Asymptotic equivalence lemma, 41\\
Asymptotic estimator, 490-492\\
Asymptotic local power, 431\\
Asymptotic normality generalized least squares (GLS) analysis, 175-176\\
generalized method of moments (GMM) estimation, 527-528\\
M-estimators, 407-409, 411-413\\
maximum likelihood estimation (MLE), 476-479\\
nonlinear regression model, 405-409, 411-413\\
normal distribution, 40-41\\
Poisson QMLE, 728-732\\
two-stage least squares (2SLS) analysis, 101-102\\
Asymptotic refinement, 439-440\\
Asymptotic size, 45

Asymptotic variance\\
maximum likelihood estimation (MLE), 479-481, 862-863\\
nonlinear regression model, 413-420\\
Asymptotically pivotal statistic, 439-440\\
Attenuation bias, 81\\
Autonomy, simultaneous equations models (SEMs), 239, 247\\
Average causal effect. See Average treatment effect (ATE)\\
Average derivatives, 579\\
Average partial effects (APE), 22-25, 73, 141-142 binary response models, 577-582, 586-589, 615\\
in cluster sampling, 872\\
corner solution responses, 668, 674-675, 679\\
Poisson distribution, 726-727, 766-769\\
Average structural function (ASF), 24-25, 602-604, 875\\
Average treatment effect (ATE), 73, 578, 903-975\\
conditional on $x, 909$\\
counterfactual setting, 904-908\\
defining, 905-906\\
discrete or limited range responses, 960-961\\
ignorability of treatment, 908-936\\
instrumental variables methods, 937-954\\
control function approach, 948-951\\
correction function approaches, 945-948\\
estimating, using IV, 937-945\\
estimating local average treatment effect by IV, 951-954\\
multiple treatments, 964-967\\
multivalued treatments, 961-963\\
panel data, 968-975\\
regression discontinuity designs, 954-959\\
fuzzy regression discontinuity design, 957-959\\
sharp regression discontinuity design, 954-957\\
unconfoundedness of, 959\\
self-selection problem, 907-908\\
on the treated, 906\\
Balanced panels, 284\\
Baseline hazard, 988\\
Berndt, Hall, Hall, and Hausman (BHHH) optimization method, 433-434\\
Bernoulli distribution, 768\\
Bernoulli GLM variance assumption, 751-753\\
Between estimator, 304\\
Binary censoring, 780-782\\
Binary endogenous explanatory variable, 594-599\\
Binary response models, 471, 561-635\\
average partial effect (APE), 577-582, 586-587, 615\\
for panel data, 608-635\\
dynamic unobserved effects, 625-630

Binary response models (cont.)\\
pooled probit and logit, 609-610\\
probit models with heterogeneity and endogenous explanatory variables, 630-632\\
semiparametric approaches, 605, 632-635\\
unobserved effects logit models under strict exogeneity, 619-625\\
unobserved effects probit models under strict exogeneity, 610-619\\
index models, 565-582\\
logit, 566, 568, 573-582\\
maximum likelihood estimation (MLE), 567-569\\
probit, 566, 568, 573-582, 587, 591, 594-599\\
reporting results, 573-582\\
testing, 569-573\\
introduction, 561\\
linear probability model (LPM), 562-565\\
specification issues, 582-608\\
binary endogenous explanatory variable, 594-599\\
conditional maximum likelihood estimation (CMLE), 590-594\\
continuous endogenous explanatory variables, 585-594\\
estimation under weaker assumptions, 604-608\\
latent variable model, 599-604\\
neglected heterogeneity, 582-585\\
with sample selection, 813-814\\
Binomial GLM variance assumption, 740\\
Binomial regression model, 739-740\\
Bivariate probit model, 595-599\\
Bootstrap\\
bias method of, 438\\
$p$-value of, 440\\
standard error of, 438-439, 581, 590\\
variance estimate of, 438\\
Bootstrapping, 438-442, 603\\
critical values and, 439-440\\
Borel measurable function, 402\\
Boundedness in probability, 37-40\\
Breusch-Pagan test, 140\\
Burr distribution, 1005\\
Cauchy-Schwartz inequality, 323\\
Causal relationship\\
ceteris paribus analysis in, 3\\
control variables in, 3-4\\
establishing, 3-4\\
nature of, 3\\
Causality, simultaneous equations models (SEMs), 239-240\\
Censored least absolute deviations (CLAD) estimates, 689, 789-790

Censored normal regression model, 786-787\\
Censored responses. See Corner solution responses; Data censoring\\
Central limit theorem (CLT), 526\\
Demoivre-Laplace theorem, 40\\
statement of, 42\\
Ceteris paribus analysis, 3\\
Chamberlain approach to unobserved effects models, 347-349, 551, 616-619, 624\\
Chamberlain-Mundlak device, 662, 709, 766, 767, 788-789, 874\\
Chamberlain's correlated random effects probit model, 616-619, 624\\
Check function, 450\\
Chi-square estimator, 229-230\\
Choice-based sampling, 801-802, 854\\
Classical errors-in-variables (CEV) assumption, 80-82, 367\\
Classical linear model (CLM), 780-781, 886, 892\\
Classical minimum distance (CMD) estimation, 545-547\\
Cluster sampling, 6, 853, 863-894\\
generalized least squares (GLS) analysis, 866, 878\\
hierarchical linear models, 876-877, 879, 881\\
with large cluster sizes, 883-884\\
with large number of clusters, 864-876\\
with small cluster sizes, 864-876\\
with small number of clusters, 884-894\\
with unit-specific panel data, 876-883\\
Coefficient of interest, 150-151\\
Coefficient of variation, 741-742\\
Competing risks models, 1019\\
Complex survey sampling, 894-899\\
asymptotic properties of regression, 896-897\\
M-estimators, 897-899\\
multistage sampling, 894-895\\
Composite errors, 291\\
Concentrated objective function, 435-436\\
Conditional covariance, 33-34\\
Conditional density, 492-494, 522\\
Conditional distribution, 522-524\\
Conditional expectations (CE), 13-36\\
conditional covariance, 33-34\\
conditional variance, 32-33\\
corner solution responses, 691\\
examples, 14-15\\
features, 14-25\\
average partial effects, 22-25\\
elasticity, 16-18\\
error form of models, 18-19\\
partial effects, 15-16, 22-25\\
properties, 19-22\\
semielasticity, 18\\
law of iterated expectations (LIE), 18-22, 30-32\\
linear projections, 25-27, 34-36\\
models with unobserved effects, 758-759\\
properties, 19-22, 30-36\\
role in econometrics, 13-14\\
structural, 13\\
variable types, 13, 14, 23-24\\
Conditional hazard, 1004\\
Conditional independence, 908\\
Conditional independence assumption, 500-502\\
Conditional information matrix equality (CIME), 478-479\\
Conditional Jensen's inequality, 31-32\\
Conditional Kullback-Leibler information inequality, 473-474\\
Conditional $\log$ likelihood for observation $i, 474$\\
Conditional logit estimators, 621-624\\
Conditional logit (CL) model, 647-649\\
Conditional maximum likelihood estimation (CMLE), 474-476, 540-543\\
binary response models, 590-594\\
described, 470-471\\
general framework, 473-476\\
models with lagged dependent variables, 498-499\\
Conditional mean independence, 908\\
Conditional moment restrictions, 542-545\\
Conditional moment tests, 485\\
Conditional probit model, 649\\
Conditional variance, 32-33\\
matrix of, 172\\
Consistency\\
consistent estimators, 42-45\\
instrumental variables estimator, 142-143\\
of conditional maximum likelihood estimation (CMLE), 475-476\\
ordinary least squares (OLS) analysis, 56-58, 108\\
analogy principle, 57-58\\
method of moments, 57-58\\
population orthogonality condition, 56-57\\
quantile estimation, 449-453\\
two-stage least squares (2SLS) analysis, 98-100, 108, 142-143\\
Contemporaneous exogeneity, 164, 165\\
Continuous exogenous explanatory variable, 585-594\\
average partial effects (APE), 586-589\\
bootstrap standard error, 438-439, 501, 590\\
control function (CF) approach, 590-594\\
limited information procedure, 591-592\\
maximum likelihood estimation (MLE), 591-594\\
probit statistic, 587\\
two-step estimation, 587-588

Continuous function, 402-403\\
Continuous mapping theorem, 41\\
Control function (CF) approach\\
average treatment effect (ATE) analysis, 948-951\\
continuous exogenous explanatory variable, 590-594\\
control function estimation for triangular systems, 268-271\\
to endogeneity in single-equation linear models, 126-129\\
Garen model, 145-146\\
nature of, 126, 271\\
Control group, difference-in-differences (DD) estimation, 147-148, 150\\
Control variable ( $x$ ), 13\\
establishing, 3-4\\
nature of, 3\\
Convergence\\
in distribution, 40-41\\
in probability, 38-40\\
of deterministic sequences, 37\\
Corner solution model, 667\\
Corner solution outcomes. See Corner solution responses\\
Corner solution responses, 667-715\\
examples, 669-671\\
motivation, 667-669\\
panel data methods, 705-715\\
dynamic unobserved effects Tobit models, 713-715\\
pooled methods, 705-707\\
unobserved effects models under strict exogeneity, 707-713\\
two-limit Tobit model, 703-705\\
type I Tobit model, 670-689\\
estimation, 676-677\\
inference, 676-677\\
reporting results, 677-680\\
specification issues, 680-689\\
useful expressions, 671-676\\
type II Tobit model, 690-703\\
exponential, 697-703\\
exponential conditional mean, 694-696\\
lognormal hurdle model, 694-696\\
specification issues, 680-689\\
truncated normal hurdle model, 692-695\\
Correction function, average treatment effect (ATE) analysis, 945-948\\
Correctly specified model\\
for the conditional density, 474\\
for the conditional mean, 397-398\\
Correlated random coefficiency (CRC) model, 142-145

Correlated random effects (CRE)\\
estimation, 287\\
framework, 286-287, 384-387, 460-461, 496, 616-619, 624, 662-663, 872-873\\
Tobit model, 708, 711-713\\
Correlated random slopes, 384-387\\
Count variables, 472, 723-724. See also Poisson regression\\
Covariance\\
conditional, 33-34\\
restrictions in simultaneous equations models (SEMs), 257-260\\
Covariate ( $x$ ). See Explanatory variable ( $x$ )\\
Cox's partial likelihood method for proportional hazard model, 1018\\
Cragg's truncated normal hurdle model, 695, 707\\
Criterion function statistic, 227, 429-430\\
Cross equation correlation, 172, 188-189\\
Cross equation restrictions, 256-257\\
Cross section data, 5\\
pooled cross section, 5-6, 146-147\\
Cut points, 655\\
Data censoring, 777-790\\
binary, 780-782\\
censored flow data, 993-1000\\
censored least absolute deviations (CLAD), 689, 789-790\\
censored normal regression model, 786-787\\
from above and below, 785-790\\
interval coding, 783-785\\
left censoring, 1000\\
right censoring, 785-786, 993\\
sample selection. See Sample selection top coding, 779-780\\
Data truncation, 777\\
Delta method, 47\\
Demoivre-Laplace theorem, 40\\
Dependent variable ( $y$ ), 13, 14\\
fixed, 9-11\\
in ordinary least squares (OLS) analysis, 53-55, 73-76\\
measurement error, 76-78\\
lagged, 290, 371-374, 497-499\\
limited, 559\\
random coefficient models, 73-76\\
sample selection on basis of, 799-802, 809-813\\
Deterministic sequences, convergence in asymptotic analysis, 37\\
Difference-in-difference-in-differences (DDD) estimate, 150-151\\
Difference-in-differences (DD) estimation, 147-151, 321\\
Discrete responses, average treatment effect (ATE), 960-961

Dispersion parameter, 512\\
Distributed lag model, 290\\
Distribution, convergence in, 40-41\\
Disturbance. See Error terms\\
Disturbance term, conditional expectations, 18-19\\
Donald and Lang (DL) method of cluster sampling, 885-889, 893-894\\
Doubly robust, 930\\
Dummy endogenous variable model, 938\\
Dummy variable estimator, 308-310, 328\\
Dummy variable regression, 307-310\\
Duration analysis, 983-1019\\
competing risks model, 1019\\
Cox's partial likelihood method for proportional hazard model, 1018\\
grouped duration data, 1010-1018\\
time-invariant covariates, 1011-1015\\
time-varying covariates, 1015-1017\\
unobserved heterogeneity, 1017-1018\\
hazard functions, 984-991\\
conditional on time-invariant covariates, 988-989\\
conditional on time-varying covariates, 989-991\\
without covariates, 984-988\\
multiple-spell data, 1018-1019\\
single-spell data with time-invariant covariates, 991-1010\\
flow sampling, 992-993\\
maximum likelihood estimation with censored flow data, 993-1000\\
stock sampling, 1000-1003\\
unobserved heterogeneity, 1003-1010\\
survival analysis, 983\\
Duration dependence, 987\\
Durbin-Watson-Hausman (DWH) test, 129-132\\
Dynamic completeness\\
conditional density, 492-494\\
conditional on the unobserved effect, 371-374\\
of conditional mean, 194-196\\
of pooled probit and logit, 609-610\\
Dynamic ignorability, 970-971\\
Dynamic unobserved effects models\\
binary response models, 625-630\\
corner solution responses, 713-715\\
Efficiency\\
asymptotic, 44, 103-104, 131, 229-231\\
generalized method of moments (GMM) estimation, 538-545\\
conditional moment restrictions, 542-545\\
general efficiency framework, 538-540\\
maximum likelihood estimation, 540-542\\
relative efficiency, 539-540\\
linear simultaneous equations models (SEMs), 260-261\\
relative efficiency of two-stage least squares (2SLS) analysis, 103-104\\
Efficiency bound, 541\\
Elasticity, 17\\
conditional expectations, 16-18\\
Endogenous switching regression, 948\\
Endogenous variables, 54-55\\
control function approach in single-equation linear models, 126-129\\
explanatory, 585-594, 630-632, 651-653, 660-662, 681-685, 753-755, 809-813, 817-819\\
exponential regression function, 742-748\\
fractional responses, 753-755\\
identification in linear systems, 241-242\\
identification in nonlinear systems, 262-263\\
probit models with heterogeneity and endogenous explanatory variables, 630-632\\
specification tests, 129-134\\
types\\
measurement error, 55, 76-82\\
omitted, 54-55, 65-76\\
simultaneity, 55\\
Equivalent structures, 246\\
Error terms. See also Omitted variables\\
degrees-of-freedom correction, 62, 106-107\\
measurement errors, 55, 76-82\\
dependent variable in ordinary least squares (OLS) analysis, 76-78\\
explanatory variable in ordinary least squares (OLS) analysis, 78-82\\
multiplicative measurement error, 78\\
nature of, 8, 18\\
standard errors\\
heteroskedasticity-robust standard error, 61\\
Huber standard error, 61\\
two-stage least squares (2SLS) analysis, 108-112\\
White standard error, 61\\
Estimable model, ordinary least squares (OLS) estimation, 53\\
Exchangeable working correlation matrix, 448-449\\
Exclusion restrictions, in simultaneous equations models (SEMs), 241-243\\
Exogenous sampling, 795\\
Exogenous variables, 54\\
explanatory variables in sample selection, 802-808, 815-817\\
fractional responses, 748-753\\
identification in linear systems, 241-242\\
logit models under strict exogeneity, 619-625\\
panel data models with unobserved effects, 494-497\\
relaxing strict exogeneity assumption, 764-766\\
stratification based on, 861-863\\
unobserved effects of probit models under strict exogeneity, 610-619\\
Expected Hessian form of the LM statistic, 425\\
Experimental data, 5\\
Experimental group, difference-in-differences (DD) estimation, 147-148\\
Explained variable (y). See Dependent variable (y)\\
Explanatory variable ( $x$ ), 13, 14\\
binary endogenous, 594-599\\
continuous exogenous, 585-594\\
endogenous, 585-594, 630-632, 651-653, 660-662, 681-686, 753-755, 809-813, 817-819\\
exogenous, 585-594, 748-753\\
fixed, 9-11\\
instrumental variables (IV) approach and, 89-90\\
in ordinary least squares (OLS) analysis, measurement error, 78-82\\
probit models with heterogeneity and endogenous explanatory variables, 630-632\\
Exponential conditional mean, 694-696\\
Exponential distribution, 986-987\\
Exponential QMLE, 741\\
Exponential regression function, 397, 742-748\\
Exponential response function, 814\\
Exponential type II Tobit (ET2T) model, 697-703, 790-791, 804, 808\\
External covariates, 990\\
$F$ statistic\\
dummy variable regression, 309\\
from OLS analysis, 62, 104-105\\
for two-stage least squares (2SLS) analysis, 105, 112\\
Factor loads, 552\\
Feasible GLS (FGLS) estimator, 176-188, 200-201, 269-270\\
random effects analysis, 296-299\\
Finite distributed lag (FDL) model, 165-166\\
First differencing (FD) methods, 315-321\\
first-difference (FD) estimator, 316-318, 321-326\\
first difference instrumental variables (FDIV) estimator, 361-365\\
inference, 315-318\\
linear unobserved effects models, random trend model, 375-377\\
policy analysis, 320-321\\
robust variance matrix, 318-319\\
testing for serial correlation, 319-320\\
transformation and, 316-318\\
First-order asymptotic distribution, 407\\
First-stage regression, two-stage least squares (2SLS) analysis, 97, 105\\
Fixed effects (FE) estimators, 287, 300-304, 310, 332, 366, 495\\
first differencing estimators versus, 321-326

Fixed effects (FE) estimators (cont.)\\
in cluster sampling, 867-870, 876, 880\\
logit estimators, 621-623\\
random effects estimators versus, 326-334, 349-358\\
Fixed effects (FE) framework, 286, 300-315\\
asymptotic inference with fixed effects, 304-307\\
dummy variable regression, 307-310\\
estimators, 287, 300-304, 310\\
fixed effects generalized least squares (FEGLS), 300, 312-315\\
Hausman test comparing random and fixed effects estimators, 328-334, 355-358\\
Poisson estimation, 762-764\\
random trend model, 375-377\\
robust variance matrix estimator, 310-312\\
robustness of standard methods, 382-384\\
serial correlation, 310-311\\
time-varying coefficients on unobserved effects, 552-554\\
Fixed effects GLS (FEGLS), 300, 312-315\\
Fixed effects instrumental variables (FEIV) estimator, 353-358, 364-365\\
Fixed effects Poisson (FEP) estimator, 763\\
Fixed effects Poisson model, 756\\
Fixed effects residuals, 306-307\\
Fixed effects 2SLS (FE2SLS) estimator, 353-358\\
Fixed effects transformation, 302-303\\
Fixed explanatory variables, 9-11\\
Flow sampling, 992-993\\
Forbidden regression, 267-268\\
Fractional logit regression, 751\\
Fractional probit regression, 751\\
Fractional responses, 748-755\\
endogenous explanatory variables, 753-755\\
exogenous explanatory variables, 748-753\\
panel (longitudinal) data, 766-769\\
Fully recursive system, in simultaneous equations models (SEMs), 259-260\\
Fully robust variance matrix estimator, 416\\
Fuzzy regression discontinuity (FRD) design, 957-959

Gamma (exponential) regression model, 740-742\\
Gamma QMLE, 741\\
Gamma regression model, 741\\
Gamma-distributed heterogeneity, 1004\\
Gauss-Markov assumptions, 308\\
Gauss-Newton optimization method, 434-435\\
General linear system of equations, 210-213\\
Generalized condition information matrix equality (GCIME), 513-514\\
Generalized estimating equation (GEE), 446-447\\
for cluster sampling, 872\\
for panel data, 514-517, 614

Generalized extreme value distribution, 651\\
Generalized Gauss-Newton optimization method, 434-435\\
Generalized information matrix equality (GIME), 417\\
Generalized instrumental variables (GIV) estimator, 207, 222-226\\
comparison with generalized method of moment, 224-226\\
comparison with three-stage least squares (3SLS) estimator, 224-226, 254-255\\
derivation, 222-224\\
Generalized inverse, 312\\
Generalized least squares (GLS) analysis cluster sampling, 866, 878\\
fixed effects (FEGLS), 300, 312-315\\
random effects analysis, 292-294\\
systems of equations, 161, 173-184, 200-201\\
asymptotic normality, 175-176\\
consistency, 173-175\\
feasible GLS (FGLS), 176-188, 200-201\\
ordinary least squares (OLS) analysis versus, 185-188\\
weighted multivariate nonlinear least squares (WMNLS) estimator, 444-449\\
Generalized linear model (GLM), 512-513\\
binomial variance assumption, 739-740\\
Poisson, 725, 729-730\\
standard error, 729-730\\
Generalized method of moments (GMM) estimation, 207, 213-222, 226-229, 232-233, 525-555, 748\\
asymptotic properties, 525-530\\
efficient estimation, 538-545\\
general weighting matrix, 213-216\\
in cluster sampling, 871-872\\
in linear unobserved effects panel data model, 345-349, 370-371, 372-374\\
in nonlinear simultaneous equations (SEMs), 270, 272\\
minimum distance estimation classical, 545-547\\
unobserved effects model, 549-551\\
optimal weighting matrix, 217-219\\
panel data application, 547-555\\
minimum distance approach to unobserved effects model, 549-551\\
nonlinear dynamic models, 547-549\\
time-varying coefficients on unobserved effects, 551-555\\
systems of nonlinear equations, 532-538\\
three-stage least squares (3SLS) analysis, 219-222, 232\\
two-stage least squares (2SLS) estimator, 216-217\\
under orthogonality conditions, 530-532

Generalized method of moments (GMM) estimator\\
comparison with generalized instrumental variables (GIV) estimator, 224-226\\
statement of, 525-530\\
testing classical hypotheses, 226-227\\
testing overidentification restrictions, 228-229, 551\\
Generalized propensity score, 963\\
Generalized residual function, 530\\
Generated instruments, 125\\
estimating, 125-126\\
GMM with generated instruments, 543-545\\
two-stage least squares (2SLS) analysis, 124-125\\
Generated regressors, 123\\
estimating, 125-126\\
ordinary least squares analysis in single-equation linear models, 123-124\\
Geometric QMLE, 738\\
GLM variance assumption, 512-513, 725\\
GMM. See Generalized method of moments (GMM) estimation\\
GMM criterion function statistic, 529-530\\
GMM distance statistic, 227, 529-530\\
GMM three-stage least squares (GMM 3SLS) estimator, 219-222, 232\\
GMM with generated instruments, 543-545\\
Goodness-of-fit measure, 573-574, 677\\
Grouped duration data, 1010-1018\\
time-invariant covariates, 1011-1015\\
time-varying covariates, 1015-1017\\
unobserved heterogeneity, 1017-1018\\
Hausman test, 324-325, 328-334\\
comparing random and fixed effects estimators, 328-334, 355-356\\
computing Hausman statistic, 331-332\\
Hausman and Taylor (HT) model, 358-361\\
key component, 330\\
Heckit procedure, 699, 805-808\\
Heckman's method (Heckit), 699, 805-808, 838-849\\
Hedonic price system, 535-536\\
Hessian form of the LM statistic, 424-427\\
Hessian of the objective function, 406-409, 414, 417, 420\\
Heterogeneity\\
neglected, 582-585, 680-681\\
probit models with heterogeneity and endogenous explanatory variables, 630-632\\
unobserved, 22, 285, 1003-1010, 1017-1018\\
Heterokurtosis-robust test, 141\\
Heteroskedasticity\\
data censoring, 781\\
for two-stage least squares (2SLS) analysis, 106-107\\
heteroskedastic probit model, 602-605, 686-687\\
heteroskedasticity-robust $t$ statistics, 62, 106-107, 136-138\\
in latent variable model, 599-604, 606, 685-687\\
ordinary least squares (OLS) analysis, 60-62\\
specification tests, 138-141\\
standard error, 61-62\\
testing after pooled ordinary least squares (POLS) analysis, 199-200\\
Heteroskedasticity-robust variance matrix estimator for NLS, 416-417\\
Hierarchical linear models (HLM), in cluster sampling, 876-877, 879, 881\\
Hierarchical model, 649\\
Histogram estimator, 457\\
Homogenous linear restrictions, 247\\
Homokurtosis, 139-140\\
Homoskedasticity\\
assumption in sample selection, 796-797\\
asymptotic efficiency, 231\\
ordinary least squares (OLS) analysis, 59-60, 131, 220-221\\
system homoskedasticity assumption, 180-182\\
Huber standard error, 61\\
Huber-White sandwich estimator, 415-416, 446-447\\
Hurdle models\\
lognormal, 694-696\\
nature of, 691\\
truncated normal, 692-694, 695\\
Hypothesis testing\\
in binary response index models, 569-573\\
multiple exclusion restrictions, 570-571\\
nonlinear hypothesis about (BETA), 571\\
tests against more general alternatives, 571-573\\
in maximum likelihood estimation (MLE), 481-482\\
nonlinear regression model, 420-431\\
behavior of statistics under alternatives, 430-431\\
change in objective function, 428-430\\
score (Lagrange multiplier) tests, 421-428\\
Wald tests, 420-421\\
Poisson quasi-maximum likelihood estimator (QMLE), 732-734\\
two-stage least squares (2SLS) analysis, 104-106\\
Identification assumption, 13-14\\
Identification problem, 91-92\\
Identified due to a nonlinearity, 265-266\\
Idiosyncratic disturbances, 285\\
Idiosyncratic errors, 285

Ignorability of treatment, 908-936\\
identification, 911-915\\
matching methods, 934-936\\
propensity score methods, 920-934\\
regression adjustment, 915-920, 930-934\\
Ignorable selection, 822\\
Imperfect proxy variables, 69\\
Incidental parameters problem, 495, 612\\
Incidental truncation, 777-778, 802-821\\
exogenous explanatory variables, 802-808\\
Tobit selection equation, 815-821\\
Independence from irrelevant alternatives (IIA) assumption, 648-649\\
Independent, not identically distributed (i.n.i.d.) sample, 6, 146-147\\
Independent identically distributed (i.i.d.) sample, 5\\
Independent variable. See Explanatory variable ( $x$ )\\
Index models, 565-582\\
binary response models\\
logit, 566, 568, 573-582\\
maximum likelihood estimation (MLE), 567-569\\
probit, 566, 568, 573-582, 587, 591, 594-599\\
reporting results, 573-582\\
testing, 569-573\\
maximum likelihood estimation (MLE) for binary response, 567-569\\
Index structure, 512\\
Indicator function, 450\\
Individual effects, 285\\
Individual heterogeneity, 285\\
Individual-specific slopes, 374-387\\
general models, 377-381\\
random trend model, 375-377\\
Influence function representation, 406-407\\
Influential observations, 451-453\\
Initial condition, 497-498\\
Initial conditions problems, 626\\
Instrumental variables (IV) estimation, 142-143, 207-233\\
average treatment effect (ATE), 937-954\\
continuous endogenous explanatory variables, 591-594\\
first difference instrumental variables (FDIV) estimator, 361-365\\
fixed effects (FE) methods, 353-358\\
instrumental variable characteristics, 90\\
nonlinear instrumental variables estimator, 531\\
random effects (RE) methods, 349-353\\
single-equation linear models, 89-114\\
examples, 93-96\\
motivation for, 89-98\\
omitted variables problem, 92-93, 112-114\\
two-stage least squares (2SLS) analysis. See\\
Two-stage least squares (2SLS) analysis\\
systems of equations, 207-233\\
examples, 207-210\\
general linear system of equations, 210-213\\
generalized instrumental variables (GIV) estimator, 207, 222-226\\
generalized linear system of equations, 210-213\\
generalized method of moments (GMM) estimation, 207, 213-222, 226-229, 232-233, 522-525, 748\\
introduction, 207\\
optimal instruments, 229-232\\
simultaneous equations model (SEM), 207\\
three-stage least squares (3SLS) analysis, 219-222, 232\\
two-stage least squares (2SLS) analysis, 232-233\\
Instrumental variables (IV) estimator\\
consistency of, 142-143\\
of (BETA), 92\\
Internal covariates, 991\\
Interval coding, 783-785\\
Interval regression, 783-785\\
Inverse Mills ratio, 672-673, 838\\
Inverse probability weighting (IPW), 778, 821-827, 840-844\\
Iteration. See Law of iterated expectations (LIE)\\
Jensen's inequality, 31-32\\
Just identified equations, 251\\
Kaplan-Meier estimator, 1014-1015\\
Kernel density estimator, 457\\
Kernel estimators, 915\\
Knowledge of the world of work (KWW) test, 72, 812-813\\
Kullback-Leibler information criterion (KLIC), 473-474, 503-505, 523-524

Lagged dependent variables, 290, 371-374, 497-499\\
Lagrange multiplier (score) statistic, 62-65, 107, 299, 421-428\\
Large-sample theory. See Asymptotic analysis\\
Latent class model, 651\\
Latent variable, 285, 471-472\\
model of, 565\\
heteroskedasticity, 599-604, 606, 685-687\\
nonnormality, 599-604\\
Law of iterated expectations (LIE), 18-22, 30-32, 34-36, 288, 414, 672\\
general statement, 19, 21\\
identification problem, 20\\
proxy variables, 23-24, 67-72\\
two-stage least squares (2SLS) analysis, 101\\
Law of large numbers\\
statement of, 61, 401\\
weak law of large numbers (WLLN), 42, 403-404, 526\\
Least absolute deviations (LAD) estimator, 404, 451-453, 871\\
Least squares. See Multivariate nonlinear regression methods; Ordinary least squares (OLS) analysis; Three-stage least squares (3SLS) analysis; Two-stage least squares (2SLS) analysis\\
Least squares linear predictor, 26-27\\
Left censoring, 1000\\
Left truncation, 1000\\
Length-biased sampling, 1000\\
Likelihood ratio (LR) statistic, 481-482, 677\\
Limited dependent variable, 559\\
Limited information maximum likelihood (LIML) estimator, 745-746\\
Limited information procedure, 591-592\\
Limited range responses, average treatment effect (ATE), 960-961\\
Limiting behavior, of estimators, 42-45\\
Linear exponential family (LEF), 509-514\\
Linear panel data models attrition, 837-845\\
linear unobserved effects. See Linear unobserved effects panel data models\\
sample selection, 827-837\\
fixed effects estimation with unbalanced panels, 828-831\\
random effects estimation with unbalanced panels, 831-832\\
testing and correcting for selection bias, 832-837\\
systems of equations, 163-167, 191-201\\
assumptions for pooled ordinary least squares, 191-194\\
contemporaneous exogeneity, 164, 165\\
dynamic completeness, 194-196\\
example, 163-166, 167-168, 169-171, 191-192\\
feasible generalized least squares under strict exogeneity, 200-201\\
finite distributed lag (FDL) model, 165-166\\
pooled ordinary least squares, 191-194, 198-199\\
robust asymptotic variance matrix, 197-198\\
sequential exogeneity, 164-165\\
strict exogeneity, 165-166\\
testing for heteroskedasticity, 199-200\\
testing for serial correlation, 198-199\\
time series persistence, 196-197\\
Linear probability model (LPM), for binary response, 562-565

Linear projections, 25-27\\
least squares linear predictor, 26-27\\
minimum mean square linear predictor, 26-27\\
partial linear model, 29\\
properties, 34-36\\
Linear simultaneous equations models (SEMs), 241-261\\
cross equation restrictions, 256-257\\
efficiency, 260-261\\
estimation after identification, 252-256\\
exclusion restrictions, 241-243\\
general linear restrictions, 245-248\\
identification, 251, 256-261\\
order condition, 245, 249, 250-251\\
rank condition, 248-251\\
reduced forms, 243-245, 255-256\\
structural equations, 241-242, 245-251, 260-261\\
Linear unobserved effects panel data models, 281-387\\
assumptions, 285-290\\
random versus fixed effects, 285-287\\
strict exogeneity, 287-288\\
comparison of estimators, 321-334\\
fixed effects versus first differencing, 321-326\\
Hausman test, 328-334\\
random effects and fixed effects estimators, 285-287, 326-334\\
correlated random slopes, 384-387\\
estimating by pooled ordinary least squares, 291\\
estimation under sequential exogeneity, 368-374\\
examples, 289-290\\
first differencing methods, 315-321\\
inference, 315-318\\
instrumental variables, 361-365\\
policy analysis, 320-321\\
robust variance matrix, 318-319\\
testing for serial correlation, 319-320\\
fixed effects methods, $300-315$\\
asymptotic inference with fixed effects, 304-307\\
consistency of fixed effects estimator, 300-304\\
dummy variable regression, 307-310\\
fixed effects versus first differencing, 321-326\\
instrumental variables, 353-358\\
random effects and fixed effects estimators, 285-287, 326-334\\
robust variance matrix estimator, 310-315\\
robustness of standard, 382-384\\
serial correlation, 310-315\\
generalized method of moments (GMM) approaches, 345-349, 370-374\\
Hausman and Taylor (HT) models, 358-361\\
models with individual-specific slopes, 374-387\\
omitted variables problem, 281-285\\
random effects models, 285-287, 291-300

Linear unobserved effects panel data models (cont.)\\
estimation and inference, 291-297\\
general feasible generalized least squares analysis, 298-299\\
instrumental variables methods, 348-353\\
robust variance matrix estimator, 297-298\\
testing for presence of unobserved effect, 299-300\\
random trend model, 375-377\\
unobserved effects, 282\\
unobserved effects models with measurement error, 365-368\\
Link function, 512\\
LM statistic\\
behavior under alternatives, 430-431\\
expected Hessian form of the, 425\\
Hessian form of the, 424-427\\
outer product of the score, 424\\
Local alternatives, 45-46, 430-431\\
Local average response (LAR), 634-635\\
Local average treatment effect (LATE), 906\\
Local linear regression, 956\\
Local power analysis, 45-46\\
Log-logistic hazard function, 987-988\\
Log-odds transformation, 749\\
Logistic distribution, 988\\
Logistic function, 397\\
Logit model, 566\\
conditional logit (CL) model, 647-649\\
fractional logit regression, 751\\
logit estimator, 568, 621-624\\
mixed logit model, 648\\
multinomial logit (MNL) model, 643-648\\
nested logit model, 649-651\\
ordered, 656-658\\
pooled, 609-610\\
pooled multinomial logit, 654\\
unobserved effects under strict exogeneity, 619-625\\
Lognormal hurdle (LH) model, 694-696\\
Longitudinal data. See Panel (longitudinal) data\\
M-estimators, 400-401, 403-405\\
asymptotic normality, 407-409, 411-413\\
in cluster sampling, 871\\
in complex survey sampling, 897-899\\
sample selection, 841, 843-845\\
two-step, 409-413\\
adjustments, 418-420\\
unweighted M-estimator in stratified sampling, 861-863\\
weighted M-estimator in stratified sampling, 856-861\\
Mahalanobis distance, 934

Matching on the covariates, 934\\
Matrix of instruments, 530\\
Maximum likelihood estimation (MLE), 469-517\\
asymptotic normality, 476-479\\
asymptotic variance, 479-481\\
binary response index model, 567-569\\
conditional maximum likelihood estimation (CMLE), 470-471, 473-476\\
continuous exogenous explanatory variable, 591-594\\
count variable, 472\\
data censoring, 780-781\\
efficiency in generalized method of moments (GMM) estimation, 540-542\\
hypothesis testing, 481-482\\
in cluster sampling, 892-893\\
index models, 567-569\\
panel data models with unobserved effects, 494-499\\
lagged dependent variables, 497-499\\
strictly exogenous explanatory variables, 494-497\\
parametric model, 471\\
partial (pooled) likelihood methods for panel data, 485-494\\
asymptotic inference, 490-492\\
inference with dynamically complete models, 492-494\\
setup, 486-490\\
Poisson regression, 472-473, 475, 477, 478, 481, 485\\
probit model, 471-472, 475, 477, 478, 480-481, 488-489, 492\\
quasi-maximum likelihood estimation, 502-517\\
general misspecification, 502-504\\
generalized estimating equations for panel data, 514-517\\
in linear exponential family, 509-514\\
model selection tests, 505-509\\
specification testing, 482-485\\
two-step estimators involving maximum likelihood, 499-502\\
first-step estimator, 500-502\\
second-step estimator as maximum likelihood estimator, 499\\
Maximum score estimator, 605-606\\
Mean independent, average treatment effect (ATE), 907\\
Measurement error\\
endogenous variable, 55\\
explanatory variable, 78-82\\
in ordinary least squares (OLS) analysis, 76-82 attenuation bias, 81\\
classical errors-in-variables (CEV) assumption, 80-82, 367\\
dependent variable error, 76-78\\
examples, 78, 81\\
independent variable error, 78-82\\
multiplicative measurement error, 78\\
proxy variable versus, 76, 92-93, 113\\
unobserved effects models, 365-368\\
Measures, 522\\
Median regression, 404-405\\
Method of moment, ordinary least squares (OLS) analysis, 57-58\\
Mills ratio, 809-812\\
Minimum chi-square estimator, 528, 545\\
Minimum distance (MD) estimates, 545-547, 889-894\\
Minimum mean square linear predictor, 26-27\\
Missing at random (MAR), 795\\
Missing completely at random (MCAR), 794-795, 827\\
Mixed models\\
in cluster sampling, 876-877\\
mixed logit model, 648\\
Mixture model, 651, 1003\\
Model selection tests, 505-509\\
Monte Carlo simulation, 436-438\\
Multinomial logit (MNL) model, 643-648\\
Multinomial probit model, 649\\
Multinomial response models, 643-654\\
endogenous explanatory variables, 651-653\\
multinomial logit (MNL) model, 643-648\\
ordered response models, 643, 655-663\\
endogenous explanatory variables, 660-662\\
ordered logit, 656-658\\
ordered probit, 655-658\\
panel data methods, 662-663\\
specification issues, 658-659\\
panel data methods, 653-654\\
probabilistic choice models, 646-651\\
Multinomial sampling, 855\\
Multiple exclusion restrictions, 570-571\\
Multiple indicator solution, 112-114\\
Multiple-spell data, 1018-1019\\
Multiple treatments, average treatment effect (ATE), 964-967\\
Multiplicative measurement error, 78\\
Multiplicative random effects model, 759-760\\
Multistage sampling, 894-895\\
Multivalued treatments, average treatment effect (ATE), 961-963\\
Multivariate nonlinear least squares estimator, 443\\
Multivariate nonlinear regression methods, 442-449\\
multivariate nonlinear least squares, 442-444\\
weighted multivariate nonlinear least squares, 444-449\\
Mundlak representation, 461, 553-554\\
N-consistent estimator, 44-45\\
N -equivalent estimator, 44-45\\
N-R-squared test, 63\\
Natural experiments, 94-95, 147\\
Negative duration dependence, 987\\
Neglected heterogeneity, 582-585\\
Tobit models, 680-681\\
Nested logit model, 649-651\\
Newey-Tauchen-White (NTW) statistic, 484-485\\
Newton-Raphson optimization method, 432-433\\
NLS estimator, 400\\
Nominal response, 643\\
Nonlinear endogenous variables, 262-263\\
Nonlinear estimation\\
generalized method of moments (GMM) approach. See Generalized method of moments (GMM) estimation\\
maximum likelihood methods. See Maximum likelihood estimation (MLE)\\
nonlinear regression model. See Nonlinear regression model\\
Nonlinear hypothesis testing, about (BETA), 571\\
Nonlinear least squares (NLS) assumption, 399-400, 408-409, 669\\
multivariate, 442-444\\
variable addition test (VAT) approach, 427-428, 573\\
weighted nonlinear least squares analysis, 409-413\\
Nonlinear least squares residuals, 416\\
Nonlinear panel data models, generalized method of moments (GMM) estimation, 547-555\\
minimum distance approach to unobserved effects model, 549-551\\
models with time-varying coefficients on unobserved effects, 551-555\\
nonlinear dynamic models, 547-549\\
Nonlinear regression model, 397-462\\
asymptotic inference, 454-459\\
asymptotic normality, 405-409, 411-413\\
asymptotic variance, 413-420\\
adjustments for two-step estimation, 418-420\\
estimation without nuisance parameters, 413-418\\
bootstrapping, 438-442\\
consistency, 410-411\\
correctly specified model for the conditional mean, 397-398\\
hypothesis testing, 420-431

Nonlinear regression model (cont.)\\
behavior of statistics under alternatives, 430-431\\
change in objective function, 428-430\\
score (Lagrange multiplier) tests, 421-428\\
Wald tests, 420-421\\
identification, 401-402\\
M-estimator, 400-401, 403-405, 409-413\\
adjustments for two-step estimation, 418-420\\
two-step, 409-413\\
median regression, 404-405\\
Monte Carlo simulation, 436-438\\
multivariate, 442-449\\
nonlinear least squares, 442-444\\
weighted multivariate nonlinear least squares, 444-449\\
nonlinear least squares (NLS) assumption, 399-400, 408-413\\
optimization methods, 431-436\\
Berndt, Hall, Hall, and Hausman algorithm, 433-434\\
concentrating parameters out of objective function, 435-436\\
generalized Gauss-Newton method, 434-435\\
Newton-Raphson method, 432-433\\
parametric model, 397\\
quantile estimation, 449-462\\
consistency, 453\\
estimation problem, 449-453\\
quantile regression panel data, 459-462\\
quantiles, 449-453\\
resampling methods, 438-442\\
uniform convergence in probability, 402-403\\
Nonlinear simultaneous equations models (SEMs), 262-271\\
control function estimation for triangular systems, 268-271\\
different instruments for different equations, 271-273\\
estimation, 266-271\\
forbidden regression, 267-268\\
three-stage least squares (3SLS) analysis, 266-267\\
two-stage least squares (2SLS) analysis, 266-267\\
identification, 262-266\\
Nonlinear 2SLS (N2SLS) estimator, 534-535\\
Nonlinear SUR estimator, 447\\
Nonnormality\\
data censoring, 781\\
in latent variable model, 599-604\\
Nonparametric bootstrap, 438\\
Nonparametric density estimator, 456-457\\
Nonparametric regression, 605\\
Nonparametric residual bootstrap, 441\\
Nonsymmetric test, 440

Normal distribution, asymptotically normal, 40-41, 43-44, 175-176, 405-409, 411-413, 728-732\\
Normalization restriction, 247\\
Normalized differences, 917\\
Objective function\\
change in, 428-430\\
concentrating parameters out of, 435-436\\
Hessian of the, 406-409, 414, 417, 420\\
OLS. See Ordinary least squares (OLS) analysis\\
OLS equation by equation, 169\\
Omitted variables\\
linear unobserved effects panel data models, 281-285\\
ordinary least squares (OLS)\\
bias, 66-67\\
ignoring omitted variables, 65-67\\
inconsistency, 66-67\\
nature of, 54-55\\
proxy variable, 67-72\\
solutions to, 65-76\\
two-stage least squares (2SLS), 112\\
Optimization methods, nonlinear regression model, 431-436\\
Berndt, Hall, Hall, and Hausman algorithm, 433-434\\
concentrating parameters out of objective function, 435-436\\
generalized Gauss-Newton method, 434-435\\
Newton-Raphson method, 432-433\\
Order condition, 99, 211\\
in simultaneous equations models (SEMs), 245, 249-251\\
Ordered response models, 643, 655-663\\
ordered logit model, 656-657\\
ordered probit model, 655-657\\
parallel regression function, 658-659\\
specification issues, 658-659\\
Ordinary least squares (OLS) analysis, 6, 10. See also Pooled ordinary least squares (POLS) analysis\\
incidental truncation in sample selection, 806-808\\
influential observations, 451-453\\
linear models estimation by, 796-798\\
linear unobserved effects panel data models, 283-284\\
OLS estimator, 127-128\\
pooled OLS estimation, 865-866\\
simultaneous equations, 253-255\\
single-equation linear models, 53-82\\
asymptotic properties, 54-65\\
estimable model, 53\\
examples, 63-65, 67, 69-72, 75\\
generated regressors, 123-124\\
omitted variable problem, 65-76\\
population model, 53\\
properties under measurement error, 55, 76-82\\
structural model, 53, 66\\
weighted least squares (WLS) in, 60-61\\
zero-mean assumption, 53-54\\
specification tests\\
endogeneity, 129-134\\
functional form, 137-138\\
heteroskedasticity, 138-141\\
overidentifying restrictions, 134-137\\
system estimation methods, 269-270\\
systems of equations, 161, 166-173, 198-200\\
asymptotic properties of system OLS, 167-172\\
feasible generalized least squares versus, 185-188\\
pooled ordinary least squares (POLS) estimator, 169-172, 191-194, 198-199\\
preliminaries, 166-167\\
testing multiple hypotheses, 172-173\\
Tobit I model, 677-680\\
two-stage least squares (2SLS) compared with, 107-112\\
Orthogonality condition\\
general method of moments (GMM) estimation, 530-532\\
population, 56-57\\
Outer product of the score LM statistic, 424\\
Overdispersion, 512-513\\
Overidentification exogenous variable, 98\\
Overidentification test statistic, 228-229, 551\\
Overidentified equations, 251\\
Overidentifying restrictions, 98, 251\\
specification tests, 134-137\\
Overlap assumption, 910\\
Panel (longitudinal) data, 6-7. See also Linear panel data models; Nonlinear panel data models\\
average treatment effect (ATE), 968-975\\
binary response model, 608-635\\
dynamic unobserved effects models, 625-630\\
pooled probit and logit, 609-610\\
probit models with heterogeneity and endogenous explanatory variables, 630-632\\
semiparametric approaches, 605, 632-635\\
unobserved effects logit models under strict exogeneity, 619-625\\
unobserved effects probit models under strict exogeneity, 610-619\\
cluster sampling with unit-specific panel data, 876-883\\
corner solution responses, 705-715\\
dynamic unobserved effects Tobit models, 713-715\\
pooled methods, 705-707\\
unobserved effects models under strict exogeneity, 707-713\\
count data, 755-769\\
conditional expectations with unobserved effects, 758-759\\
fixed effects Poisson estimation, 762-764\\
fractional response model, 766-769\\
pooled QMLE, 756-758\\
random effects methods, 759-762\\
relaxing strict exogeneity assumption, 764-766\\
fractional response models for panel data, 766-769\\
models with lagged dependent variables, 497-499\\
models with unobserved effects, 494-499\\
strictly exogenous explanatory variables, 494-497\\
multinomial response models, 653-654\\
ordered response models, 662-663\\
partial (pooled) likelihood methods, 485-494\\
Parallel regression assumption, 658-659\\
Parameter space, 397\\
Parametric bootstrap, 440-441\\
Parametric model, 14, 397, 471\\
Partial effect at the average (PEA), 575-577\\
Partial effects\\
average partial effects (APE), 22-25, 73, 141-142\\
binary response models, 577-582, 586-589, 615\\
conditional expectations, 15-16, 22-25\\
estimating, 3-4\\
nature of, 3\\
structural, 583-585\\
Partial linear model, 29\\
Partial log likelihood, 486\\
Partial (pooled) likelihood methods for panel data, 485-494\\
asymptotic inference, 490-492\\
inference with dynamically complete models, 492-494\\
setup, 486-490\\
Partial (pooled) maximum likelihood estimator (PMLE), 487-489\\
Partial QMLE, 514\\
Participation decision, 690-691\\
Partitioned projection formula, 36\\
Pearson dispersion estimator, 513\\
Pearson residuals, 513\\
Percentile- $t$ method, 440\\
Piecewise-constant proportional hazard, 1013-1014\\
Point-wise convergence in probability, 402\\
Poisson GLM variance assumption, 725\\
Poisson quasi-maximum likelihood estimator (QMLE), 727-732, 741, 745

Poisson random effects model, 760\\
Poisson regression, 724-736\\
assumptions, 724-727\\
asymptotic normality, 728-732\\
consistency, 727-728\\
fixed effects model, 756, 762-764\\
hypothesis testing, 732-734\\
maximum likelihood estimation (MLE), 472-473, 475, 477, 478, 481, 485\\
quantities of interest, 724-727\\
random effects model, 760\\
specification testing, 734-736\\
Poisson variance assumption, 725\\
Policy analysis\\
difference-in-differences estimation, 147-151\\
first differencing, 320-321\\
fixed effects estimation, 315\\
Pooled bivariate probit, 631-632\\
Pooled cross section, 5-6\\
over time, 146-147\\
Pooled IV probit, 631-632\\
Pooled least absolute deviations, 461-462\\
Pooled log likelihood, 486\\
Pooled logit, 609-610\\
Pooled multinomial logit, 654\\
Pooled negative binomial analysis, 761\\
Pooled nonlinear least squares (PNLS) estimator, 443-444\\
Pooled ordered probit, 662-663\\
Pooled ordinary least squares (POLS) analysis, 198-199. See also Ordinary least squares (OLS) analysis\\
assumptions, 191-194\\
in cluster sampling, 867-870, 877-878, 882-884\\
OLS test statistic, 332-333\\
pooled ordinary least squares (POLS) estimator, 169-172, 191-194, 198-199\\
testing for heteroskedasticity, 199-200\\
testing for serial correlation, 198-199\\
unobserved effects models, 291, 365-366\\
Pooled Poisson QMLE, 756-758\\
Pooled probit, 609-610, 662-663\\
Pooled probit estimator, 488-489\\
Pooled QMLE, 514\\
Pooled quantile regression, 459-460\\
Pooled weighted nonlinear least squares (PWNLS) estimator, 447\\
Poorly identified model, 265-266, 402\\
Population-averaged (PA) model, 614-615\\
Population-averaged parameters, 612-613\\
Population model, 5\\
ordinary least squares (OLS) estimation, 53\\
Population orthogonality condition, 56-57\\
Positive duration dependence, 987

Prais-Winston estimator, 201\\
Primary sampling unit (PSU), 894\\
Probabilistic choice models, 646-651\\
Probability\\
boundedness in, 37-40\\
convergence in asymptotic analysis, 38\\
with probability approaching one (w.p.a.1.), 39-40\\
Probability limit (plim), 38\\
Probit model, 471-472, 475, 477, 478, 480-481, 488-489, 492, 566\\
attrition in linear panel data models, 838\\
bivariate, 595-599\\
conditional, 649\\
fractional probit regression, 751, 753\\
heteroskedastic, 602-605, 686-687\\
incidental truncation in sample selection, 802-814\\
index models, 566, 568, 573-582, 587, 591, 594-599\\
instrumental variables probit (IV probit), 591\\
lognormal hurdle model, 696\\
multinomial, 649\\
ordered, 655-658\\
pooled, 609-610, 662-663\\
probit estimator, 568\\
random effects ordered probit, 663\\
reporting results for, 573-582\\
unobserved effects under strict exogeneity, 610-619\\
with heterogeneity and endogenous explanatory variables, 630-632\\
Propensity scores, 911, 920-934\\
matching of, 936\\
Proportional hazard models, 988\\
Proportional hazard with time-varying covariates, 991\\
Proxy variables\\
examples, 69-72\\
formal requirements, 67-69\\
imperfect, 69\\
law of iterated expectations (LIE), 23-24, 67-72\\
measurement error versus, $76,92-93,113$\\
ordinary least squares (OLS) analysis, 67-72\\
Pseudo-maximum likelihood estimator, 503\\
Pseudo-R-squared measures, 574-575, 580\\
QLR (quasi-likelihood ratio) statistic, 429-430, 532\\
QMLE. See Quasi-maximum likelihood estimation (QMLE)\\
Quantile estimation, 449-462\\
Quantile regression, 404-405\\
for panel data, 459-462\\
quantile regression estimator, 451

Quasi-likelihood ratio (QLR) statistic, 429-430, 532\\
Quasi-log-likelihood (pseudo-log-likelihood) function, 503\\
Quasi-maximum likelihood estimation (QMLE) in linear exponential family, 509-514\\
maximum likelihood estimation (MLE), 502-517\\
general misspecification, 502-504\\
generalized estimating equations for panel data, 514-517\\
in linear exponential family, 509-514\\
model selection tests, 505-509\\
partial, 514\\
Poisson, 727-732, 745\\
pooled, 514, 756-758\\
random effects analysis, 760-761\\
Quasi-time demeaning, 327, 328\\
$R$-absolute loss function, 450\\
Random coefficient models, 73-76\\
correlated, 141-146\\
described, 74-75\\
example, 75\\
Random effects (RE)\\
analysis, 291-300, 495, 759-762\\
estimation and inference, 291-297\\
estimators in, 287, 294-297, 326-328, 552-554, 866-868\\
fixed effects analysis versus, 326-328, 349-358\\
framework of, 286-287\\
general feasible generalized least squares (FGLS) analysis, 298-299\\
Hausman test comparing random and fixed effects estimators, 328-334, 355-358\\
in cluster sampling, 867-871, 878-879, 882\\
multiplicative random effects model, 759-760\\
Poisson model, 760\\
quasi-MLE random effects analysis, 760-761\\
robust variance matrix estimator, 297-298\\
testing for presence of unobserved effect, 299-300, 552-554\\
weighted multivariate nonlinear least squares estimator (WMNLS), 761-762\\
logit model, 619-620\\
probit estimator, 613-614\\
probit model, 616-619, 663\\
structure of, 294\\
Tobit model, 709, 711-713\\
2SLS estimator, 352-353\\
Random effects instrumental variables (REIV) estimator, 349-353\\
Random growth model, 375-377\\
Random sampling, 5\\
limit theorems, 41-42\\
random sequence, 38-39\\
Random trend model, 375-377\\
Rank condition for identification, 92, 99-100, 210-211, 248-251\\
Rao's score principle, 422-423\\
Recursive system, in simultaneous equations models (SEMs), 259-260\\
Reduced-form equation\\
in simultaneous equations models (SEMs), 243-245, 255-256\\
nature of, 90-91, 243\\
Redundancy, of proxy variables, 67-68\\
Regress and (y). See Dependent variable (y)\\
Regression. See Nonlinear regression model; Poisson regression; Seemingly unrelated regressions (SUR) model\\
Regression adjustment, 915\\
average treatment effect (ATE), 915-920\\
ignorability of treatment, 915-920, 930-934\\
Regression discontinuity designs, 954-959\\
fuzzy regression discontinuity (FRD) design, 957-959\\
sharp regression discontinuity (SRD), 954-957\\
unconfoundedness, 959\\
Regressor (x). See Explanatory variable ( $x$ )\\
Regularity conditions, 4-5\\
Resampling methods, bootstrapping, 438-442\\
RESET, 137-138, 482, 734\\
Response probability, 561\\
Response surface analysis, 437-438\\
Response variable (y). See Dependent variable (y)\\
Restricted model estimator, 62-63\\
Right censoring, 785-786, 993\\
Risk set, 1014\\
Robust variance matrix estimator, 172, 297-298\\
first-differencing methods, 318-319\\
fixed effects serial correlation, 310-315\\
Rubin causal model (RCM), 903\\
Sample selection, 6, 790-845\\
attrition, 837-845\\
exponential response function, 814\\
general attrition, 837-845\\
incidental truncation, 802-814\\
binary response model, 813-814\\
endogenous explanatory variables, 809-813, 817-819\\
exogenous explanatory variables, 802-808, 815-817\\
Tobin selection equation, 815-821\\
inverse probability weighting (IPW) for missing data, 821-827, 840-844\\
linear models, 792-798\\
linear panel data models, 827-845

Sample selection (cont.)\\
fixed effects estimation with unbalanced panels, 828-831\\
random effects estimation with unbalanced panels, 831-832\\
sample selection bias, 832-837\\
nonlinear models, 798-799\\
structural Tobit model, 819-821\\
truncated regression, 799-802\\
Sampling. See Cluster sampling; Sample selection; Stratified sampling\\
Sampling weights, 857\\
Sargan-Hausman test, 135\\
Score of the log likelihood, 476-478\\
Score statistic, 62-65, 107, 299, 421-428\\
Second-stage regression, two-stage least squares (2SLS) analysis, 97, 105\\
Secondary sampling units, 894\\
Seemingly unrelated regressions (SUR) model, systems of equations, 7, 185-191\\
example, 161-163, 167-169, 186-189\\
ordinary least squares versus feasible generalized least squares, $185-188$\\
singular variance matrices, 189-191\\
systems with cross equation restrictions, 188-189\\
Selected sample, 790\\
Selection indicator, 793\\
Selection mechanisms, 790\\
Selection model, 692\\
Selection on observables, 795-796, 909\\
Self-selection problem, 289-290, 907-908\\
Semielasticity, conditional expectations, 18\\
Semiparametric estimators, 605, 632-634\\
Semiparametric method, 688-689\\
Semirobust variance matrix estimator, 415-416\\
Sequential exogeneity\\
estimation in unobserved effects model, 368-374\\
systems of equations, 164-165\\
Sequential moment restrictions, 368-371, 765-766\\
Sequentially exogenous conditional on the unobserved effect, 368-371\\
Sequentially exogenous covariates, 990\\
Serial correlation, 172\\
first-differencing methods, 319-320\\
inference with fixed effects, 305-306\\
robust variance matrix estimator, 310-315\\
testing after pooled ordinary least squares (POLS) analysis, 198-199\\
Series estimators, 915\\
Sharp regression discontinuity (SRD) design, 954-957\\
Simultaneity, endogenous variable, 55\\
Simultaneous equations models (SEMs), 207, 239-273\\
autonomy, 239, 247\\
causality, 239-240\\
examples, 240-241\\
linear equations, 241-261\\
covariance restrictions, 257-260\\
cross equation restrictions, 256-257\\
efficiency, 260-261\\
estimation after identification, 252-256\\
exclusion restrictions, 241-243\\
general linear restrictions, 245-248\\
identification, 251, 256-261\\
order condition, 245, 249-251\\
rank condition, 248-251\\
reduced forms, 243-245, 255-256\\
structural equations, 241-242, 245-251\\
nonlinear equations, 262-271\\
control function estimation for triangular systems, 268-271\\
different instruments for different equations, 271-273\\
estimation, 266-268\\
identification, 262-266\\
scope, 239-241\\
structural equations, 241-242, 245-251, 260-261\\
Single-equation linear models\\
control function approach to endogeneity, 126-129, 145-146\\
correlated random coefficient models, 141-146\\
difference-to-difference estimation, 147-151\\
estimation with generated regressors and instruments, 123-126\\
instrumental variables estimation, 89-114\\
measurement error problem, 112-114\\
motivation for, 89-98\\
omitted variables problem, 112\\
two-stage least squares (2SLS), 96-112\\
ordinary least squares (OLS) analysis. See\\
Ordinary least squares (OLS) analysis, Singleequation linear models\\
overview, 53-54\\
pooled cross sections over time, 146-147\\
specification tests, 129-141\\
endogeneity, 129-134\\
functional form, 137-138\\
heteroskedasticity, 138-141\\
overidentifying restrictions, 134-137\\
two-stage least squares (2SLS) analysis, 96-112\\
asymptotic efficiency, 103-104\\
asymptotic normality, 101-102\\
consistency, 98-100, 108\\
heteroskedasticity-robust inference, 106-107\\
hypothesis testing, 104-106\\
potential pitfalls, 107-112\\
Single-spell data, 991-1010

Slutsky's theorem, 39, 47\\
Smearing estimate, 696\\
Smith-Blundell procedure, 684-685\\
Spatial correlation, 6\\
Specification tests, 129-141\\
endogeneity, 129-134, 594-599\\
functional form, 137-138\\
heteroskedasticity, 138-141\\
in maximum likelihood estimation (MLE), 482-483\\
in ordered models, 658-659\\
nonlinear regression model, 421-422\\
overidentifying restrictions, 134-137\\
Poisson quasi-maximum likelihood estimator (QMLE), 734-736\\
Spillover effect, 8-9\\
Stable unit treatment value assumption (SUTVA), 905\\
Standard censored regression model. See Type I Tobit model\\
Standard error\\
asymptotic, 44\\
bootstrap, 438-439, 581, 590\\
GLM standard error, 729-730\\
heteroskedasticity-robust standard error, 61\\
Huber, 61\\
two-stage least squares (2SLS) analysis, 108-112\\
White, 61\\
Standard stratified sampling (SS sampling), 854-856, 860\\
Standardized residual, 418-419, 570-571\\
State dependence, 371-374, 626\\
Static models, 163-164\\
Stochastic analysis, 4-11\\
asymptotic analysis, 7\\
data structures, 4-7\\
cluster sampling, 6\\
cross section data, 5-6\\
experimental data, 5\\
panel (longitudinal) data, 6-7\\
random sampling assumption, 5\\
spatial correlation, 6\\
examples, 7-9\\
setting selection, 4-5\\
Stock sampling, $1000-1003$\\
Stratified sampling, 6, 853-863\\
stratification based on exogenous variables, 861-863\\
variable probability sampling, 854-856\\
weighted estimators to account for stratification, 856-861\\
Strict exogeneity, 165, 325-326, 329\\
strictly exogenous conditional on the unobserved effect, 287-288, 495, 610-619\\
strictly exogenous corner solution responses, 707-713\\
strictly exogenous covariates, 990\\
Strong ignorability, 911\\
Structural conditional expectation, 13\\
Structural equations\\
estimating Tobit equations with sample selection, 819-821\\
simultaneous equations models (SEMs), 241-242, 245-251, 260-261\\
Structural error, ordinary least squares (OLS) analysis, 66-67\\
Structural model, ordinary least squares (OLS) analysis, 53\\
Structural partial effects, 583-585\\
Subject-specific (SS) model, 614-615\\
Survival analysis, 983\\
Survivor function, 985\\
Symmetric test, 440\\
System homoskedasticity assumption, 180-182\\
System instrumental variables (SIV), 207\\
System ordinary least squares (SOLS), 167-172\\
OLS equation by equation, 169\\
pooled ordinary least squares (POLS) estimator, 169-172, 191-194, 198-199\\
system ordinary least squares estimator of (BETA), 168-169\\
System ordinary least squares (SOLS) estimator of (BETA), 168-169\\
Systems of equations, 161-233\\
examples, 161-166\\
generalized least squares (GLS) analysis, 161, 173-184, 200-201\\
asymptotic normality, 175-176\\
consistency, 173-175\\
feasible GLS (FGLS), 176-188, 200-201\\
instrumental variables estimation, 207-233\\
examples, 207-210\\
general linear system of equations, 210-213\\
generalized instrumental variables (GIV) estimator, 207, 222-226\\
generalized method of moments (GMM) estimation, 207, 213-222, 226-229, 232-233\\
introduction, 207\\
optimal instruments, 229-232\\
simultaneous equations model (SEM), 207\\
three-stage least squares (3SLS) analysis, 219-222, 232\\
two-stage least squares (2SLS) analysis, 232-233\\
introduction, 161\\
linear panel data model, 163-167, 191-201\\
assumptions for pooled ordinary least squares, 191-194

Systems of equations (cont.)\\
contemporaneous exogeneity, 164-165\\
dynamic completeness, 194-196\\
example, 163-171\\
sequential exogeneity, 164-165\\
strict exogeneity, 165-166\\
time series persistence, 196-197\\
nonlinear equations, 532-538\\
ordinary least squares (OLS) analysis, 161, 166-173, 198-200\\
preliminaries, 166-167\\
testing multiple hypotheses, 172-173\\
seemingly unrelated regressions (SUR) model, 185-191\\
example, 161-163, 167, 169, 186-188\\
ordinary least squares versus generalized least squares analysis, 185-188\\
singular variance matrices, 189-191\\
systems with cross equation restrictions, 188-189\\
$T$ statistics\\
bootstrapping critical values for, 439-440\\
heteroskedasticity-robust $t$ statistics, 62\\
Test statistics, asymptotic properties of, 45-47\\
Three-stage least squares (3SLS) analysis, 219-222, 232\\
comparison with generalized instrumental variables (GIV) estimator, 224-226, 254-255\\
comparison with generalized method of moments (GMM) analysis, 345-347, 372-374\\
estimation in nonlinear simultaneous equations models (SEMs), 266-268, 272-273\\
in simultaneous equations, 254-255\\
nonlinear 3SLS (N3SLS) estimator, 531-532\\
traditional 3SLS estimator, 224-226\\
two-stage least squares (2SLS) analysis versus, 233-234\\
Threshold parameters, 655\\
Time-constant data, fixed effects methods, 301-302, 328\\
Time-demeaning matrix, 303-304, 880-882\\
Time series persistence, 196-197\\
Time-varying variances, 172\\
coefficients in unobserved effects, 551-555\\
Tobin's method. See Tobit models\\
Tobit models. See also Type I Tobit model; Type\\
II Tobit model; Type III Tobit model\\
cluster sampling, 875-876\\
correlated random effects, 708, 711-713\\
random effects, 709, 711-713\\
Tobit selection equation, 817-819\\
truncated Tobit model, 801, 815-821\\
two-limit, 703-705, 787-788\\
unobserved effects, 708

Top coding, 779-780\\
Traditional random effects probit model, 612-613\\
Traditional 3SLS estimator, 224-226\\
Treatment effects, average treatment effect (ATE), 73\\
Treatment group, difference-in-differences (DD) estimation, 147-148, 150\\
Triangular systems control function estimation, 268-271\\
nature of, 268-269\\
Truncated normal hurdle (TNH) model, 692-695\\
Truncated normal regression model, 801\\
Truncated regression, 799-802\\
Truncated Tobit model, 801, 815-821\\
Two-limit Tobit model, 703-705, 787-788\\
Two-part models, 691\\
2SLS. See Two-stage least squares (2SLS) analysis\\
Two-stage least squares (2SLS) analysis, 6, 96-112\\
estimation in nonlinear simultaneous equations models (SEMs), 266-268, 271-273\\
examples, 102, 105-106\\
first-stage regression, 97, 105\\
general treatment, 98-112\\
asymptotic efficiency, 103-104\\
asymptotic normality, 101-102\\
consistency, 98-100\\
heteroskedasticity-robust inference, 106-107\\
hypothesis testing, 104-106\\
potential pitfalls, 107-112\\
linear models estimation by, 792-798\\
nonlinear system 2SLS estimator, 531\\
ordinary least squares (OLS) analysis compared with, 107-112\\
second-stage regression, 97, 105\\
in simultaneous equations models (SEMs), 254-255, 258-259\\
specification tests, 129-141\\
endogeneity, 129-134\\
functional form, 137-138\\
heteroskedasticity, 138-141\\
overidentifying restrictions, 134-137\\
three-stage least squares (3SLS) analysis versus, 233-234, 254-255\\
two-stage least squares (2SLS) estimator, 96-98, 100, 124-125\\
with generated instruments, 124-125\\
Two-stage least squares (2SLS) residuals, 101-102\\
Two-step M-estimators, 409-413\\
Two-step maximum likelihood estimator, 499\\
Two-step partial MLE, 499\\
Type I extreme value distribution, 646-647, 998\\
Type I Tobit model, 670-689\\
data censoring, 782, 787-789\\
estimation, 676-677\\
inference, 676-677\\
reporting results, 677-680\\
specification issues, 680-689\\
testing against TNH model, 701-703\\
useful expressions, 671-676\\
Type II Tobit model, 690-703\\
exponential, 697-703\\
exponential conditional mean, 694-696\\
incidental truncation, 804-806\\
lognormal hurdle model, 694-696\\
specification issues, 680-689\\
truncated normal hurdle model, 692-694\\
Type III Tobit model, 815-817\\
Uncentered R-squared, 63\\
Unconditional hazard, 1005\\
Unconditional information matrix equality (UIME), 479\\
Unconditional maximum likelihood estimation (MLE), 541-542\\
Unconditional variance matrix, 172, 182-183, 305, 366\\
Unconfoundedness, 908, 959\\
assumption of, 822, 959\\
Underdispersion, 512-513, 725\\
Unidentified equations, 251\\
Uniform convergence in probability, 402-403\\
Uniform weak law of large numbers (UWLLN), 403-404\\
Unobserved component, 285\\
Unobserved effects, 282\\
attrition in linear panel data models, 837\\
Chamberlain approach, 347-349, 551, 616-619, 624\\
dynamic unobserved effects models, 625-630, 713-715\\
logit models under strict exogeneity, 619-625\\
minimum distance approach to unobserved effects model, 549-551\\
models of conditional expectations with, 758-759\\
probit models under strict exogeneity, 610-619\\
strictly exogenous conditional on the unobserved effect, 287-288\\
strictly exogenous corner solution responses, 707-713\\
testing for presence of, 299-300\\
time-varying coefficients on unobserved effects, 551-555\\
unobserved effects model (UEM), described, 285-287. See also Linear unobserved effects panel data models\\
Unobserved heterogeneity, 22, 285, 1003-1010, 1017-1018

Unordered response, 643\\
Unstructured working correlation matrix, 448\\
Variable addition test (VAT), 427-428, 573\\
Variable probability sampling (VP sampling), 854-856, 858-863\\
Variables\\
in conditional expectations, 13, 14, 23-24\\
control, 3-4\\
dependent (y). See Dependent variable (y)\\
endogenous. See Endogenous variables\\
exogenous, 54\\
explanatory. See Explanatory variable ( $x$ )\\
fixed explanatory, 9-11\\
independent ( $x$ ). See Explanatory variable ( $x$ )\\
indicators of $q$, 112-114\\
indicators of unobservables, 112-114\\
interactions between unobservable and observable, 73-76\\
omitted, 54-55\\
in ordinary least squares (OLS) analysis, 65-76\\
Variance\\
asymptotic. See Asymptotic variance\\
bootstrap estimate, 438\\
conditional, 32-33, 172\\
ordinary least squares (OLS) analysis heteroskedasticity, 60-62\\
homoskedasticity, 59-60, 131\\
Poisson variance assumption, 725\\
robust variance matrix estimator, 172, 297-298\\
Vuong model selection test, 505-509\\
Wage offer function, 7-8\\
Wald statistic, 46-47, 62, 104, 107, 136\\
behavior under alternatives, 430-431\\
for pooled OLS, 332-333\\
generalized method of moments (GMM) under orthogonality condition, 532\\
in bootstrap samples, 440\\
Wald tests, 420-421, 423-424\\
Weak consistency, 43\\
Weak instruments, 108\\
Weak law of large numbers (WLLN), 42, 403-404, 526\\
Weibull distribution, 987, 995-998, 1004, 1007, 1018\\
Weighted exogenous sample MLE (WESMLE), 860-861\\
Weighted least squares (WLS) analysis, 60-61 in cluster sampling, 890-891\\
Weighted M-estimator, in stratified sampling, 856-861\\
Weighted multivariate nonlinear least squares\\
(WMNLS) estimator, 444-449, 614, 761-762

Weighted nonlinear least squares (WNLS) estimator, 409-413\\
adjustments, 418-420\\
asymptotic normality, 411-413\\
consistency, 410-411\\
variable addition test (VAT) approach, 427-428\\
White standard error, 61\\
White test, 140\\
Wild bootstrap, 441\\
Willingness to pay (WTP), 780

With probability approaching one (w.p.a.1.), 39-40\\
Working correlation matrix, 448\\
Working variance matrix, 446\\
$X$ variable. See Explanatory variable ( $x$ )\\
$Y$ variable, 13, 14\\
Zero-conditional mean, 398



\end{document}
